\documentclass[12pt, a4paper]{article}
\usepackage[utf8]{inputenc}
\usepackage[T1]{fontenc}
\usepackage[spanish,shorthands=off]{babel}
\usepackage{amsmath}
\usepackage{amssymb}
\usepackage{amsfonts}
\usepackage{geometry}
\geometry{a4paper, margin=2.5cm}
\usepackage{tikz}
\usepackage{pgfplots}
\pgfplotsset{compat=1.18}
\usepackage{listings}
\usepackage{xcolor}

\definecolor{codegreen}{rgb}{0,0.6,0}
\definecolor{codegray}{rgb}{0.5,0.5,0.5}
\definecolor{codepurple}{rgb}{0.58,0,0.82}
\definecolor{backcolour}{rgb}{0.95,0.95,0.92}

\lstdefinestyle{pythonstyle}{
    backgroundcolor=\color{backcolour},   
    commentstyle=\color{codegreen},
    keywordstyle=\color{magenta},
    numberstyle=\tiny\color{codegray},
    stringstyle=\color{codepurple},
    basicstyle=\ttfamily\footnotesize,
    breakatwhitespace=false,         
    breaklines=true,                 
    captionpos=b,                    
    keepspaces=true,                 
    numbers=left,                    
    numbersep=5pt,                  
    showspaces=false,                
    showstringspaces=false,
    showtabs=false,                  
    tabsize=4,
    language=Python
}
\lstset{style=pythonstyle}

\begin{document}

\subsection*{Enunciado}
Dos atletas A y B parten desde el origen de un sistema de referencia y corren sobre una pista recta horizontal con rapideces constantes de $2\text{ m/s}$ y $5\text{ m/s}$, respectivamente. El atleta A inicia su movimiento en $t = 0\text{ s}$, mientras que el atleta B parte en $t = 5\text{ s}$.

Análisis conceptual: Antes de realizar cálculos, responda: Sin efectuar cálculos, ¿qué atleta tiene mayor rapidez? ¿Esto garantiza que uno alcanzará al otro? Explique.

\subsection*{Solución}

\subsubsection*{Análisis de las rapideces relativas}
La rapidez es una magnitud escalar que indica la distancia recorrida por unidad de tiempo. El enunciado establece explícitamente los valores para ambos corredores: el atleta A tiene una rapidez de $2\text{ m/s}$ y el atleta B de $5\text{ m/s}$. Al comparar estas magnitudes, es evidente que el atleta B cubre una distancia significativamente mayor en el mismo intervalo de tiempo.

\subsubsection*{Evaluación de la cinemática de persecución}
En la física conceptual, un problema de persecución unidimensional depende de dos factores: la ventaja inicial (ya sea espacial o temporal) y la velocidad relativa entre los cuerpos. El atleta A tiene una ventaja temporal de $5\text{ s}$, lo que se traduce en una ventaja espacial inicial. Sin embargo, al correr sobre la misma trayectoria recta, el hecho de que B posea una rapidez superior garantiza matemáticamente la intersección. Esto ocurre porque la tasa a la que B avanza es mayor que la tasa a la que A huye, cerrando indefectiblemente la brecha espacial creada durante los primeros segundos.

\subsubsection*{Conclusión}
El atleta B tiene mayor rapidez ($5\text{ m/s} > 2\text{ m/s}$). Esto sí garantiza que alcanzará al atleta A. La justificación reside en que, al moverse en la misma trayectoria y tener una tasa de desplazamiento superior, el atleta B recortará progresivamente la distancia de separación inicial hasta interceptarlo.

\newpage
\subsection*{Enunciado}
 Dos atletas A y B parten desde el origen de un sistema de referencia y corren sobre una pista recta horizontal con rapideces constantes de $2\text{ m/s}$ y $5\text{ m/s}$, respectivamente. El atleta A inicia su movimiento en $t = 0\text{ s}$, mientras que el atleta B parte en $t = 5\text{ s}$.

Análisis conceptual: Durante los primeros $5\text{ s}$, ¿qué atleta está en movimiento? ¿Qué ocurre con la separación entre ellos en ese intervalo?

\subsection*{Solución}

\subsubsection*{Análisis del intervalo de tiempo inicial}
El modelado temporal del sistema nos indica que el cronómetro global ($t$) comienza a correr en el instante cero ($t=0$). En este punto, el atleta A arranca. Por el contrario, la condición inicial del atleta B dicta que su estado de reposo absoluto se mantiene hasta que el cronómetro marque $t=5\text{ s}$. Por lo tanto, durante toda la ventana temporal de $0$ a $5$ segundos, solo existe un cuerpo con energía cinética en el sistema.

\subsubsection*{Evolución de la separación espacial}
Dado que el atleta B permanece anclado en el origen del sistema de referencia ($x=0$) mientras el atleta A se desplaza en línea recta a razón de $2$ metros cada segundo, la distancia geométrica que los separa está dictada exclusivamente por el avance de A. La separación experimenta un crecimiento lineal continuo durante este periodo.

\subsubsection*{Conclusión}
Durante los primeros $5\text{ s}$, únicamente el atleta A está en movimiento. En ese intervalo, la separación entre ellos aumenta de manera constante, acumulando una ventaja espacial a favor de A antes de que B inicie su carrera.

\newpage
\subsection*{Enunciado}
 Dos atletas A y B parten desde el origen de un sistema de referencia y corren sobre una pista recta horizontal con rapideces constantes de $2\text{ m/s}$ y $5\text{ m/s}$, respectivamente. El atleta A inicia su movimiento en $t = 0\text{ s}$, mientras que el atleta B parte en $t = 5\text{ s}$.

Análisis conceptual: Cuando el atleta B comienza a moverse, ¿la distancia entre ambos aumenta o disminuye? Justifique usando sus rapideces.

\subsection*{Solución}

\subsubsection*{Análisis de la tasa de cambio de posición relativa}
En el instante $t=5\text{ s}$, el sistema entra en una nueva fase dinámica: ambos cuerpos se encuentran simultáneamente en movimiento sobre el mismo eje coordenado. Para determinar el comportamiento de la brecha espacial entre ellos, debemos recurrir al concepto de velocidad relativa.

\subsubsection*{Justificación cinemática}
La distancia entre dos cuerpos en persecución se reduce si y solo si el cuerpo perseguidor avanza más rápido que el cuerpo perseguido. En este escenario, el atleta A (que va adelante) intenta aumentar la distancia avanzando a $2\text{ m/s}$. Simultáneamente, el atleta B (que va detrás) intenta cerrar la distancia avanzando a $5\text{ m/s}$. La velocidad relativa de acercamiento es la diferencia de sus rapideces: $5\text{ m/s} - 2\text{ m/s} = 3\text{ m/s}$. Esto significa que el esfuerzo de B sobrepasa la huida de A por un margen neto de $3$ metros cada segundo.

\subsubsection*{Conclusión}
La distancia entre ambos disminuye. Se justifica porque la rapidez del atleta B ($5\text{ m/s}$) es estrictamente mayor que la del atleta A ($2\text{ m/s}$), estableciendo una tasa relativa de acercamiento positiva que recorta la separación espacial continuamente.

\newpage
\subsection*{Enunciado}
Dos atletas A y B parten desde el origen de un sistema de referencia y corren sobre una pista recta horizontal con rapideces constantes de $2\text{ m/s}$ y $5\text{ m/s}$, respectivamente. El atleta A inicia su movimiento en $t = 0\text{ s}$, mientras que el atleta B parte en $t = 5\text{ s}$.

Análisis conceptual: Mientras ambos atletas están en movimiento, ¿quién se acerca a quién? Explique en función de sus velocidades.

\subsection*{Solución}

\subsubsection*{Definición del marco de referencia relativo}
En la mecánica clásica, el movimiento es relativo al observador. Si fijamos nuestro marco de referencia en un observador externo estático, observamos a ambos atletas alejándose del origen. Sin embargo, la pregunta exige analizar la interacción entre ellos. 

\subsubsection*{Explicación dinámica del acercamiento}
Si situamos imaginariamente nuestro marco de referencia sobre el atleta A, este se considerará a sí mismo en reposo. Desde esta perspectiva, observaría al origen alejándose hacia atrás a $2\text{ m/s}$, pero vería al atleta B acercándose por detrás a una velocidad aparente de $3\text{ m/s}$ ($5\text{ m/s} - 2\text{ m/s}$). Dado que A se encuentra en una posición espacial adelantada (mayor valor en $x$) y B está detrás con un vector velocidad de mayor magnitud apuntando en la misma dirección, es B quien acorta activamente el espacio intermedio.

\subsubsection*{Conclusión}
El atleta B se acerca al atleta A. En función de sus velocidades, al estar B ubicado detrás de A, su mayor velocidad le permite recorrer mayor espacio en el mismo tiempo, superando la tasa de alejamiento de A y cerrando la brecha entre ellos.

\newpage
\subsection*{Enunciado}
 Dos atletas A y B parten desde el origen de un sistema de referencia y corren sobre una pista recta horizontal con rapideces constantes de $2\text{ m/s}$ y $5\text{ m/s}$, respectivamente. El atleta A inicia su movimiento en $t = 0\text{ s}$, mientras que el atleta B parte en $t = 5\text{ s}$.

Análisis conceptual: En el instante en que un atleta alcanza al otro, ¿qué condición deben cumplir sus posiciones?

\subsection*{Solución}

\subsubsection*{Concepto de intersección cinemática}
En el lenguaje cotidiano, "alcanzar" o "encontrarse" implica que dos entidades comparten el mismo espacio en el mismo momento. En la formalización de la cinemática, la posición de un objeto puntual se define mediante una coordenada en un eje espacial, generalmente denotada como $x(t)$, la cual representa la distancia desde un origen común en un instante $t$.

\subsubsection*{Formalización matemática del encuentro}
Para que el evento de "alcance" se concrete de manera efectiva en el plano físico, no basta con que los atletas recorran la misma cantidad de metros en total (ya que partieron en tiempos distintos); es un requisito absoluto que, medidos desde el origen del sistema de coordenadas ($x=0$), ambos ocupen el mismo punto en la recta. Por lo tanto, si $x_A$ es la coordenada de A y $x_B$ es la coordenada de B, ambas funciones de posición deben igualarse.

\subsubsection*{Conclusión}
La condición fundamental que deben cumplir es que sus posiciones espaciales deben ser idénticas. Matemáticamente, en el instante exacto del encuentro, la ecuación de posición del atleta A debe ser igual a la ecuación de posición del atleta B ($x_A = x_B$) medido desde el mismo sistema de referencia.

\newpage
\subsection*{Enunciado}
Dos atletas A y B parten desde el origen de un sistema de referencia y corren sobre una pista recta horizontal con rapideces constantes de $2\text{ m/s}$ y $5\text{ m/s}$, respectivamente. El atleta A inicia su movimiento en $t = 0\text{ s}$, mientras que el atleta B parte en $t = 5\text{ s}$.

Luego, resuelva: Determine el instante en que el corredor B alcanza al corredor A.

\subsection*{Solución}

\subsubsection*{Planteamiento de las ecuaciones de movimiento (MRU)}
Para resolver analíticamente el problema, construimos el modelo matemático del Movimiento Rectilíneo Uniforme para cada atleta utilizando la ecuación general: $x(t) = x_0 + v(t - t_0)$.
Para el atleta A: Parte del origen ($x_0 = 0$) en el instante inicial ($t_0 = 0$) con una velocidad $v_A = 2\text{ m/s}$.
$$x_A(t) = 0 + 2(t - 0) \implies x_A(t) = 2t$$
Para el atleta B: Parte del origen ($x_0 = 0$) pero con un retraso temporal ($t_0 = 5$) con una velocidad $v_B = 5\text{ m/s}$.
$$x_B(t) = 0 + 5(t - 5) \implies x_B(t) = 5t - 25$$

\subsubsection*{Resolución del sistema para el punto de encuentro}
Aplicando la condición conceptual deducida previamente, igualamos las ecuaciones de posición ($x_A = x_B$) para hallar el instante de tiempo $t$ en el que ocurre el alcance.
$$2t = 5t - 25$$
Agrupamos los términos dependientes del tiempo:
$$25 = 5t - 2t$$
$$25 = 3t$$
$$t = \frac{25}{3} \approx 8.33\text{ s}$$

\subsubsection*{Representación gráfica del encuentro}
Para consolidar la comprensión pedagógica de cómo la pendiente más pronunciada (mayor velocidad) del corredor B le permite interceptar la trayectoria del corredor A a pesar de su retraso inicial, se adjunta el siguiente modelo en Python de la gráfica de Posición vs. Tiempo.

\begin{lstlisting}[language=Python]
import matplotlib.pyplot as plt
import numpy as np

fig, ax = plt.subplots(figsize=(7, 4.5))
t = np.linspace(0, 10, 100)
x_A = 2 * t
x_B = np.piecewise(t, [t < 5, t >= 5], [0, lambda t: 5 * (t - 5)])

ax.plot(t, x_A, label='Atleta A (v = 2 m/s)', color='blue', linewidth=2)
ax.plot(t, x_B, label='Atleta B (v = 5 m/s, t0 = 5 s)', color='red', linewidth=2)

t_enc = 25/3
x_enc = 50/3
ax.plot(t_enc, x_enc, 'ko', markersize=6)
ax.annotate(f'Punto de Alcance\n(t = {t_enc:.2f} s)', xy=(t_enc, x_enc),
            xytext=(t_enc - 4, x_enc + 4),
            arrowprops=dict(facecolor='black', shrink=0.05, width=1.5, headwidth=6),
            fontweight='bold')

ax.set_xlabel('Tiempo Global (s)', fontweight='bold')
ax.set_ylabel('Posicion desde el origen (m)', fontweight='bold')
ax.set_title('Grafica Cinematica: Posicion vs Tiempo')
ax.grid(True, linestyle='--', alpha=0.6)
ax.legend(loc='upper left')

plt.show()
\end{lstlisting}

\subsubsection*{Conclusión}
El corredor B alcanza al corredor A en el instante exacto de $t = 8.33\text{ s}$ (o en forma fraccionaria, a los $25/3$ de segundo) medido desde que el cronómetro global inició con la partida del primer atleta.

\newpage
\subsection*{Enunciado}
 Dos atletas A y B parten desde el origen de un sistema de referencia y corren sobre una pista recta horizontal con rapideces constantes de $2\text{ m/s}$ y $5\text{ m/s}$, respectivamente. El atleta A inicia su movimiento en $t = 0\text{ s}$, mientras que el atleta B parte en $t = 5\text{ s}$.

Luego, resuelva: Calcule la distancia recorrida por A desde que parte hasta que es alcanzado por B.

\subsection*{Solución}

\subsubsection*{Aplicación del modelo cinemático}
Una vez que hemos determinado el instante preciso en el que ocurre la intersección ($t = 25/3\text{ s}$), podemos calcular la ubicación espacial de ese evento sustituyendo este valor temporal en la ecuación de movimiento de cualquiera de los dos atletas. Dado que el problema solicita la distancia recorrida por el atleta A, utilizaremos su ecuación específica.

\subsubsection*{Cálculo de la distancia (desplazamiento)}
Dado que el atleta A parte del origen ($x_0 = 0$) y corre en una sola dirección sin retroceder, la magnitud de su posición final en el momento del alcance equivale exactamente a la distancia total recorrida.
Utilizando la ecuación $x_A(t) = 2t$:
$$d_A = x_A\left(\frac{25}{3}\right) = 2 \cdot \left(\frac{25}{3}\right)$$
$$d_A = \frac{50}{3} \approx 16.67\text{ m}$$

*(Nota de comprobación: Si evaluamos el avance de B, este corre durante $25/3 - 5 = 10/3$ segundos. Multiplicando esto por su rapidez de $5\text{ m/s}$, obtenemos $5 \cdot 10/3 = 50/3\text{ m}$, lo que confirma que ambos se encuentran en la misma posición).*

\subsubsection*{Conclusión}
La distancia total recorrida por el atleta A desde que parte hasta que es alcanzado por B es de $16.67\text{ m}$ (exactamente $50/3\text{ m}$).

\newpage
\subsection*{Enunciado}
 Dos automóviles A y B están separados inicialmente por una distancia horizontal de 2000 m. El automóvil A parte directamente hacia B con una rapidez inicial de 20 m/s y una aceleración constante de magnitud 4 m/s² en el mismo sentido de su velocidad. Ocho segundos más tarde, el automóvil B parte desde el reposo y se mueve hacia A con aceleración constante de magnitud 7 m/s².

Análisis conceptual: ¿Cómo se mueve el automóvil A durante los primeros 8 s? Describa qué ocurre con su rapidez y con la distancia que lo separa de B.

\subsection*{Solución}

\subsubsection*{Análisis del intervalo de movimiento solitario}
Durante los primeros 8 s, el automóvil A se mueve bajo las condiciones de un Movimiento Rectilíneo Uniformemente Variado (MRUV) acelerado. Dado que posee una velocidad inicial de 20 m/s y una aceleración a favor del movimiento, su rapidez no es constante, sino que aumenta de forma continua a razón de 4 m/s por cada segundo que transcurre. 
Respecto a la distancia, como el automóvil A avanza activamente en dirección al automóvil B (el cual permanece estacionario en este lapso), la separación geométrica entre ambos disminuye progresivamente desde los 2000 m iniciales.

\newpage
\subsection*{Enunciado}
 Dos automóviles A y B están separados inicialmente por una distancia horizontal de 2000 m. El automóvil A parte directamente hacia B con una rapidez inicial de 20 m/s y una aceleración constante de magnitud 4 m/s² en el mismo sentido de su velocidad. Ocho segundos más tarde, el automóvil B parte desde el reposo y se mueve hacia A con aceleración constante de magnitud 7 m/s².
 Análisis conceptual: Antes de que el automóvil B comience a moverse, ¿la separación entre ambos aumenta o disminuye? Explique.

\subsection*{Solución}

\subsubsection*{Variación de la separación espacial}
La separación disminuye. La explicación cinemática es directa: la distancia entre dos puntos se acorta si uno de ellos se mueve en dirección al otro mientras este último permanece fijo. En esta ventana temporal (de 0 a 8 segundos), el automóvil B tiene una velocidad de 0 m/s, mientras que el automóvil A avanza hacia él con una velocidad cada vez mayor.

\newpage
\subsection*{Enunciado}
 Dos automóviles A y B están separados inicialmente por una distancia horizontal de 2000 m. El automóvil A parte directamente hacia B con una rapidez inicial de 20 m/s y una aceleración constante de magnitud 4 m/s² en el mismo sentido de su velocidad. Ocho segundos más tarde, el automóvil B parte desde el reposo y se mueve hacia A con aceleración constante de magnitud 7 m/s².
 Análisis conceptual: Cuando ambos automóviles están en movimiento, ¿qué ocurre con la distancia entre ellos? Justifique en función de la dirección de sus movimientos.

\subsection*{Solución}

\subsubsection*{Superposición de movimientos convergentes}
La distancia entre ellos disminuye a un ritmo mucho más acelerado. A partir del segundo 8, el sistema entra en una fase de trayectorias convergentes. El automóvil A continúa desplazándose hacia la derecha (hacia B) y el automóvil B comienza a desplazarse hacia la izquierda (hacia A). Al tener direcciones de movimiento diametralmente opuestas y apuntando una hacia la otra, las velocidades de ambos vehículos se suman en términos relativos, cerrando la brecha espacial de forma conjunta y agresiva.

\newpage
\subsection*{Enunciado}
Dos automóviles A y B están separados inicialmente por una distancia horizontal de 2000 m. El automóvil A parte directamente hacia B con una rapidez inicial de 20 m/s y una aceleración constante de magnitud 4 m/s² en el mismo sentido de su velocidad. Ocho segundos más tarde, el automóvil B parte desde el reposo y se mueve hacia A con aceleración constante de magnitud 7 m/s².
Análisis conceptual: En el caso del automóvil A, la aceleración tiene el mismo sentido que la velocidad. ¿Qué implica esto sobre su rapidez?

\subsection*{Solución}

\subsubsection*{Relación vectorial entre velocidad y aceleración}
En la mecánica clásica, la aceleración es la tasa de cambio del vector velocidad. Cuando el vector aceleración y el vector velocidad son colineales y poseen el mismo sentido, la aceleración actúa a favor de la inercia del cuerpo. Esto implica que la rapidez (la magnitud escalar de la velocidad) experimenta un incremento continuo. El vehículo viaja cada vez más rápido en la misma dirección.

\newpage
\subsection*{Enunciado}
 Dos automóviles A y B están separados inicialmente por una distancia horizontal de 2000 m. El automóvil A parte directamente hacia B con una rapidez inicial de 20 m/s y una aceleración constante de magnitud 4 m/s² en el mismo sentido de su velocidad. Ocho segundos más tarde, el automóvil B parte desde el reposo y se mueve hacia A con aceleración constante de magnitud 7 m/s².
 Análisis conceptual: El automóvil B parte desde el reposo y acelera hacia A. ¿Cómo cambia su rapidez con el tiempo?

\subsection*{Solución}

\subsubsection*{Evolución de la rapidez desde el reposo}
Partir desde el reposo significa que la velocidad inicial de B en el instante $t = 8$ s es exactamente 0 m/s. Al aplicarse una aceleración constante de 7 m/s² hacia A, su rapidez deja de ser nula y comienza a aumentar linealmente con el tiempo. Cada segundo que el automóvil B está en movimiento, su rapidez se incrementa en 7 m/s.

\newpage
\subsection*{Enunciado}
 Dos automóviles A y B están separados inicialmente por una distancia horizontal de 2000 m. El automóvil A parte directamente hacia B con una rapidez inicial de 20 m/s y una aceleración constante de magnitud 4 m/s² en el mismo sentido de su velocidad. Ocho segundos más tarde, el automóvil B parte desde el reposo y se mueve hacia A con aceleración constante de magnitud 7 m/s².
 Análisis conceptual: En el instante del encuentro, ¿qué condición deben cumplir las posiciones de ambos automóviles?

\subsection*{Solución}

\subsubsection*{Condición matemática de intersección}
Para que un encuentro físico ocurra, ambos vehículos deben ocupar el mismo punto geométrico en el espacio de manera simultánea. En términos analíticos, si establecemos un único sistema de coordenadas de referencia (por ejemplo, situando el origen en el punto de partida de A), la condición absoluta es que la ecuación de posición del automóvil A sea exactamente igual a la ecuación de posición del automóvil B. 
$$x_A(t) = x_B(t)$$

\newpage
\subsection*{Enunciado}
Dos automóviles A y B están separados inicialmente por una distancia horizontal de 2000 m. El automóvil A parte directamente hacia B con una rapidez inicial de 20 m/s y una aceleración constante de magnitud 4 m/s² en el mismo sentido de su velocidad. Ocho segundos más tarde, el automóvil B parte desde el reposo y se mueve hacia A con aceleración constante de magnitud 7 m/s².
Determine el instante en que ambos automóviles se encuentran.

\subsection*{Solución}

\subsubsection*{Planteamiento de las ecuaciones cinemáticas}
Estableceremos el origen de coordenadas ($x = 0$ m) en la posición inicial del automóvil A, y el instante $t = 0$ s como el momento en que A inicia su movimiento. 

Para el automóvil A (movimiento hacia la derecha, acelerado):
* Posición inicial: $x_{0A} = 0$ m
* Velocidad inicial: $v_{0A} = 20$ m/s
* Aceleración: $a_A = 4$ m/s²
La ecuación de posición es $x(t) = x_0 + v_0 t + \frac{1}{2} a t^2$:
$$x_A(t) = 0 + 20t + \frac{1}{2}(4)t^2$$
$$x_A(t) = 20t + 2t^2$$

Para el automóvil B (movimiento hacia la izquierda, acelerado con retraso temporal):
* Posición inicial: $x_{0B} = 2000$ m
* Velocidad inicial: $v_{0B} = 0$ m/s (reposo)
* Aceleración: $a_B = -7$ m/s² (negativa porque apunta hacia la izquierda)
* Retraso temporal: $t_0 = 8$ s. Su movimiento solo existe para $t \ge 8$.
La ecuación de posición modificada por el retraso es $x(t) = x_0 + v_0(t - t_0) + \frac{1}{2} a (t - t_0)^2$:
$$x_B(t) = 2000 + 0(t - 8) + \frac{1}{2}(-7)(t - 8)^2$$
$$x_B(t) = 2000 - 3.5(t - 8)^2$$

\subsubsection*{Resolución del sistema para el tiempo de alcance}
Aplicamos la condición de encuentro $x_A(t) = x_B(t)$:
$$20t + 2t^2 = 2000 - 3.5(t - 8)^2$$
Desarrollamos el binomio al cuadrado:
$$20t + 2t^2 = 2000 - 3.5(t^2 - 16t + 64)$$
$$20t + 2t^2 = 2000 - 3.5t^2 + 56t - 224$$
Agrupamos todos los términos en el lado izquierdo para formar una ecuación cuadrática igualada a cero:
$$2t^2 + 3.5t^2 + 20t - 56t - 2000 + 224 = 0$$
$$5.5t^2 - 36t - 1776 = 0$$
Aplicamos la fórmula general para ecuaciones de segundo grado ($t = \frac{-b \pm \sqrt{b^2 - 4ac}}{2a}$):
$$t = \frac{-(-36) \pm \sqrt{(-36)^2 - 4(5.5)(-1776)}}{2(5.5)}$$
$$t = \frac{36 \pm \sqrt{1296 + 39072}}{11}$$
$$t = \frac{36 \pm \sqrt{40368}}{11}$$
Dado que buscamos un tiempo futuro ($t > 8$), tomamos la raíz positiva ($\sqrt{40368} \approx 200.918$):
$$t = \frac{36 + 200.918}{11} \approx \frac{236.918}{11} \approx 21.538 \text{ s}$$

\subsubsection*{Conclusión}
Ambos automóviles se encuentran en el instante $t = 21.54$ s, medido desde el momento en que el automóvil A comenzó a moverse.

\newpage
\subsection*{Enunciado}
Dos automóviles A y B están separados inicialmente por una distancia horizontal de 2000 m. El automóvil A parte directamente hacia B con una rapidez inicial de 20 m/s y una aceleración constante de magnitud 4 m/s² en el mismo sentido de su velocidad. Ocho segundos más tarde, el automóvil B parte desde el reposo y se mueve hacia A con aceleración constante de magnitud 7 m/s².
resuelva: Determine la distancia medida desde la posición inicial de A hasta el punto donde ocurre el encuentro.

\subsection*{Solución}

\subsubsection*{Cálculo de la posición de encuentro}
Puesto que nuestro sistema de coordenadas ubicó el origen ($0$ m) exactamente en la posición inicial de A, la coordenada final del encuentro representa de manera directa la distancia recorrida por A. 
Sustituimos el tiempo de encuentro ($t = 21.538$ s) en la ecuación de posición del automóvil A:
$$x_A(t) = 20t + 2t^2$$
$$x_A(21.538) = 20(21.538) + 2(21.538)^2$$
$$x_A = 430.76 + 2(463.885)$$
$$x_A = 430.76 + 927.77$$
$$x_A \approx 1358.53 \text{ m}$$

*(Podemos validar este resultado comprobando la posición de B: $x_B = 2000 - 3.5(21.538 - 8)^2 = 2000 - 3.5(183.277) = 2000 - 641.47 = 1358.53$ m. La intersección es perfecta).*

\subsubsection*{Representación Gráfica del Sistema de Trayectorias}
Para visualizar cómo el automóvil A avanza acelerando de manera continua mientras que el automóvil B permanece estático los primeros 8 segundos antes de acelerar bruscamente en sentido contrario, se presenta la siguiente simulación gráfica generada en Python. Se puede observar claramente la intersección de ambas curvas parabólicas.

\begin{lstlisting}[language=Python]
import matplotlib.pyplot as plt
import numpy as np

fig, ax = plt.subplots(figsize=(8, 5))
t = np.linspace(0, 25, 200)

x_A = 20 * t + 2 * t**2
x_B = np.piecewise(t, [t < 8, t >= 8], [2000, lambda t: 2000 - 3.5 * (t - 8)**2])

ax.plot(t, x_A, label='Automovil A', color='blue', linewidth=2.5)
ax.plot(t, x_B, label='Automovil B (Retraso 8 s)', color='red', linewidth=2.5)

t_enc = 21.54
x_enc = 1358.53
ax.plot(t_enc, x_enc, 'ko', markersize=7)
ax.annotate(f'Encuentro\n(t = {t_enc:.2f} s,\n x = {x_enc:.1f} m)',
            xy=(t_enc, x_enc), xytext=(t_enc - 8, x_enc + 200),
            arrowprops=dict(facecolor='black', shrink=0.05, width=1.5, headwidth=7),
            fontweight='bold', fontsize=10)

ax.set_xlabel('Tiempo de recorrido (s)', fontweight='bold')
ax.set_ylabel('Posicion desde Origen A (m)', fontweight='bold')
ax.set_title('Grafica Cinematica: Posicion vs Tiempo (MRUV Convergente)')
ax.grid(True, linestyle='--', alpha=0.7)
ax.legend(loc='center left')

plt.show()
\end{lstlisting}

\subsubsection*{Conclusión}
La distancia medida desde la posición inicial de A hasta el punto exacto donde ocurre el encuentro es de $1358.53$ m.

\newpage

\subsection*{Enunciado}

Una vía férrea recta y horizontal de 200 km conecta las ciudades A y B. En $t = 0 \text{ h}$, un tren parte desde A y se mueve hacia B con rapidez constante de 40 km/h. En el mismo instante, un tren parte desde B y se mueve hacia A con rapidez constante de 60 km/h. Determine el instante en el que ambos trenes se encuentran.

\vspace{0.5cm}

\begin{center}
\begin{tikzpicture}
\draw[thick, <->] (0,0) -- (10,0);
\filldraw[black] (1,0) circle (2pt) node[below=0.2cm] {Ciudad A ($x=0$)};
\filldraw[black] (9,0) circle (2pt) node[below=0.2cm] {Ciudad B ($x=200 \text{ km}$)};
\draw[->, thick, blue] (1, 0.5) -- (2.5, 0.5) node[midway, above] {$v_A = 40 \text{ km/h}$};
\draw[->, thick, red] (9, 0.5) -- (7.5, 0.5) node[midway, above] {$v_B = 60 \text{ km/h}$};
\end{tikzpicture}
\end{center}

\subsection*{Solución}

\subsubsection*{Conceptos Físicos Involucrados}

De acuerdo con la pedagogía de la física conceptual, el movimiento es relativo y debe describirse siempre respecto a un sistema de referencia espacial y temporal. Cuando un objeto se mueve a lo largo de una trayectoria recta cubriendo distancias iguales en intervalos de tiempo iguales, decimos que tiene una velocidad constante y experimenta un Movimiento Rectilíneo Uniforme. En este tipo de movimiento, la aceleración es nula y la posición en cualquier instante de tiempo depende linealmente de la velocidad y de la posición inicial.

Para modelar matemáticamente este fenómeno, debemos asignar signos a las velocidades dependiendo de la dirección de movimiento de cada objeto respecto al eje de coordenadas elegido.

\subsubsection*{Definición del Sistema de Referencia}

Establecemos el origen de nuestro sistema de coordenadas espaciales en la Ciudad A. Por lo tanto, la posición inicial del tren A es nula, y la posición inicial del tren B está a 200 km de distancia en el eje positivo de las equis. 

El tren A se mueve en la dirección positiva del eje, por lo que su velocidad es matemáticamente positiva. El tren B se mueve en la dirección opuesta, acercándose al origen, por lo que su velocidad debe ser considerada matemáticamente negativa.

Datos iniciales para el Tren A:
\begin{itemize}
    \item Posición inicial $x_{A0} = 0 \text{ km}$
    \item Velocidad $v_A = 40 \text{ km/h}$
\end{itemize}

Datos iniciales para el Tren B:
\begin{itemize}
    \item Posición inicial $x_{B0} = 200 \text{ km}$
    \item Velocidad $v_B = -60 \text{ km/h}$
\end{itemize}

\subsubsection*{Ecuaciones de Movimiento}

La ecuación general de posición para un objeto a velocidad constante es:
\begin{equation}
x(t) = x_0 + v \cdot t
\end{equation}

Sustituyendo los valores definidos en nuestro sistema de referencia, obtenemos las ecuaciones de posición para cada tren en función del tiempo:

Para el Tren A:
\begin{equation}
x_A(t) = 0 + 40t = 40t
\end{equation}

Para el Tren B:
\begin{equation}
x_B(t) = 200 - 60t
\end{equation}

\subsubsection*{Condición de Encuentro}

El concepto de "encuentro" en cinemática significa que dos objetos ocupan la misma posición exacta en el mismo instante de tiempo. Por lo tanto, la condición matemática para resolver el problema es igualar las ecuaciones de posición de ambos trenes.

\begin{equation}
x_A(t) = x_B(t)
\end{equation}

\subsubsection*{Desarrollo Matemático}

Sustituimos las expresiones obtenidas en la condición de encuentro para despejar la variable del tiempo temporal:

\begin{equation}
40t = 200 - 60t
\end{equation}

Agrupamos los términos con la variable temporal en un solo lado de la ecuación sumando sesenta unidades de tiempo en ambos lados:

\begin{equation}
40t + 60t = 200
\end{equation}

\begin{equation}
100t = 200
\end{equation}

Despejamos el tiempo dividiendo ambos lados de la ecuación por cien:

\begin{equation}
t = \frac{200}{100}
\end{equation}

\begin{equation}
t = 2.0 \text{ h}
\end{equation}

\subsubsection*{Representación Gráfica del Encuentro}

Una de las formas más efectivas de comprender el movimiento relativo es mediante una gráfica de posición en función del tiempo. El punto donde las rectas se cruzan representa el instante y lugar exacto del encuentro. A continuación, se presenta el código en Python para generar dicha gráfica pedagógica.

\begin{lstlisting}[language=Python]
import matplotlib.pyplot as plt
import numpy as np

t = np.linspace(0, 3, 100)
x_A = 40 * t
x_B = 200 - 60 * t

plt.figure(figsize=(8, 5))
plt.plot(t, x_A, label='Tren A (hacia Ciudad B)', color='blue', linewidth=2)
plt.plot(t, x_B, label='Tren B (hacia Ciudad A)', color='red', linewidth=2)

plt.axhline(0, color='black', linewidth=0.5)
plt.axvline(0, color='black', linewidth=0.5)

plt.plot(2, 80, 'go', markersize=8, label='Punto de encuentro (2 h, 80 km)')
plt.vlines(2, 0, 80, linestyle='dashed', color='green')
plt.hlines(80, 0, 2, linestyle='dashed', color='green')

plt.title('Posicion vs Tiempo para los Trenes A y B')
plt.xlabel('Tiempo (h)')
plt.ylabel('Posicion (km)')
plt.legend()
plt.grid(True)
plt.show()
\end{lstlisting}

\subsubsection*{Conclusión}

El instante en el que ambos trenes se encuentran es a las 2.0 horas después de haber partido.

\newpage

\subsection*{Enunciado}

Una vía férrea recta y horizontal de 200 km conecta las ciudades A y B. En $t = 0 \text{ h}$, un tren parte desde A y se mueve hacia B con rapidez constante de 40 km/h. En el mismo instante, un tren parte desde B y se mueve hacia A con rapidez constante de 60 km/h. Determine la distancia entre la ciudad A y el punto de encuentro.

\vspace{0.5cm}

\begin{center}
\begin{tikzpicture}
\draw[thick, <->] (0,0) -- (10,0);
\filldraw[black] (1,0) circle (2pt) node[below=0.2cm] {Ciudad A ($x=0$)};
\filldraw[black] (9,0) circle (2pt) node[below=0.2cm] {Ciudad B ($x=200 \text{ km}$)};
\draw[->, thick, blue] (1, 0.5) -- (2.5, 0.5) node[midway, above] {$v_A = 40 \text{ km/h}$};
\draw[->, thick, red] (9, 0.5) -- (7.5, 0.5) node[midway, above] {$v_B = 60 \text{ km/h}$};
\filldraw[green!70!black] (4,0) circle (3pt) node[below=0.2cm] {Encuentro};
\draw[dashed, thick, green!70!black] (1, -1) -- (4, -1) node[midway, below] {$d = ?$};
\draw[dashed, thick, gray] (1,0) -- (1,-1.2);
\draw[dashed, thick, gray] (4,0) -- (4,-1.2);
\end{tikzpicture}
\end{center}

\subsection*{Solución}

\subsubsection*{Conceptos Físicos Involucrados}

La distancia es una magnitud escalar que representa la longitud de la trayectoria recorrida. Dado que hemos establecido nuestro sistema de referencia con el origen ubicado exactamente en la Ciudad A, cualquier valor positivo de la coordenada espacial en nuestro sistema representa directamente la distancia medida desde dicha ciudad. 

Por lo tanto, para encontrar a qué distancia de la Ciudad A ocurre el encuentro, basta con evaluar la función de posición de cualquiera de los dos trenes en el instante de tiempo en el que coinciden.

\subsubsection*{Recapitulación del Instante de Encuentro}

Como se demostró al analizar el cruce de ambos trenes, la condición cinemática de posiciones iguales estableció que el encuentro ocurre transcurridas dos horas desde la partida inicial. 

\begin{equation}
t_{encuentro} = 2.0 \text{ h}
\end{equation}

\subsubsection*{Desarrollo Matemático}

Utilizamos la ecuación de posición del Tren A, ya que este parte exactamente desde nuestro punto de interés topográfico. Su ecuación de movimiento es:

\begin{equation}
x_A(t) = 40t
\end{equation}

Sustituimos el tiempo de encuentro en la ecuación:

\begin{equation}
x_A(2.0) = 40 \cdot (2.0)
\end{equation}

Realizamos el producto matemático:

\begin{equation}
x_A(2.0) = 80 \text{ km}
\end{equation}

Para asegurar la rigurosidad conceptual, podemos verificar este resultado sustituyendo el mismo instante de tiempo en la ecuación de movimiento del Tren B. Ambas ecuaciones deben arrojar la misma posición en el espacio si el tiempo calculado es correcto.

\begin{equation}
x_B(2.0) = 200 - 60 \cdot (2.0)
\end{equation}

\begin{equation}
x_B(2.0) = 200 - 120
\end{equation}

\begin{equation}
x_B(2.0) = 80 \text{ km}
\end{equation}

Ambas ecuaciones confirman que en el instante de dos horas, la coordenada espacial compartida es de ochenta kilómetros. Al estar el origen en la Ciudad A, esta coordenada corresponde directamente a la distancia solicitada.

\subsubsection*{Conclusión}

La distancia entre la ciudad A y el punto de encuentro es de 80 km.
\newpage

\subsection*{Enunciado}

Una vía férrea recta y horizontal de 200 km conecta las ciudades A y B. Considere que, respecto al contexto original, el primer tren parte desde la ciudad A y se mueve hacia B con rapidez constante de 40 km/h, pero con un retraso de 30 minutos ($0.5 \text{ h}$) respecto al instante $t = 0 \text{ h}$. En $t = 0 \text{ h}$, el segundo tren parte desde B y se mueve hacia A con rapidez constante de 60 km/h. Determine el instante en el que ambos trenes se encuentran.

\vspace{0.5cm}

\begin{center}
\begin{tikzpicture}
\draw[thick, <->] (0,0) -- (10,0);
\filldraw[black] (1,0) circle (2pt) node[below=0.2cm] {Ciudad A ($x=0$)};
\filldraw[black] (9,0) circle (2pt) node[below=0.2cm] {Ciudad B ($x=200 \text{ km}$)};
\draw[->, thick, blue] (1, 0.5) -- (2.5, 0.5) node[midway, above] {$v_A = 40 \text{ km/h (retraso 0.5 h)}$};
\draw[->, thick, red] (9, 0.5) -- (7.5, 0.5) node[midway, above] {$v_B = 60 \text{ km/h}$};
\end{tikzpicture}
\end{center}

\subsection*{Solución}

\subsubsection*{Resultado Directo}

El encuentro ocurre en $t = 2.2 \text{ h}$.

\subsubsection*{Modelado Cinemático Inicial}

Los parámetros lógicos y las condiciones espaciales iniciales son validados como contexto establecido. Se asume el desfase temporal en la posición del origen para reconfigurar la cinemática relativa del sistema.

\subsubsection*{Desarrollo Matemático}

Mediante la igualación de las funciones de posición adaptadas al retraso, se ejecuta la transformación algebraica directa hacia la variable temporal:

\begin{equation}
40(t - 0.5) = 200 - 60t \implies t = 2.2 \text{ h}
\end{equation}

\subsubsection*{Conclusión}

El instante en el que ambos trenes se encuentran es 2.2 h.

\newpage

\subsection*{Enunciado}

Una vía férrea recta y horizontal de 200 km conecta las ciudades A y B. Considere que el primer tren parte desde la ciudad A y se mueve hacia B con rapidez constante de 40 km/h, pero con un retraso de 30 minutos ($0.5 \text{ h}$) respecto al instante $t = 0 \text{ h}$. En $t = 0 \text{ h}$, el segundo tren parte desde B y se mueve hacia A con rapidez constante de 60 km/h. Determine la distancia entre la ciudad A y el punto de encuentro.

\vspace{0.5cm}

\begin{center}
\begin{tikzpicture}
\draw[thick, <->] (0,0) -- (10,0);
\filldraw[black] (1,0) circle (2pt) node[below=0.2cm] {Ciudad A ($x=0$)};
\filldraw[black] (9,0) circle (2pt) node[below=0.2cm] {Ciudad B ($x=200 \text{ km}$)};
\draw[->, thick, blue] (1, 0.5) -- (2.5, 0.5) node[midway, above] {$v_A = 40 \text{ km/h}$};
\draw[->, thick, red] (9, 0.5) -- (7.5, 0.5) node[midway, above] {$v_B = 60 \text{ km/h}$};
\filldraw[green!70!black] (3.4,0) circle (3pt) node[below=0.2cm] {Encuentro};
\draw[dashed, thick, green!70!black] (1, -1) -- (3.4, -1) node[midway, below] {$d = ?$};
\draw[dashed, thick, gray] (1,0) -- (1,-1.2);
\draw[dashed, thick, gray] (3.4,0) -- (3.4,-1.2);
\end{tikzpicture}
\end{center}

\subsection*{Solución}

\subsubsection*{Resultado Directo}

La distancia es de 68 km.

\subsubsection*{Evaluación de Coordenadas}

El tiempo de intersección calculado previamente es validado como parámetro operativo constante para el análisis posicional.

\subsubsection*{Desarrollo Matemático}

Sustituyendo el tiempo ajustado directamente en la función de desplazamiento del origen, se obtiene la distancia final:

\begin{equation}
x_A = 40(2.2 - 0.5) = 68 \text{ km}
\end{equation}

\subsubsection*{Conclusión}

La distancia entre la ciudad A y el punto de encuentro es de 68 km.

\newpage

\subsection*{Enunciado}

Un automovilista viaja por una carretera recta horizontal a 90 km/h cuando de repente ve un obstáculo 50 m adelante. Si tarda 0.4 s en aplicar los frenos, y cuando lo hace desacelera uniformemente a razón de 5 m/s$^2$. Determine si choca o no con el obstáculo. En caso afirmativo, determine cuál es la rapidez del impacto.

\vspace{0.5cm}

\begin{center}
\begin{tikzpicture}
\draw[thick] (0,0) -- (10,0);
\draw[fill=blue!50] (0.5,0) rectangle (2,1);
\filldraw[black] (0.8,0) circle (0.2);
\filldraw[black] (1.7,0) circle (0.2);
\draw[fill=red!60] (8.5,0) rectangle (9,1.5);
\node[above] at (1.25, 1) {Auto};
\node[above] at (8.75, 1.5) {Obstáculo};
\draw[<->, dashed] (2, -0.5) -- (8.5, -0.5) node[midway, below] {$d = 50 \text{ m}$};
\draw[->, thick, blue] (2.2, 0.5) -- (3.5, 0.5) node[midway, above] {$v_0 = 90 \text{ km/h}$};
\draw[<-, thick, red] (1.25, 1.5) -- (2.5, 1.5) node[midway, above] {$a = 5 \text{ m/s}^2$};
\end{tikzpicture}
\end{center}

\subsection*{Solución}

\subsubsection*{Conversión de Unidades}

En la física conceptual, es fundamental trabajar en un sistema de unidades coherente para evitar discrepancias numéricas. El Sistema Internacional utiliza metros por segundo para la velocidad y metros por segundo al cuadrado para la aceleración. Convertimos la velocidad inicial multiplicando por el factor de conversión correspondiente a metros y segundos.

\begin{equation}
v_0 = 90 \text{ \frac{km}{h}} \cdot \frac{1000 \text{ m}}{1 \text{ km}} \cdot \frac{1 \text{ h}}{3600 \text{ s}} = 25 \text{ \frac{m}{s}}
\end{equation}

\subsubsection*{Fase de Reacción y Principio de Inercia}

Según la primera ley de Newton, un objeto en movimiento tiende a permanecer en su estado de movimiento uniforme a menos que una fuerza externa actúe sobre él. Durante el tiempo de reacción del conductor, los frenos aún no se han aplicado, por lo que no hay fuerza neta de frenado. El vehículo se mueve por inercia describiendo un Movimiento Rectilíneo Uniforme. Calculamos la distancia recorrida de forma ciega en esta fase.

\begin{equation}
d_r = v_0 \cdot t_r
\end{equation}

\begin{equation}
d_r = 25 \cdot 0.4 = 10 \text{ m}
\end{equation}

El automóvil avanza diez metros antes de que los componentes mecánicos comiencen a frenar el vehículo. Esto significa que la distancia útil restante para detenerse antes de impactar el obstáculo se reduce considerablemente.

\begin{equation}
d_{\text{restante}} = 50 - 10 = 40 \text{ m}
\end{equation}

\subsubsection*{Fase de Frenado y Determinación del Choque}

Al aplicar los frenos, el vehículo experimenta un Movimiento Rectilíneo Uniformemente Variado con una aceleración contraria al vector velocidad. Para saber si choca, calcularemos la distancia teórica que necesitaría para detenerse por completo y la compararemos con la distancia útil restante.

Utilizamos la ecuación cinemática que relaciona el cuadrado de las velocidades, la aceleración constante y la distancia de frenado:

\begin{equation}
v_f^2 = v_0^2 + 2 \cdot a \cdot d_{\text{parada}}
\end{equation}

Sustituyendo la velocidad final nula, la velocidad inicial y la aceleración negativa que provee el problema:

\begin{equation}
0 = (25)^2 + 2 \cdot (-5) \cdot d_{\text{parada}}
\end{equation}

\begin{equation}
0 = 625 - 10 \cdot d_{\text{parada}}
\end{equation}

\begin{equation}
10 \cdot d_{\text{parada}} = 625 \implies d_{\text{parada}} = 62.5 \text{ m}
\end{equation}

Como la distancia estrictamente necesaria para detenerse es mayor a los cuarenta metros disponibles físicamente, el vehículo inevitablemente chocará con el obstáculo antes de poder disipar toda su energía cinética.

\subsubsection*{Cálculo de la Rapidez de Impacto}

Para encontrar la rapidez con la que el auto golpea el obstáculo, utilizamos la misma ecuación cinemática del paso anterior, pero ahora evaluamos la velocidad en el instante en que el vehículo ha recorrido exactamente los cuarenta metros restantes.

\begin{equation}
v_{\text{impacto}}^2 = v_0^2 + 2 \cdot a \cdot d_{\text{restante}}
\end{equation}

\begin{equation}
v_{\text{impacto}}^2 = (25)^2 + 2 \cdot (-5) \cdot (40)
\end{equation}

\begin{equation}
v_{\text{impacto}}^2 = 625 - 400 = 225
\end{equation}

Aplicando la raíz cuadrada para despejar la rapidez:

\begin{equation}
v_{\text{impacto}} = \sqrt{225} = 15 \text{ \frac{m}{s}}
\end{equation}

\subsubsection*{Análisis Gráfico del Movimiento}

Una gráfica de velocidad en función del tiempo permite al estudiante visualizar claramente las dos fases del movimiento: la fase de velocidad constante por inercia y la fase de desaceleración lineal. El área bajo la curva hasta el instante de la colisión representa el desplazamiento total de cincuenta metros.

\begin{lstlisting}[language=Python]
import matplotlib.pyplot as plt
import numpy as np

t_reaccion = np.linspace(0, 0.4, 50)
v_reaccion = np.full_like(t_reaccion, 25)

t_frenado = np.linspace(0.4, 2.4, 100)
v_frenado = 25 - 5 * (t_frenado - 0.4)

plt.figure(figsize=(8, 5))
plt.plot(t_reaccion, v_reaccion, label='Fase de reaccion (MRU)', color='blue', linewidth=2)
plt.plot(t_frenado, v_frenado, label='Fase de frenado (MRUV)', color='red', linewidth=2)

plt.fill_between(t_reaccion, v_reaccion, color='blue', alpha=0.2)
plt.fill_between(t_frenado, v_frenado, color='red', alpha=0.2)

plt.axvline(2.4, color='black', linestyle='dashed')
plt.plot(2.4, 15, 'ko', markersize=6)
plt.text(2.45, 15, 'Impacto (15 m/s)', verticalalignment='center')

plt.title('Velocidad vs Tiempo del Automovil')
plt.xlabel('Tiempo (s)')
plt.ylabel('Velocidad (m/s)')
plt.legend()
plt.grid(True)
plt.show()
\end{lstlisting}

\subsubsection*{Conclusión}

El automóvil sí choca con el obstáculo, y la rapidez del impacto es de 15 m/s.

\newpage

\subsection*{Enunciado}

Un automovilista viaja por una carretera recta horizontal a 90 km/h cuando de repente ve un obstáculo 50 m adelante. Tarda 0.4 s en aplicar los frenos. Determine con qué magnitud debería frenar de manera uniforme para evitar el choque.

\vspace{0.5cm}

\begin{center}
\begin{tikzpicture}
\draw[thick] (0,0) -- (10,0);
\draw[fill=blue!50] (0.5,0) rectangle (2,1);
\filldraw[black] (0.8,0) circle (0.2);
\filldraw[black] (1.7,0) circle (0.2);
\draw[fill=red!60] (8.5,0) rectangle (9,1.5);
\node[above] at (1.25, 1) {Auto};
\node[above] at (8.75, 1.5) {Obstáculo};
\draw[<->, dashed] (2, -0.5) -- (8.5, -0.5) node[midway, below] {$d = 50 \text{ m}$};
\draw[->, thick, blue] (2.2, 0.5) -- (3.5, 0.5) node[midway, above] {$v_0 = 90 \text{ km/h}$};
\draw[<-, thick, red] (4.5, 1) -- (5.8, 1) node[midway, above] {$|a_{\text{requerida}}| = ?$};
\end{tikzpicture}
\end{center}

\subsection*{Solución}

\subsubsection*{Análisis de la Distancia Disponible}

La pedagogía de la física requiere desglosar el problema en eventos cronológicos. Al igual que en el escenario anterior, el conductor posee un tiempo de reacción inherente humano durante el cual el vehículo sigue moviéndose a velocidad constante dictada por su inercia. 

Convertimos nuevamente la velocidad inicial para homogeneizar las unidades:

\begin{equation}
v_0 = 90 \text{ \frac{km}{h}} = 25 \text{ \frac{m}{s}}
\end{equation}

La distancia consumida durante el tiempo de reacción de cuatro décimas de segundo es una constante física de este problema:

\begin{equation}
d_r = 25 \cdot 0.4 = 10 \text{ m}
\end{equation}

Dado que el obstáculo se encuentra a cincuenta metros de la posición en la que el conductor lo percibe visualmente, la distancia máxima que el vehículo tiene disponible para recorrer durante el frenado sin impactar es la diferencia directa entre la distancia total y la distancia perdida en la reacción.

\begin{equation}
d_{\text{frenado}} = 50 - 10 = 40 \text{ m}
\end{equation}

\subsubsection*{Cálculo de la Aceleración Requerida}

Para que el automóvil logre evitar el choque justo a tiempo, su velocidad final debe ser cero exactamente al agotar los cuarenta metros disponibles. Planteamos un Movimiento Rectilíneo Uniformemente Variado configurando estos límites de frontera.

Recurrimos a la ecuación cinemática independiente del tiempo:

\begin{equation}
v_f^2 = v_0^2 + 2 \cdot a \cdot d_{\text{frenado}}
\end{equation}

Sustituyendo los valores conocidos para forzar la condición límite:

\begin{equation}
0 = (25)^2 + 2 \cdot a \cdot (40)
\end{equation}

Resolvemos algebraicamente para despejar la variable de aceleración requerida:

\begin{equation}
0 = 625 + 80 \cdot a
\end{equation}

\begin{equation}
-80 \cdot a = 625
\end{equation}

\begin{equation}
a = -\frac{625}{80} = -7.8125 \text{ \frac{m}{s}^2}
\end{equation}

El signo negativo indica en nuestro sistema de referencia que el vector de aceleración es opuesto al vector de velocidad, lo que físicamente concuerda con una acción de frenado continuo.

\subsubsection*{Magnitud de la Desaceleración}

El enunciado solicita de manera explícita la magnitud con la que debe frenar. La magnitud de un vector corresponde a su valor absoluto algebraico, desprovisto de información direccional referencial. Aproximando a dos cifras decimales como es convención en los problemas clásicos de cinemática:

\begin{equation}
|a| = 7.81 \text{ \frac{m}{s}^2}
\end{equation}

\subsubsection*{Conclusión}

Debería frenar de manera uniforme con una magnitud de 7.81 m/s$^2$ para evitar el choque.

\newpage

\subsection*{Enunciado}

Un policía motorizado estacionado al costado de una carretera recta horizontal observa pasar a un auto que se mueve a exceso de velocidad con rapidez constante de 28.0 m/s. 1.5 segundos después, el motorizado parte del reposo y persigue al auto acelerando uniformemente a razón de 3.0 m/s$^2$. Determine el intervalo de tiempo total, desde que el motorizado observa al auto hasta que lo alcanza.

\vspace{0.5cm}

\begin{center}
\begin{tikzpicture}
\draw[thick] (0,0) -- (10,0);
\filldraw[black] (0.5,0) circle (2pt) node[below=0.2cm] {$x = 0$ (Poste de control)};

\draw[fill=red!60] (0.5, 0.8) rectangle (1.5, 1.3);
\node[white, font=\tiny, font=\bfseries] at (1, 1.05) {Auto};
\draw[->, thick, red] (1.6, 1.05) -- (3.0, 1.05) node[midway, above] {$v_a = 28.0 \text{ m/s}$};

\draw[fill=blue!60] (0.5, 0.1) rectangle (1.5, 0.6);
\node[white, font=\tiny, font=\bfseries] at (1, 0.35) {Policía};
\draw[->, thick, blue] (1.6, 0.35) -- (2.5, 0.35) node[midway, below] {$a_p = 3.0 \text{ m/s}^2$};

\node[align=center, text=blue] at (5, -0.8) {El motorizado arranca en $t = 1.5 \text{ s}$};
\draw[dashed, gray] (0.5, -1) -- (0.5, 1.5);
\end{tikzpicture}
\end{center}

\subsection*{Solución}

\subsubsection*{Conceptos Físicos Involucrados}

De acuerdo con los principios de la física conceptual para la cinemática, la descripción de cualquier movimiento requiere establecer rigurosamente un marco de referencia espacial y un "reloj" maestro que mida el tiempo de manera unificada. En este problema interactúan dos tipos de movimiento a lo largo de una misma dimensión espacial. 

El automóvil experimenta un Movimiento Rectilíneo Uniforme dado que se traslada con rapidez constante, implicando que su aceleración es cero. En contraste, el policía parte del reposo y aumenta su velocidad gradualmente, lo que constituye un Movimiento Rectilíneo Uniformemente Variado condicionado por un retraso temporal respecto a nuestro reloj maestro.

\subsubsection*{Definición del Sistema de Referencia}

Situamos el origen de coordenadas espaciales exactamente en el punto donde se encuentra estacionado el policía. Para el dominio temporal, fijamos el instante inicial nulo en el momento exacto en que el automóvil cruza la posición del policía. 

Datos del Automóvil:
\begin{itemize}
    \item Posición inicial nula.
    \item Velocidad constante de 28.0 m/s.
    \item Tiempo de inicio de movimiento respecto al reloj maestro es cero.
\end{itemize}

Datos del Policía Motorizado:
\begin{itemize}
    \item Posición inicial nula.
    \item Velocidad inicial nula dado que parte del reposo.
    \item Aceleración constante de 3.0 m/s$^2$.
    \item Retraso temporal de 1.5 segundos. El policía solo se mueve para instantes de tiempo estrictamente mayores a 1.5 segundos.
\end{itemize}

\subsubsection*{Ecuaciones de Movimiento}

Construimos las funciones dependientes del tiempo para la posición de cada vehículo basándonos en los modelos cinemáticos estándar.

Para el automóvil, la posición obedece a una función lineal:
\begin{equation}
x_a(t) = 28.0 \cdot t
\end{equation}

Para el policía, debemos incorporar el desfase temporal. El tiempo efectivo durante el cual el policía está acelerando no es el tiempo total medido por nuestro reloj, sino el tiempo total menos el retraso de observación. La ecuación general de posición para un objeto acelerado desde el reposo se adapta de la siguiente manera:
\begin{equation}
x_p(t) = \frac{1}{2} \cdot a_p \cdot (t - t_{\text{retraso}})^2
\end{equation}

Sustituyendo los valores numéricos correspondientes al policía:
\begin{equation}
x_p(t) = \frac{1}{2} \cdot 3.0 \cdot (t - 1.5)^2 = 1.5 \cdot (t - 1.5)^2
\end{equation}

\subsubsection*{Condición de Alcance}

El concepto de "alcanzar" a otro móvil se traduce matemáticamente en la intersección de sus trayectorias espacio-temporales. Es decir, buscamos el instante de tiempo en el cual ambas funciones de posición producen el mismo valor de distancia respecto al origen. Igualamos las ecuaciones:

\begin{equation}
x_a(t) = x_p(t)
\end{equation}

\subsubsection*{Desarrollo Matemático}

Sustituimos las expresiones algebraicas en la condición de alcance:

\begin{equation}
28.0t = 1.5(t - 1.5)^2
\end{equation}

Expandimos el binomio al cuadrado en el lado derecho de la igualdad:

\begin{equation}
28.0t = 1.5(t^2 - 3.0t + 2.25)
\end{equation}

Distribuimos el factor multiplicativo en los términos del polinomio:

\begin{equation}
28.0t = 1.5t^2 - 4.5t + 3.375
\end{equation}

Transponemos todos los términos a un solo lado para estructurar una ecuación cuadrática en su forma estándar, igualada a cero:

\begin{equation}
1.5t^2 - 4.5t - 28.0t + 3.375 = 0
\end{equation}

\begin{equation}
1.5t^2 - 32.5t + 3.375 = 0
\end{equation}

Aplicamos la fórmula general para la resolución de ecuaciones cuadráticas:

\begin{equation}
t = \frac{-b \pm \sqrt{b^2 - 4ac}}{2a}
\end{equation}

Sustituyendo los coeficientes obtenidos:

\begin{equation}
t = \frac{-(-32.5) \pm \sqrt{(-32.5)^2 - 4(1.5)(3.375)}}{2(1.5)}
\end{equation}

Calculamos el valor del discriminante:

\begin{equation}
t = \frac{32.5 \pm \sqrt{1056.25 - 20.25}}{3.0}
\end{equation}

\begin{equation}
t = \frac{32.5 \pm \sqrt{1036}}{3.0}
\end{equation}

Extraemos la raíz cuadrada y evaluamos las dos posibles soluciones. La raíz de 1036 es aproximadamente 32.187.

Primera solución evaluando el signo positivo:
\begin{equation}
t_1 = \frac{32.5 + 32.187}{3.0} = \frac{64.687}{3.0} \approx 21.56 \text{ s}
\end{equation}

Segunda solución evaluando el signo negativo:
\begin{equation}
t_2 = \frac{32.5 - 32.187}{3.0} = \frac{0.313}{3.0} \approx 0.10 \text{ s}
\end{equation}

Interpretación física de las raíces: La solución de 0.10 segundos es matemáticamente válida para las parábolas abstractas, pero físicamente inaceptable para nuestro modelo, ya que el dominio de validez para la ecuación del policía exige que el tiempo sea mayor a 1.5 segundos. Por descarte lógico, la única solución real factible es 21.56 segundos.

\subsubsection*{Representación Gráfica del Problema}

La visualización de las curvas de posición clarifica la intersección matemática. La línea recta ascendente del automóvil es superada eventualmente por el crecimiento parabólico de la posición del policía.

\begin{lstlisting}[language=Python]
import matplotlib.pyplot as plt
import numpy as np

t_total = np.linspace(0, 25, 400)
x_auto = 28.0 * t_total

x_policia = np.zeros_like(t_total)
mask = t_total > 1.5
x_policia[mask] = 1.5 * (t_total[mask] - 1.5)**2

plt.figure(figsize=(9, 6))
plt.plot(t_total, x_auto, label='Auto (MRU)', color='red', linewidth=2)
plt.plot(t_total, x_policia, label='Policia (MRUV)', color='blue', linewidth=2)

plt.axvline(21.56, color='black', linestyle='dashed', alpha=0.5)
plt.axhline(603.74, color='black', linestyle='dashed', alpha=0.5)
plt.plot(21.56, 603.74, 'ko', markersize=7)

plt.text(14, 630, 'Punto de alcance\n(21.56 s)', fontsize=10, 
         bbox=dict(facecolor='white', edgecolor='black', boxstyle='round,pad=0.5'))

plt.title('Grafica de Posicion vs Tiempo de la Persecucion')
plt.xlabel('Tiempo medido desde que el auto cruza al policia (s)')
plt.ylabel('Posicion respecto al origen (m)')
plt.legend()
plt.grid(True)
plt.show()
\end{lstlisting}

\subsubsection*{Conclusión}

El intervalo de tiempo total, desde que el motorizado observa al auto hasta que lo alcanza, es de 21.56 s.

\newpage

\subsection*{Enunciado}

Un policía motorizado estacionado al costado de una carretera recta horizontal observa pasar a un auto que se mueve a exceso de velocidad con rapidez constante de 28.0 m/s. 1.5 segundos después, el motorizado parte del reposo y persigue al auto acelerando uniformemente a razón de 3.0 m/s$^2$. Determine la distancia recorrida por el motorizado hasta alcanzar al auto.

\vspace{0.5cm}

\begin{center}
\begin{tikzpicture}
\draw[thick] (0,0) -- (10,0);
\filldraw[black] (0.5,0) circle (2pt);

\draw[fill=blue!60] (0.5, 0.1) rectangle (1.5, 0.6);
\node[white, font=\tiny, font=\bfseries] at (1, 0.35) {Policía};

\draw[fill=blue!60, opacity=0.3] (8.5, 0.1) rectangle (9.5, 0.6);
\draw[fill=red!60, opacity=0.3] (8.5, 0.8) rectangle (9.5, 1.3);

\draw[<->, dashed, thick, blue] (1.0, -0.6) -- (9.0, -0.6) node[midway, below] {$d = ?$ (Distancia de alcance)};
\draw[dashed, gray] (1.0, 0) -- (1.0, -0.8);
\draw[dashed, gray] (9.0, 0) -- (9.0, -0.8);

\node[align=center, font=\small] at (9.0, 1.6) {Punto de\\Encuentro};
\end{tikzpicture}
\end{center}

\subsection*{Solución}

\subsubsection*{Conceptos Físicos Involucrados}

La distancia total recorrida desde un punto de partida hasta un destino en una trayectoria rectilínea unidireccional coincide numéricamente con el valor de la posición final del móvil, asumiendo que el marco de referencia tiene su origen en el lugar de inicio. En un problema de persecución donde dos entidades se interceptan, la coordenada espacial evaluada en el tiempo del encuentro debe ser idéntica para ambas.

Al existir esta identidad de posición, pedagógicamente podemos seleccionar la ecuación matemática que resulte más sencilla de procesar algebraicamente para minimizar el riesgo de errores en el cálculo.

\subsubsection*{Recuperación de Datos Temporales}

Del análisis del sistema de ecuaciones resuelto para determinar la condición temporal de intersección, extraemos el instante preciso en el cual el evento toma lugar en nuestra línea del tiempo principal.

\begin{equation}
t_{\text{alcance}} \approx 21.562 \text{ s}
\end{equation}

Es imperativo utilizar una cantidad razonable de cifras decimales significativas procedentes del cálculo de la raíz cuadrática previa para asegurar la exactitud del producto final.

\subsubsection*{Desarrollo Matemático}

Como se justificó en los conceptos físicos, evaluaremos la ecuación de posición del automóvil, ya que presenta una estructura lineal, carece de desfasajes temporales y depende únicamente de un producto escalar directo. 

La ecuación de posición del auto es:
\begin{equation}
x_a(t) = 28.0 \cdot t
\end{equation}

Sustituimos el tiempo de alcance hallado en la función:
\begin{equation}
x_a(21.562) = 28.0 \cdot (21.562)
\end{equation}

Realizamos la multiplicación aritmética:
\begin{equation}
x_a(21.562) = 603.736 \text{ m}
\end{equation}

Para garantizar una total consistencia interna, validamos este resultado utilizando la ecuación independiente del policía. Recordando que el policía aceleró durante un tiempo efectivo menor al tiempo total debido a su retraso de 1.5 segundos:

\begin{equation}
t_{\text{efectivo}} = 21.562 - 1.5 = 20.062 \text{ s}
\end{equation}

Evaluamos la posición del policía con su tiempo efectivo de aceleración:
\begin{equation}
x_p = \frac{1}{2} \cdot a_p \cdot (t_{\text{efectivo}})^2
\end{equation}

\begin{equation}
x_p = \frac{1}{2} \cdot 3.0 \cdot (20.062)^2 = 1.5 \cdot 402.484 = 603.726 \text{ m}
\end{equation}

La minúscula discrepancia en la tercera cifra decimal es una consecuencia directa del truncamiento de la raíz irracional. Para fines de estandarización académica y pedagógica, redondeamos el resultado a dos cifras decimales de precisión, confirmando la convergencia de ambos modelos a la misma distancia kilométrica fraccional.

\subsubsection*{Conclusión}

La distancia recorrida por el motorizado hasta alcanzar al auto es de 603.74 m.

\newpage

\subsection*{Enunciado}

¿Qué distancia recorrerá una madre que camina a 3 m/s hacia su hijo, quien camina hacia ella a 2 m/s, si inicialmente están separados por 50 metros? Ambos se desplazan en direcciones contrarias y en línea recta.

\begin{itemize}
    \item[a)] 20 m
    \item[b)] \textbf{30 m}
    \item[c)] 50 m
    \item[d)] 60 m
    \item[e)] 80 m
\end{itemize}

\vspace{0.5cm}

\begin{center}
\begin{tikzpicture}
\draw[step=0.5cm, blue!20, very thin] (-1,-1) grid (7,2);
\draw[->, thick, green!60!black] (0,0) -- (6,0) node[right, black] {\large x (+)};

\draw[thick, fill=white] (-0.3, 0) rectangle (0.3, 0.6);
\node[above] at (0, 0.7) {\Large M};
\node[below] at (0, -0.2) {\Large 0};

\draw[thick, fill=white] (4.7, 0) rectangle (5.3, 0.6);
\node[above] at (5, 0.7) {\Large N};
\node[below] at (5, -0.2) {\Large +50};
\end{tikzpicture}
\end{center}

\subsection*{Solución}

\subsubsection*{Conceptos Físicos Involucrados}

Desde la perspectiva de la física conceptual, el movimiento de los cuerpos debe analizarse respecto a un marco de referencia espacial. Dado que la madre y el hijo se mueven a lo largo de una línea recta cubriendo distancias iguales en tiempos iguales, ambos experimentan un Movimiento Rectilíneo Uniforme. En este régimen cinemático, la velocidad es constante, la aceleración es nula y la posición depende directamente del tiempo transcurrido, la velocidad y la posición inicial.

Para modelar matemáticamente este encuentro, es imprescindible establecer una convención de signos vectorial que diferencie la dirección en la que se mueve cada persona a lo largo del eje unidimensional.

\subsubsection*{Definición del Sistema de Referencia}

Situamos el origen de nuestro sistema de coordenadas espaciales en la posición inicial de la madre, a la cual denotaremos con la letra M. Por consiguiente, la posición inicial del hijo, denotado con la letra N, se ubica en la coordenada positiva de cincuenta metros.

La madre avanza en el sentido positivo del eje espacial, por lo que su vector velocidad tiene un valor positivo. El hijo se mueve en dirección contraria, acercándose al origen para encontrarse con su madre, lo que exige asignar un signo negativo a su velocidad en nuestras ecuaciones.

Datos iniciales para la Madre (M):
\begin{itemize}
    \item Posición inicial $x_{M0} = 0 \text{ m}$
    \item Velocidad $v_M = 3 \text{ m/s}$
\end{itemize}

Datos iniciales para el Hijo (N):
\begin{itemize}
    \item Posición inicial $x_{N0} = 50 \text{ m}$
    \item Velocidad $v_N = -2 \text{ m/s}$
\end{itemize}

\subsubsection*{Ecuaciones de Movimiento}

Utilizamos la ecuación fundamental de la cinemática para velocidad constante, donde la posición en función del tiempo es igual a la posición inicial más el producto de la velocidad por el tiempo.

Para la Madre:
\begin{equation}
x_M(t) = 0 + 3t = 3t
\end{equation}

Para el Hijo:
\begin{equation}
x_N(t) = 50 - 2t
\end{equation}

\subsubsection*{Condición de Encuentro}

El instante temporal en el que ambas personas se cruzan se define matemáticamente como el momento en el que la coordenada espacial de la madre es exactamente igual a la coordenada espacial del hijo. Igualamos sus funciones de posición:

\begin{equation}
x_M(t) = x_N(t)
\end{equation}

\subsubsection*{Desarrollo Matemático}

Sustituimos las expresiones cinemáticas estructuradas previamente para encontrar el valor de la incógnita del tiempo:

\begin{equation}
3t = 50 - 2t
\end{equation}

Agrupamos los términos dependientes del tiempo en el lado izquierdo de la ecuación sumando el término negativo de la derecha en ambos miembros:

\begin{equation}
3t + 2t = 50
\end{equation}

\begin{equation}
5t = 50
\end{equation}

Despejamos el tiempo dividiendo la ecuación por cinco:

\begin{equation}
t = 10 \text{ s}
\end{equation}

El encuentro se produce diez segundos después de que ambos inician su marcha. El problema solicita la distancia que recorrerá la madre. Dado que ella parte desde el origen del sistema de coordenadas, su posición en el instante de encuentro equivale exactamente a la distancia que ha recorrido. Evaluamos la función de posición de la madre para este instante temporal:

\begin{equation}
x_M(10) = 3 \cdot (10)
\end{equation}

\begin{equation}
x_M(10) = 30 \text{ m}
\end{equation}

Como validación pedagógica, podemos corroborar que las posiciones coinciden evaluando el desplazamiento del hijo. En diez segundos, el hijo recorre veinte metros hacia la izquierda, y partiendo desde la posición de cincuenta metros, se ubicará también en la coordenada espacial de treinta metros.

\subsubsection*{Representación Gráfica del Encuentro}

El modelado gráfico de las rectas de posición facilita la comprensión conceptual de los cruces cinemáticos. A continuación, se presenta la implementación computacional en Python para generar el diagrama de la intersección que valida los resultados analíticos.

\begin{lstlisting}[language=Python]
import matplotlib.pyplot as plt
import numpy as np

t_eval = np.linspace(0, 15, 100)
pos_madre = 3 * t_eval
pos_hijo = 50 - 2 * t_eval

plt.figure(figsize=(8, 5))
plt.plot(t_eval, pos_madre, label='Madre (M)', color='blue', linewidth=2)
plt.plot(t_eval, pos_hijo, label='Hijo (N)', color='red', linewidth=2)

plt.axvline(10, color='black', linestyle='dashed', alpha=0.5)
plt.axhline(30, color='black', linestyle='dashed', alpha=0.5)
plt.plot(10, 30, 'go', markersize=8)

plt.text(10.5, 32, 'Punto de encuentro\n(10 s, 30 m)', fontsize=10)

plt.title('Posicion vs Tiempo para la Madre y el Hijo')
plt.xlabel('Tiempo (s)')
plt.ylabel('Posicion (m)')
plt.legend()
plt.grid(True)
plt.show()
\end{lstlisting}

\subsubsection*{Conclusión}

La distancia que recorrerá la madre hasta encontrarse con su hijo es de 30 m (opción b).

\newpage

\subsection*{Enunciado}

Un automóvil puede acelerar desde 0 hasta 20 m/s en 5 segundos. Entonces, el tiempo mínimo que tarda en recorrer 200 m partiendo desde el reposo es:

\begin{itemize}
    \item[a)] 5 s
    \item[b)] $\sqrt{5} \text{ s}$
    \item[c)] \textbf{10 s}
    \item[d)] $\sqrt{10} \text{ s}$
    \item[e)] 20 s
\end{itemize}

\vspace{0.5cm}

\begin{center}
\begin{tikzpicture}
\draw[thick] (0,0) -- (10,0);

\draw[fill=blue!30, rounded corners=2pt] (0.5,0.2) rectangle (2.0,1.0);
\draw[fill=blue!20] (1.5,0.6) rectangle (1.9,0.9);
\filldraw[black] (0.8,0.2) circle (0.2);
\filldraw[black] (1.7,0.2) circle (0.2);
\node[above] at (1.25, 1.0) {$v_0 = 0 \text{ m/s}$};

\draw[->, thick, red] (2.2, 0.6) -- (3.5, 0.6) node[midway, above] {$a = ?$};

\draw[<->, dashed, thick, black!70] (1.25, -0.6) -- (8.75, -0.6) node[midway, below] {$d = 200 \text{ m}$};
\draw[dashed, gray] (1.25, 0) -- (1.25, -0.8);
\draw[dashed, gray] (8.75, 0) -- (8.75, -0.8);

\draw[fill=blue!30, opacity=0.3, rounded corners=2pt] (8.0,0.2) rectangle (9.5,1.0);
\filldraw[black, opacity=0.3] (8.3,0.2) circle (0.2);
\filldraw[black, opacity=0.3] (9.2,0.2) circle (0.2);
\node[above, text=gray] at (8.75, 1.0) {$t = ?$};

\end{tikzpicture}
\end{center}

\subsection*{Solución}

\subsubsection*{Conceptos Físicos Involucrados}

En el marco de la física conceptual, la aceleración se define como la tasa a la cual un objeto cambia su velocidad en el tiempo. Si un vehículo requiere cubrir una distancia específica en el menor tiempo posible, debe maximizar su aceleración desde el instante inicial y mantenerla constante durante todo el trayecto. 

Al mantener una aceleración uniforme a lo largo del tiempo, el automóvil describe un Movimiento Rectilíneo Uniformemente Variado. La posición del vehículo en cualquier instante dependerá cuadráticamente del tiempo, debido a que la velocidad está aumentando continuamente.

\subsubsection*{Determinación de la Aceleración}

El enunciado proporciona los datos iniciales necesarios para descubrir la capacidad de aceleración del automóvil. Sabemos que cambia su estado de reposo a una velocidad de veinte metros por segundo en un lapso de cinco segundos.

Utilizamos la definición matemática de la aceleración constante, que es el cociente entre la variación de la velocidad y el intervalo de tiempo:

\begin{equation}
a = \frac{\Delta v}{\Delta t} = \frac{v_f - v_0}{\Delta t}
\end{equation}

Sustituimos los valores numéricos correspondientes a la prueba de rendimiento del vehículo:

\begin{equation}
a = \frac{20 \text{ m/s} - 0 \text{ m/s}}{5 \text{ s}}
\end{equation}

\begin{equation}
a = 4 \text{ \frac{m}{s}^2}
\end{equation}

Esto significa que, en cada segundo que transcurre, la velocidad del automóvil se incrementa en cuatro metros por segundo.

\subsubsection*{Relación entre Posición y Tiempo}

Una vez conocida la capacidad de aceleración del automóvil, podemos proyectar su posición futura. La ecuación cinemática que relaciona el desplazamiento con el tiempo para un cuerpo que parte del reposo es:

\begin{equation}
x(t) = x_0 + v_0 t + \frac{1}{2} a t^2
\end{equation}

Fijando nuestro sistema de referencia en la posición inicial del automóvil nula y recordando que parte del reposo, los dos primeros términos de la ecuación se anulan. La expresión se simplifica a:

\begin{equation}
x(t) = \frac{1}{2} a t^2
\end{equation}

\subsubsection*{Desarrollo Matemático}

Nuestro objetivo es encontrar el tiempo exacto en el cual la posición del automóvil alcanza los doscientos metros. Sustituimos la distancia objetivo y la aceleración calculada en nuestra ecuación simplificada:

\begin{equation}
200 = \frac{1}{2} \cdot (4) \cdot t^2
\end{equation}

Realizamos la división de la constante:

\begin{equation}
200 = 2 \cdot t^2
\end{equation}

Despejamos el término cuadrático dividiendo ambos lados de la ecuación por dos:

\begin{equation}
t^2 = \frac{200}{2}
\end{equation}

\begin{equation}
t^2 = 100
\end{equation}

Finalmente, extraemos la raíz cuadrada para obtener el valor del tiempo. En este contexto físico, descartamos la raíz negativa, pues el tiempo transcurre de manera estrictamente positiva:

\begin{equation}
t = \sqrt{100}
\end{equation}

\begin{equation}
t = 10 \text{ s}
\end{equation}

\subsubsection*{Representación Gráfica del Movimiento Acelerado}

Para consolidar el aprendizaje, graficar la función de posición en un movimiento uniformemente variado revela una curva parabólica característica, ilustrando cómo el vehículo recorre cada vez mayor distancia en el mismo intervalo de tiempo debido al aumento de velocidad.

\begin{lstlisting}[language=Python]
import matplotlib.pyplot as plt
import numpy as np

t = np.linspace(0, 12, 100)
x = 0.5 * 4 * t**2

plt.figure(figsize=(8, 5))
plt.plot(t, x, label='Posicion $x(t) = 2t^2$', color='blue', linewidth=2)

plt.axvline(10, color='black', linestyle='dashed', alpha=0.5)
plt.axhline(200, color='black', linestyle='dashed', alpha=0.5)
plt.plot(10, 200, 'ro', markersize=8)

plt.text(6.5, 210, 'Meta alcanzada\n(10 s, 200 m)', fontsize=11)

plt.title('Distancia vs Tiempo del Automovil')
plt.xlabel('Tiempo (s)')
plt.ylabel('Distancia Recorrida (m)')
plt.legend()
plt.grid(True)
plt.show()
\end{lstlisting}

\subsubsection*{Conclusión}

El tiempo mínimo que tarda en recorrer 200 m partiendo desde el reposo es de 10 s (opción c).

\newpage

\subsection*{Enunciado}

Un automóvil que se mueve a lo largo de una carretera recta horizontal con una rapidez de 30 m/s empieza a frenar uniformemente a un ritmo de 4 m/s$^2$ en un punto ubicado a 100 m de un obstáculo. Entonces es correcto afirmar que el automóvil:

\begin{itemize}
    \item[a)] se detiene 20 m antes de chocar con el obstáculo
    \item[b)] choca con el obstáculo con una rapidez de 20 m/s
    \item[c)] \textbf{choca con el obstáculo con una rapidez de 10 m/s}
    \item[d)] se detiene 10 m antes de chocar con el obstáculo
    \item[e)] se detiene justamente antes de chocar con el obstáculo
\end{itemize}

\vspace{0.5cm}

\begin{center}
\begin{tikzpicture}
\draw[thick] (0,0) -- (12,0);
\draw[fill=blue!50] (0.5,0) rectangle (2,1);
\filldraw[black] (0.8,0) circle (0.2);
\filldraw[black] (1.7,0) circle (0.2);
\draw[fill=red!60] (10.5,0) rectangle (11,1.5);
\node[above] at (1.25, 1) {Auto};
\node[above] at (10.75, 1.5) {Obstáculo};
\draw[<->, dashed] (2, -0.5) -- (10.5, -0.5) node[midway, below] {$d = 100 \text{ m}$};
\draw[->, thick, blue] (2.2, 0.5) -- (3.5, 0.5) node[midway, above] {$v_0 = 30 \text{ m/s}$};
\draw[<-, thick, red] (1.25, 1.5) -- (2.5, 1.5) node[midway, above] {$a = 4 \text{ m/s}^2$};
\end{tikzpicture}
\end{center}

\subsection*{Solución}

\subsubsection*{Conceptos Físicos Involucrados}

Desde la perspectiva de la física conceptual, la aceleración es una magnitud vectorial que mide el cambio de la velocidad en el tiempo. Cuando un vehículo frena, experimenta una aceleración en dirección opuesta a su movimiento (una desaceleración). Esto configura un Movimiento Rectilíneo Uniformemente Variado donde la energía cinética del objeto se disipa progresivamente. 

Para resolver problemas de colisiones potenciales, el enfoque pedagógico estándar dicta primero determinar si el espacio físico disponible es suficiente para que el vehículo reduzca su velocidad a cero antes de llegar a la posición del obstáculo.

\subsubsection*{Análisis de la Distancia de Parada Teórica}

Primero, calcularemos la distancia que el automóvil necesitaría para detenerse por completo si no existiera el obstáculo. En este escenario hipotético, la velocidad final es nula. 

Utilizamos la ecuación cinemática que relaciona las velocidades al cuadrado, la aceleración y el desplazamiento, ya que no disponemos del tiempo transcurrido:

\begin{equation}
v_f^2 = v_0^2 + 2 \cdot a \cdot \Delta x
\end{equation}

Establecemos nuestro sistema de referencia asumiendo que el automóvil se mueve en la dirección positiva. Por lo tanto, la velocidad inicial es positiva y la aceleración de frenado debe ser matemáticamente negativa.

\begin{itemize}
    \item Velocidad inicial $v_0 = 30 \text{ m/s}$
    \item Velocidad final $v_f = 0 \text{ m/s}$
    \item Aceleración $a = -4 \text{ m/s}^2$
\end{itemize}

Sustituimos estos valores en nuestra ecuación:

\begin{equation}
0 = (30)^2 + 2 \cdot (-4) \cdot \Delta x_{\text{parada}}
\end{equation}

\begin{equation}
0 = 900 - 8 \cdot \Delta x_{\text{parada}}
\end{equation}

Despejamos la distancia de parada aislando el término espacial:

\begin{equation}
8 \cdot \Delta x_{\text{parada}} = 900
\end{equation}

\begin{equation}
\Delta x_{\text{parada}} = \frac{900}{8} = 112.5 \text{ m}
\end{equation}

\subsubsection*{Determinación del Choque}

Al comparar la distancia requerida para detenerse con la distancia disponible en la realidad del problema, observamos lo siguiente:

\begin{equation}
112.5 \text{ m} > 100 \text{ m}
\end{equation}

Dado que el automóvil requiere más de cien metros para detenerse, es físicamente inevitable que colisione con el obstáculo. Esto descarta las opciones que sugieren que el vehículo se detiene antes o justo a tiempo.

\subsubsection*{Cálculo de la Rapidez de Impacto}

Sabiendo que el choque ocurre, debemos averiguar qué velocidad posee el vehículo exactamente cuando ha recorrido los cien metros que lo separan del obstáculo. Reutilizamos la misma ecuación cinemática, pero ahora nuestra incógnita es la velocidad final y el desplazamiento está fijado en la posición del obstáculo.

\begin{equation}
v_{\text{impacto}}^2 = v_0^2 + 2 \cdot a \cdot d_{\text{obstáculo}}
\end{equation}

Sustituimos los datos correspondientes al momento del choque:

\begin{equation}
v_{\text{impacto}}^2 = (30)^2 + 2 \cdot (-4) \cdot (100)
\end{equation}

Realizamos las operaciones aritméticas:

\begin{equation}
v_{\text{impacto}}^2 = 900 - 800
\end{equation}

\begin{equation}
v_{\text{impacto}}^2 = 100
\end{equation}

Extraemos la raíz cuadrada para obtener la magnitud de la velocidad, es decir, la rapidez:

\begin{equation}
v_{\text{impacto}} = \sqrt{100} = 10 \text{ m/s}
\end{equation}

\subsubsection*{Análisis Gráfico de Velocidad vs Posición}

Para consolidar la comprensión, podemos graficar cómo disminuye la velocidad en función de la distancia recorrida. El código en Python que se presenta a continuación modela la curva de desaceleración y resalta el punto exacto de la interrupción del movimiento debido a la colisión.

\begin{lstlisting}[language=Python]
import matplotlib.pyplot as plt
import numpy as np

distancia_teorica = np.linspace(0, 112.5, 200)
velocidad = np.sqrt(900 - 8 * distancia_teorica)

plt.figure(figsize=(8, 5))
plt.plot(distancia_teorica, velocidad, label='Curva de frenado', color='blue', linewidth=2)

plt.axvline(100, color='red', linestyle='dashed', linewidth=1.5, label='Posicion del obstaculo')
plt.axhline(0, color='black', linewidth=0.5)

plt.plot(100, 10, 'ko', markersize=8)
plt.text(85, 12, 'Impacto (100 m, 10 m/s)', fontsize=10, weight='bold')

plt.fill_between(distancia_teorica, velocidad, color='blue', alpha=0.1, where=(distancia_teorica<=100))

plt.title('Velocidad en funcion de la Distancia Recorrida')
plt.xlabel('Distancia (m)')
plt.ylabel('Velocidad (m/s)')
plt.legend()
plt.grid(True)
plt.show()
\end{lstlisting}

\subsubsection*{Conclusión}

El automóvil choca con el obstáculo con una rapidez de 10 m/s (opción c).

\newpage

\subsection*{Enunciado}

Un policía motorizado estacionado al costado de una carretera recta horizontal observa pasar a un auto que se mueve con rapidez constante a 28.0 m/s. Después de 1.5 segundos, el motorizado parte del reposo y persigue al auto acelerando uniformemente a razón de 3.0 m/s$^2$. Determine el intervalo de tiempo total, desde que el motorizado observa al auto hasta que lo alcanza.

\vspace{0.5cm}

\begin{center}
\begin{tikzpicture}
\draw[step=0.5cm, cyan!20, very thin] (-1,-1.5) grid (10,1.5);
\draw[->, thick, green!60!black] (-1,0) -- (10,0);

\draw[thick, fill=white] (0,-0.4) rectangle (0.8,0.4);
\node[above] at (0.4, 0.5) {\large P, $t = 1.5 \text{ s}$};
\node[below] at (0.4, -0.5) {\large 0};
\draw[->, thick, red] (0.8, 0) -- (2.2, 0) node[midway, below] {\large $\vec{a}$};

\draw[thick, fill=white] (3.5,-0.4) rectangle (4.3,0.4);
\node[above] at (3.9, 0.5) {\large A, $t = 1.5 \text{ s}$};
\draw[->, thick, blue] (4.3, 0) -- (5.8, 0) node[midway, below] {\large $\vec{v}_A$};

\draw[thick, fill=white] (8,-0.4) rectangle (8.8,0.4);
\draw[thick, fill=white] (8.8,-0.4) rectangle (9.6,0.4);
\node[above] at (8.8, 0.5) {\large P, A, $t = ?$};
\node[below] at (8.8, -0.5) {\large $x_P = x_A$};
\end{tikzpicture}
\end{center}

\subsection*{Solución}

\subsubsection*{Conceptos Físicos Involucrados}

Acorde a la pedagogía de la física conceptual, el análisis de sistemas cinemáticos con múltiples móviles requiere establecer un marco de referencia espacial y temporal absoluto. El automóvil, al mantener una velocidad inalterable, describe un Movimiento Rectilíneo Uniforme. El policía, al partir del reposo y aumentar su velocidad, describe un Movimiento Rectilíneo Uniformemente Variado. 

El aspecto fundamental de este problema es la asimetría temporal. El reloj maestro del sistema se inicia en $t = 0$ cuando el automóvil cruza el origen. Sin embargo, el policía permanece en reposo durante los primeros 1.5 segundos, por lo que su intervalo efectivo de aceleración es estrictamente menor al tiempo transcurrido para el auto.

\subsubsection*{Definición del Sistema de Referencia}

Situamos el origen espacial en la posición del policía estacionado. El instante $t = 0$ corresponde al momento en que el auto cruza dicho origen.

Para el Automóvil (A):
\begin{itemize}
    \item Posición inicial $x_{A0} = 0 \text{ m}$
    \item Velocidad $v_A = 28.0 \text{ m/s}$
    \item Tiempo de inicio $t_{0A} = 0 \text{ s}$
\end{itemize}

Para el Policía (P):
\begin{itemize}
    \item Posición inicial $x_{P0} = 0 \text{ m}$
    \item Velocidad inicial $v_{P0} = 0 \text{ m/s}$
    \item Aceleración $a_P = 3.0 \text{ m/s}^2$
    \item Retraso temporal de inicio $t_{retraso} = 1.5 \text{ s}$
\end{itemize}

\subsubsection*{Ecuaciones de Movimiento}

Construimos las funciones de posición adaptadas a las condiciones de cada vehículo. 

Para el automóvil, la relación es directa y lineal:
\begin{equation}
x_A(t) = 28.0t
\end{equation}

Para el policía, la función cuadrática de posición debe incorporar el desfase del tiempo, ya que solo comienza a desplazarse cuando el tiempo total supera los 1.5 segundos:
\begin{equation}
x_P(t) = v_{P0}(t - 1.5) + \frac{1}{2} a_P (t - 1.5)^2
\end{equation}

Como parte del reposo, el término de la velocidad inicial se anula:
\begin{equation}
x_P(t) = \frac{1}{2} (3.0) (t - 1.5)^2 = 1.5(t - 1.5)^2
\end{equation}

\subsubsection*{Condición de Alcance}

Para que el policía intercepte al automóvil, ambos deben ocupar la misma coordenada espacial en el mismo instante de tiempo. Se igualan ambas ecuaciones de trayectoria:
\begin{equation}
x_A = x_P
\end{equation}

\subsubsection*{Desarrollo Matemático}

Sustituimos las ecuaciones y desarrollamos el álgebra correspondiente:
\begin{equation}
28.0t = 1.5(t - 1.5)^2
\end{equation}

Expandimos el binomio cuadrado perfecto:
\begin{equation}
28.0t = 1.5(t^2 - 3.0t + 2.25)
\end{equation}

Distribuimos el factor constante:
\begin{equation}
28.0t = 1.5t^2 - 4.5t + 3.375
\end{equation}

Agrupamos todos los términos para formar una ecuación cuadrática igualada a cero:
\begin{equation}
1.5t^2 - 4.5t - 28.0t + 3.375 = 0
\end{equation}

\begin{equation}
1.5t^2 - 32.5t + 3.375 = 0
\end{equation}

Aplicamos la fórmula resolvente cuadrática:
\begin{equation}
t = \frac{-(-32.5) \pm \sqrt{(-32.5)^2 - 4(1.5)(3.375)}}{2(1.5)}
\end{equation}

\begin{equation}
t = \frac{32.5 \pm \sqrt{1056.25 - 20.25}}{3.0}
\end{equation}

\begin{equation}
t = \frac{32.5 \pm \sqrt{1036}}{3.0}
\end{equation}

Considerando que la raíz cuadrada de 1036 es aproximadamente 32.187, calculamos la raíz con significado físico (la positiva y mayor al tiempo de retraso):
\begin{equation}
t = \frac{32.5 + 32.187}{3.0} \approx 21.56 \text{ s}
\end{equation}

\subsubsection*{Representación Gráfica de la Intersección}

La siguiente gráfica programada en Python ilustra el momento exacto en que la curva parabólica del movimiento acelerado supera a la recta constante, confirmando visualmente la solución analítica.

\begin{lstlisting}[language=Python]
import matplotlib.pyplot as plt
import numpy as np

t = np.linspace(0, 25, 200)
x_A = 28.0 * t
x_P = np.where(t > 1.5, 1.5 * (t - 1.5)**2, 0)

plt.figure(figsize=(8, 5))
plt.plot(t, x_A, label='Auto $x_A(t)$', color='blue', linewidth=2)
plt.plot(t, x_P, label='Policia $x_P(t)$', color='red', linewidth=2)

plt.axvline(21.56, color='black', linestyle='dashed', alpha=0.5)
plt.plot(21.56, 28 * 21.56, 'ko', markersize=6)

plt.title('Trayectorias de Persecucion')
plt.xlabel('Tiempo total (s)')
plt.ylabel('Posicion desde el origen (m)')
plt.legend()
plt.grid(True)
plt.show()
\end{lstlisting}

\subsubsection*{Conclusión}

El intervalo de tiempo total, desde que el motorizado observa al auto hasta que lo alcanza, es de 21.56 s.

\newpage

\subsection*{Enunciado}

Un policía motorizado estacionado al costado de una carretera recta horizontal observa pasar a un auto que se mueve con rapidez constante a 28.0 m/s. Después de 1.5 segundos, el motorizado parte del reposo y persigue al auto acelerando uniformemente a razón de 3.0 m/s$^2$. Determine la distancia recorrida por el motorizado hasta alcanzar al auto.

\vspace{0.5cm}

\begin{center}
\begin{tikzpicture}
\draw[step=0.5cm, cyan!20, very thin] (-1,-1) grid (10,1);
\draw[->, thick, green!60!black] (-1,0) -- (10,0);

\draw[thick, fill=white] (0,-0.4) rectangle (0.8,0.4);
\node[above] at (0.4, 0.5) {\large P ($t=0$)};

\draw[thick, fill=white] (8,-0.4) rectangle (8.8,0.4);
\draw[thick, fill=white] (8.8,-0.4) rectangle (9.6,0.4);
\node[above] at (8.8, 0.5) {\large Alcance};

\draw[<->, dashed, thick, black!80] (0.4, -0.6) -- (8.8, -0.6) node[midway, below, fill=cyan!20] {\large $d = ?$ (Distancia total)};
\end{tikzpicture}
\end{center}

\subsection*{Solución}

\subsubsection*{Conceptos Físicos Involucrados}

La distancia total cubierta durante la persecución es idéntica a la coordenada espacial en el instante de la captura, debido a que el sistema de referencia tiene su origen fijado en el punto exacto donde inicia la métrica. 

Por el principio de identidad posicional en un encuentro cinemático, podemos evaluar la distancia requerida sustituyendo el tiempo de intersección en la ecuación de movimiento de cualquiera de los dos vehículos. Pedagógicamente, es recomendable utilizar la función de comportamiento lineal por su menor complejidad algebraica.

\subsubsection*{Recuperación del Tiempo de Intersección}

Del análisis previo del evento, el tiempo total medido desde la observación inicial hasta el instante en que el policía iguala la posición del automóvil es:
\begin{equation}
t_{\text{alcance}} = 21.562 \text{ s}
\end{equation}

Conservar tres cifras decimales en este parámetro intermedio es una práctica recomendada en física computacional para evitar la propagación del error por truncamiento hacia la respuesta final.

\subsubsection*{Desarrollo Matemático}

Utilizamos la ecuación de posición del automóvil, la cual establece que la distancia es el producto de la rapidez constante por el tiempo total transcurrido.

\begin{equation}
x_A(t) = 28.0 \cdot t
\end{equation}

Evaluamos la función para el instante de alcance:
\begin{equation}
x_A(21.562) = 28.0 \cdot (21.562)
\end{equation}

\begin{equation}
x_A(21.562) = 603.736 \text{ m}
\end{equation}

Para verificar la robustez del modelo, comprobamos evaluando el lado del motorizado. El tiempo que el motorizado estuvo realmente en movimiento es $21.562 - 1.5 = 20.062$ segundos.

\begin{equation}
x_P = \frac{1}{2} \cdot a_P \cdot (t - 1.5)^2
\end{equation}

\begin{equation}
x_P = 1.5 \cdot (20.062)^2 = 1.5 \cdot 402.484 = 603.726 \text{ m}
\end{equation}

La convergencia de ambos valores (con una variación microscópica residual de los decimales de la raíz) valida formalmente el cálculo. Redondeando convencionalmente a dos cifras decimales para la presentación final.

\subsubsection*{Conclusión}

La distancia recorrida por el motorizado hasta alcanzar al auto es de 603.74 m.

\newpage

\subsection*{Enunciado}

Durante una prueba de manejo, un automóvil se mueve a lo largo de una carretera recta y horizontal. El vehículo parte del reposo y acelera uniformemente a una tasa de $2 \text{ m/s}^2$ hasta recorrer $100 \text{ m}$. Luego, mantiene su rapidez constante durante otros $100 \text{ m}$. Finalmente, acelera nuevamente a una tasa de $2 \text{ m/s}^2$ hasta recorrer $100 \text{ m}$ adicionales. Determine cuál es la rapidez del automóvil al finalizar la prueba.

\vspace{0.5cm}

\begin{center}
\begin{tikzpicture}
\draw[thick, ->] (0,0) -- (12,0) node[right] {$x$};

\filldraw[black] (0,0) circle (2pt) node[below=0.2cm] {$0 \text{ m}$};
\filldraw[black] (3.5,0) circle (2pt) node[below=0.2cm] {$100 \text{ m}$};
\filldraw[black] (7,0) circle (2pt) node[below=0.2cm] {$200 \text{ m}$};
\filldraw[black] (10.5,0) circle (2pt) node[below=0.2cm] {$300 \text{ m}$};

\draw[->, thick, red] (0.5, 0.5) -- (2, 0.5) node[midway, above] {$a_1 = 2 \text{ m/s}^2$};
\draw[->, thick, blue] (4.5, 0.5) -- (6, 0.5) node[midway, above] {$v_2 = \text{cte}$};
\draw[->, thick, red] (8, 0.5) -- (9.5, 0.5) node[midway, above] {$a_3 = 2 \text{ m/s}^2$};

\node[above] at (0, 0.2) {$v_0=0$};
\node[above] at (10.5, 0.2) {$v_f=?$};
\end{tikzpicture}
\end{center}

\subsection*{Solución}

\subsubsection*{Conceptos Físicos Involucrados}

En el marco de la física conceptual, los movimientos complejos pueden entenderse fragmentándolos en intervalos cinemáticos más simples. El automóvil experimenta tres fases de movimiento distintas: una fase de aceleración inicial (Movimiento Rectilíneo Uniformemente Variado), una fase de inercia o velocidad constante (Movimiento Rectilíneo Uniforme) y una fase final de aceleración (MRUV). 

La condición de frontera fundamental en este tipo de modelado matemático es que el estado final de un tramo espacial (su velocidad y posición final) se convierte exactamente en el estado inicial del tramo inmediatamente posterior.

\subsubsection*{Análisis del Primer Tramo (Aceleración Inicial)}

El vehículo inicia su trayectoria desde el estado de reposo absoluto. Extraemos los parámetros físicos de este primer intervalo:
\begin{itemize}
    \item Velocidad inicial $v_0 = 0 \text{ m/s}$
    \item Aceleración constante $a_1 = 2 \text{ m/s}^2$
    \item Desplazamiento del tramo $\Delta x_1 = 100 \text{ m}$
\end{itemize}

Utilizamos la ecuación cinemática independiente del tiempo para calcular la velocidad que alcanza el vehículo justo al cruzar la marca de los primeros cien metros:
\begin{equation}
v_1^2 = v_0^2 + 2 \cdot a_1 \cdot \Delta x_1
\end{equation}

Sustituyendo los valores conocidos:
\begin{equation}
v_1^2 = 0^2 + 2 \cdot (2) \cdot (100)
\end{equation}

\begin{equation}
v_1^2 = 400
\end{equation}

\begin{equation}
v_1 = \sqrt{400} = 20 \text{ m/s}
\end{equation}

\subsubsection*{Análisis del Segundo Tramo (Velocidad Constante)}

Según la primera ley de Newton, al suspenderse la aceleración (fuerza neta nula), el vehículo mantiene su estado de movimiento constante por inercia a lo largo de los siguientes cien metros.
\begin{itemize}
    \item Aceleración $a_2 = 0 \text{ m/s}^2$
    \item Desplazamiento del tramo $\Delta x_2 = 100 \text{ m}$
\end{itemize}

Dado que no hay variación de velocidad, la velocidad al finalizar este segundo tramo es idéntica a la que tenía al iniciarlo:
\begin{equation}
v_2 = v_1 = 20 \text{ m/s}
\end{equation}

\subsubsection*{Análisis del Tercer Tramo (Aceleración Final)}

El vehículo retoma su aceleración a partir de la marca de los doscientos metros. Para este último tramo, su velocidad de entrada es la velocidad que mantuvo durante la fase previa.
\begin{itemize}
    \item Velocidad inicial del tramo $v_{03} = 20 \text{ m/s}$
    \item Aceleración constante $a_3 = 2 \text{ m/s}^2$
    \item Desplazamiento del tramo $\Delta x_3 = 100 \text{ m}$
\end{itemize}

Aplicamos nuevamente la ecuación cuadrática independiente del tiempo para encontrar la velocidad terminal de la prueba:
\begin{equation}
v_f^2 = v_{03}^2 + 2 \cdot a_3 \cdot \Delta x_3
\end{equation}

Sustituimos y resolvemos algebraicamente:
\begin{equation}
v_f^2 = (20)^2 + 2 \cdot (2) \cdot (100)
\end{equation}

\begin{equation}
v_f^2 = 400 + 400 = 800
\end{equation}

Extraemos la raíz cuadrada. Para preservar la exactitud matemática del modelo, expresamos el resultado en forma de radical simplificado y su aproximación decimal:
\begin{equation}
v_f = \sqrt{800} = \sqrt{400 \cdot 2} = 20\sqrt{2} \text{ m/s}
\end{equation}

\subsubsection*{Conclusión}

La rapidez del automóvil al finalizar la prueba es de $20\sqrt{2} \text{ m/s}$ (aproximadamente $28.28 \text{ m/s}$).

\newpage

\subsection*{Enunciado}

Durante una prueba de manejo, un automóvil se mueve a lo largo de una carretera recta y horizontal. El vehículo parte del reposo y acelera uniformemente a una tasa de $2 \text{ m/s}^2$ hasta recorrer $100 \text{ m}$. Luego, mantiene su rapidez constante durante otros $100 \text{ m}$. Finalmente, acelera nuevamente a una tasa de $2 \text{ m/s}^2$ hasta recorrer $100 \text{ m}$ adicionales. Determine cuál es el tiempo total de la prueba.

\vspace{0.5cm}

\begin{center}
\begin{tikzpicture}
\draw[thick, ->] (0,0) -- (12,0) node[right] {$x$};
\draw[dashed, gray] (3.5, 0) -- (3.5, -1);
\draw[dashed, gray] (7, 0) -- (7, -1);

\draw[<->, thick, black!70] (0, -0.6) -- (3.5, -0.6) node[midway, below] {$\Delta t_1 = ?$};
\draw[<->, thick, black!70] (3.5, -0.6) -- (7, -0.6) node[midway, below] {$\Delta t_2 = ?$};
\draw[<->, thick, black!70] (7, -0.6) -- (10.5, -0.6) node[midway, below] {$\Delta t_3 = ?$};

\filldraw[black] (0,0) circle (2pt) node[above=0.2cm] {$0 \text{ m}$};
\filldraw[black] (3.5,0) circle (2pt) node[above=0.2cm] {$100 \text{ m}$};
\filldraw[black] (7,0) circle (2pt) node[above=0.2cm] {$200 \text{ m}$};
\filldraw[black] (10.5,0) circle (2pt) node[above=0.2cm] {$300 \text{ m}$};
\end{tikzpicture}
\end{center}

\subsection*{Solución}

\subsubsection*{Conceptos Físicos Involucrados}

El tiempo total empleado para recorrer los trescientos metros es una magnitud escalar aditiva. Matemáticamente, equivale a la suma estricta de las duraciones temporales que el vehículo consume en completar cada tramo individual. Para determinar estos lapsos parciales, se utilizarán las definiciones fundamentales de velocidad y aceleración en función de las condiciones geométricas y cinemáticas calculadas previamente.

\subsubsection*{Duración del Primer Tramo}

El automóvil acelera desde cero hasta alcanzar veinte metros por segundo. Sabiendo que la aceleración es la tasa de cambio de la velocidad en el tiempo, podemos despejar el intervalo de tiempo:
\begin{equation}
a = \frac{\Delta v}{\Delta t} \implies \Delta t = \frac{v_f - v_0}{a}
\end{equation}

Evaluamos para los parámetros del primer intervalo:
\begin{equation}
\Delta t_1 = \frac{20 \text{ m/s} - 0 \text{ m/s}}{2 \text{ m/s}^2} = 10 \text{ s}
\end{equation}

\subsubsection*{Duración del Segundo Tramo}

Durante los cien metros intermedios, el vehículo describe un Movimiento Rectilíneo Uniforme con una velocidad inalterable de veinte metros por segundo. La relación entre desplazamiento, velocidad y tiempo es directa:
\begin{equation}
\Delta x = v \cdot \Delta t \implies \Delta t = \frac{\Delta x}{v}
\end{equation}

\begin{equation}
\Delta t_2 = \frac{100 \text{ m}}{20 \text{ m/s}} = 5 \text{ s}
\end{equation}

\subsubsection*{Duración del Tercer Tramo}

Para la fase de aceleración final, disponemos de la velocidad inicial del tramo ($20 \text{ m/s}$) y la velocidad final alcanzada ($20\sqrt{2} \text{ m/s}$). Recurrimos a la misma ecuación de definición de aceleración utilizada en la primera fase:
\begin{equation}
\Delta t_3 = \frac{v_f - v_{03}}{a_3}
\end{equation}

Sustituyendo los valores analíticos:
\begin{equation}
\Delta t_3 = \frac{20\sqrt{2} - 20}{2}
\end{equation}

Factorizando y reduciendo la expresión aritmética:
\begin{equation}
\Delta t_3 = 10\sqrt{2} - 10 \text{ s}
\end{equation}

Para efectos prácticos y suma escalar, calculamos su aproximación decimal utilizando $\sqrt{2} \approx 1.414$:
\begin{equation}
\Delta t_3 \approx 10(1.414) - 10 = 14.14 - 10 = 4.14 \text{ s}
\end{equation}

\subsubsection*{Cálculo del Tiempo Total}

La medición absoluta del reloj desde el inicio de la prueba hasta la marca de los trescientos metros se consolida sumando los tres intervalos parciales. Es indispensable ejecutar rigurosamente la aritmética para evitar errores conceptuales de adición.

\begin{equation}
t_{\text{total}} = \Delta t_1 + \Delta t_2 + \Delta t_3
\end{equation}

\begin{equation}
t_{\text{total}} = 10 \text{ s} + 5 \text{ s} + 4.14 \text{ s}
\end{equation}

\begin{equation}
t_{\text{total}} = 19.14 \text{ s}
\end{equation}

\subsubsection*{Análisis Gráfico del Movimiento}

La representación de la velocidad en función del tiempo es una de las herramientas analíticas más poderosas de la física conceptual. La pendiente de la recta indica la magnitud de la aceleración, y el área bajo la curva corresponde exactamente a la distancia recorrida. El siguiente código en Python genera dicho modelo visual.

\begin{lstlisting}[language=Python]
import matplotlib.pyplot as plt
import numpy as np

# Definicion de los intervalos de tiempo absoluto
t1 = np.linspace(0, 10, 100)
v1 = 2 * t1

t2 = np.linspace(10, 15, 100)
v2 = np.full_like(t2, 20)

t3 = np.linspace(15, 19.14, 100)
v3 = 20 + 2 * (t3 - 15)

plt.figure(figsize=(9, 6))

plt.plot(t1, v1, label='Tramo 1: Aceleracion', color='blue', linewidth=2.5)
plt.plot(t2, v2, label='Tramo 2: Velocidad Constante', color='green', linewidth=2.5)
plt.plot(t3, v3, label='Tramo 3: Aceleracion Final', color='red', linewidth=2.5)

plt.fill_between(t1, v1, color='blue', alpha=0.1)
plt.fill_between(t2, v2, color='green', alpha=0.1)
plt.fill_between(t3, v3, color='red', alpha=0.1)

plt.axvline(10, color='gray', linestyle='dashed', alpha=0.5)
plt.axvline(15, color='gray', linestyle='dashed', alpha=0.5)

plt.plot(19.14, 28.28, 'ko', markersize=6)
plt.text(14.5, 29, 'Fin (19.14 s, 28.28 m/s)', fontsize=10, fontweight='bold')

plt.title('Grafica de Velocidad vs Tiempo de la Prueba de Manejo')
plt.xlabel('Tiempo Total (s)')
plt.ylabel('Velocidad (m/s)')
plt.legend()
plt.grid(True)
plt.show()
\end{lstlisting}

\subsubsection*{Conclusión}

El tiempo total de la prueba es de 19.14 s.

\newpage

\subsection*{Enunciado}

A partir de la siguiente situación representada: se lanzan verticalmente hacia arriba tres bolas de masas diferentes. La Bola A tiene una masa de 1 kg y una rapidez inicial de 10 m/s. La Bola B tiene una masa de 1.5 kg y una rapidez inicial de 15 m/s. La Bola C tiene una masa de 0.8 kg y una rapidez inicial de 3 m/s. 

Antes de calcular, ¿cuál de las tres bolas consideras que tendrá mayor rapidez después de 1 segundo? Explica tu razonamiento a partir de la velocidad inicial y del efecto de la gravedad.

\vspace{0.5cm}

\begin{center}
\begin{tikzpicture}
\filldraw[fill=yellow!80!orange, draw=black, thick] (0,0) circle (0.4) node {\large A};
\draw[->, thick, green!60!black] (0,0.5) -- (0,2.5) node[left, midway] {$10 \text{ m/s}$};
\node[below] at (0,-0.6) {$1 \text{ kg}$};

\filldraw[fill=yellow!80!orange, draw=black, thick] (3,0) circle (0.4) node {\large B};
\draw[->, thick, green!60!black] (3,0.5) -- (3,3.5) node[left, midway] {$15 \text{ m/s}$};
\node[below] at (3,-0.6) {$1.5 \text{ kg}$};

\filldraw[fill=yellow!80!orange, draw=black, thick] (6,0) circle (0.4) node {\large C};
\draw[->, thick, green!60!black] (6,0.5) -- (6,1.5) node[left, midway] {$3 \text{ m/s}$};
\node[below] at (6,-0.6) {$0.8 \text{ kg}$};
\end{tikzpicture}
\end{center}

\subsection*{Solución}

\subsubsection*{Conceptos Físicos Involucrados}

De acuerdo con la física conceptual de Hewitt, en ausencia de resistencia del aire, todos los objetos en caída libre cerca de la superficie terrestre experimentan una aceleración gravitacional constante dirigida hacia el centro de la Tierra. Para facilitar el análisis conceptual, esta aceleración se aproxima a un valor de $10 \text{ m/s}^2$ hacia abajo. Esto significa que la gravedad resta exactamente $10 \text{ m/s}$ de velocidad cada segundo a cualquier objeto que se mueva hacia arriba.

Además, es crucial distinguir entre velocidad y rapidez. La velocidad es un vector que incluye dirección y sentido, mientras que la rapidez es el valor absoluto o magnitud de esa velocidad.

\subsubsection*{Análisis del Efecto de la Gravedad}

Analizamos intuitivamente lo que le ocurre a la velocidad de cada bola tras transcurrir un segundo de acción gravitatoria constante. 

La bola A parte con $10 \text{ m/s}$ hacia arriba. Al pasar un segundo, la gravedad le resta exactamente esos $10 \text{ m/s}$, por lo que su velocidad se reduce a cero. 

La bola B parte con $15 \text{ m/s}$ hacia arriba. Tras un segundo, la gravedad le quita $10 \text{ m/s}$, dejándole un remanente positivo de $5 \text{ m/s}$ en ascenso.

La bola C parte con una velocidad mucho menor, de $3 \text{ m/s}$ hacia arriba. La gravedad consumirá esta velocidad ascendente en apenas una fracción de segundo. Durante el resto del segundo evaluado, la gravedad comenzará a acelerar la bola hacia abajo. Dado que pasa más de medio segundo acelerando a favor de la gravedad, su velocidad final será negativa, pero su magnitud absoluta (la rapidez) crecerá significativamente en esa dirección. Al realizar la estimación mental, la bola C alcanza los $7 \text{ m/s}$ hacia abajo, superando en magnitud a la bola B.

\subsubsection*{Conclusión}

Antes de calcular de forma exacta, se deduce que la bola C tendrá la mayor rapidez después de 1 segundo, debido a que invierte su dirección rápidamente y acelera hacia abajo durante la mayor parte de ese segundo.

\newpage

\subsection*{Enunciado}

Se lanzan verticalmente hacia arriba tres bolas de masas diferentes. La Bola A tiene una masa de 1 kg y una rapidez inicial de 10 m/s. La Bola B tiene una masa de 1.5 kg y una rapidez inicial de 15 m/s. La Bola C tiene una masa de 0.8 kg y una rapidez inicial de 3 m/s. 

Clasifica, de mayor a menor, las rapideces de las bolas 1 segundo después de ser lanzadas. Justifica tu respuesta mediante el análisis del movimiento vertical.

\vspace{0.5cm}

\begin{center}
\begin{tikzpicture}
\filldraw[fill=yellow!80!orange, draw=black, thick] (0,0) circle (0.4) node {\large A};
\draw[->, thick, green!60!black] (0,0.5) -- (0,2.5) node[left, midway] {$10 \text{ m/s}$};
\node[below] at (0,-0.6) {$1 \text{ kg}$};

\filldraw[fill=yellow!80!orange, draw=black, thick] (3,0) circle (0.4) node {\large B};
\draw[->, thick, green!60!black] (3,0.5) -- (3,3.5) node[left, midway] {$15 \text{ m/s}$};
\node[below] at (3,-0.6) {$1.5 \text{ kg}$};

\filldraw[fill=yellow!80!orange, draw=black, thick] (6,0) circle (0.4) node {\large C};
\draw[->, thick, green!60!black] (6,0.5) -- (6,1.5) node[left, midway] {$3 \text{ m/s}$};
\node[below] at (6,-0.6) {$0.8 \text{ kg}$};
\end{tikzpicture}
\end{center}

\subsection*{Solución}

\subsubsection*{Modelo Cinemático}

Para justificar cuantitativamente la clasificación, utilizamos la ecuación cinemática que determina la velocidad final de un objeto sujeto a una aceleración constante. Establecemos un sistema de referencia unidimensional donde la dirección hacia arriba es positiva, lo que asigna un valor matemáticamente negativo a la aceleración de la gravedad.

\begin{equation}
v_f = v_0 + g \cdot t
\end{equation}

Utilizaremos la constante conceptual de $g = -10 \text{ m/s}^2$ y el tiempo fijo de $t = 1 \text{ s}$.

\subsubsection*{Evaluación de Velocidades}

Evaluamos la ecuación de movimiento para cada una de las masas de forma independiente.

Para la Bola A:
\begin{equation}
v_A = 10 + (-10) \cdot (1) = 0 \text{ m/s}
\end{equation}

Para la Bola B:
\begin{equation}
v_B = 15 + (-10) \cdot (1) = 5 \text{ m/s}
\end{equation}

Para la Bola C:
\begin{equation}
v_C = 3 + (-10) \cdot (1) = -7 \text{ m/s}
\end{equation}

\subsubsection*{Análisis de Rapidez y Representación Gráfica}

La rapidez es el módulo del vector velocidad y carece de dirección, por lo que tomamos el valor absoluto de cada resultado matemático: $|v_A| = 0 \text{ m/s}$, $|v_B| = 5 \text{ m/s}$, y $|v_C| = 7 \text{ m/s}$. El código de Python a continuación genera una gráfica que ilustra pedagógicamente cómo las rectas de velocidad cruzan el origen para cambiar de dirección y aumentar su módulo.

\begin{lstlisting}[language=Python]
import matplotlib.pyplot as plt
import numpy as np

t = np.linspace(0, 1.2, 100)
v_A = 10 - 10 * t
v_B = 15 - 10 * t
v_C = 3 - 10 * t

plt.figure(figsize=(8, 5))
plt.plot(t, v_A, label='Bola A', color='orange', linewidth=2)
plt.plot(t, v_B, label='Bola B', color='blue', linewidth=2)
plt.plot(t, v_C, label='Bola C', color='green', linewidth=2)

plt.axhline(0, color='black', linestyle='dashed', linewidth=1)
plt.axvline(1, color='red', linestyle='dotted', linewidth=1.5, label='t = 1 s')

plt.plot(1, 0, 'ko')
plt.plot(1, 5, 'ko')
plt.plot(1, -7, 'ko')

plt.title('Velocidad vs Tiempo de las Tres Bolas')
plt.xlabel('Tiempo (s)')
plt.ylabel('Velocidad (m/s)')
plt.legend()
plt.grid(True)
plt.show()
\end{lstlisting}

\subsubsection*{Conclusión}

Clasificadas de mayor a menor, las rapideces son: Bola C ($7 \text{ m/s}$) > Bola B ($5 \text{ m/s}$) > Bola A ($0 \text{ m/s}$).

\newpage

\subsection*{Enunciado}

Se lanzan verticalmente hacia arriba tres bolas de masas diferentes. La Bola A tiene una masa de 1 kg y una rapidez inicial de 10 m/s. La Bola B tiene una masa de 1.5 kg y una rapidez inicial de 15 m/s. La Bola C tiene una masa de 0.8 kg y una rapidez inicial de 3 m/s. 

Compara las aceleraciones de las tres bolas 1 segundo después del lanzamiento. ¿Son iguales o diferentes? Explica si la masa influye o no en la aceleración durante este tipo de movimiento.

\vspace{0.5cm}

\begin{center}
\begin{tikzpicture}
\filldraw[fill=yellow!80!orange, draw=black, thick] (0,0) circle (0.4) node {\large A};
\draw[->, thick, green!60!black] (0,0.5) -- (0,2.5) node[left, midway] {$10 \text{ m/s}$};
\node[below] at (0,-0.6) {$1 \text{ kg}$};

\filldraw[fill=yellow!80!orange, draw=black, thick] (3,0) circle (0.4) node {\large B};
\draw[->, thick, green!60!black] (3,0.5) -- (3,3.5) node[left, midway] {$15 \text{ m/s}$};
\node[below] at (3,-0.6) {$1.5 \text{ kg}$};

\filldraw[fill=yellow!80!orange, draw=black, thick] (6,0) circle (0.4) node {\large C};
\draw[->, thick, green!60!black] (6,0.5) -- (6,1.5) node[left, midway] {$3 \text{ m/s}$};
\node[below] at (6,-0.6) {$0.8 \text{ kg}$};
\end{tikzpicture}
\end{center}

\subsection*{Solución}

\subsubsection*{Segunda Ley de Newton y Gravitación}

Para comprender la dinámica de los objetos en caída libre, es necesario recurrir a la Segunda Ley de Newton, la cual establece que la aceleración de un cuerpo es directamente proporcional a la fuerza neta que actúa sobre él e inversamente proporcional a su masa inercial.

\begin{equation}
a = \frac{F_{\text{neta}}}{m}
\end{equation}

Cuando un objeto ha sido lanzado por el aire y se desprecia la resistencia aerodinámica, la única fuerza externa significativa que actúa sobre él es la fuerza de atracción gravitacional de la Tierra, comúnmente llamada peso ($P$). El peso se define como el producto de la masa del objeto por la constante del campo gravitatorio terrestre.

\begin{equation}
F_{\text{neta}} = P = m \cdot g
\end{equation}

\subsubsection*{Independencia de la Masa}

Sustituimos la fuerza del peso en la ecuación fundamental de la dinámica para determinar la aceleración real que experimentan las bolas independientemente de sus propiedades físicas iniciales.

\begin{equation}
a = \frac{m \cdot g}{m}
\end{equation}

Al realizar el álgebra en esta relación, la variable de la masa se cancela tanto en el numerador como en el denominador.

\begin{equation}
a = g
\end{equation}

Este postulado de Galileo demuestra matemáticamente que la inercia del objeto (su resistencia a cambiar de movimiento, dictada por su masa) se equilibra de manera exacta con el incremento proporcional en la fuerza de atracción gravitatoria. Un objeto de mayor masa pesa más, pero también es más difícil de acelerar. Ambos efectos se anulan mutuamente en el vacío espacial.

\subsubsection*{Conclusión}

Las aceleraciones de las tres bolas son exactamente iguales (aproximadamente $10 \text{ m/s}^2$ hacia abajo). La masa no influye en absoluto en la aceleración de los cuerpos bajo la única influencia de la gravedad.

\newpage

\subsection*{Enunciado}

Se lanzan verticalmente hacia arriba tres bolas de masas diferentes. La Bola A tiene una masa de 1 kg y una rapidez inicial de 10 m/s. La Bola B tiene una masa de 1.5 kg y una rapidez inicial de 15 m/s. La Bola C tiene una masa de 0.8 kg y una rapidez inicial de 3 m/s. 

Una de las bolas puede encontrarse, en ese instante (1 segundo después del lanzamiento), detenida momentáneamente, otra aun ascendiendo y otra ya descendiendo. Identifica cuál corresponde a cada caso y explica por qué.

\vspace{0.5cm}

\begin{center}
\begin{tikzpicture}
\filldraw[fill=yellow!80!orange, draw=black, thick] (0,0) circle (0.4) node {\large A};
\draw[->, thick, green!60!black] (0,0.5) -- (0,2.5) node[left, midway] {$10 \text{ m/s}$};
\node[below] at (0,-0.6) {$1 \text{ kg}$};

\filldraw[fill=yellow!80!orange, draw=black, thick] (3,0) circle (0.4) node {\large B};
\draw[->, thick, green!60!black] (3,0.5) -- (3,3.5) node[left, midway] {$15 \text{ m/s}$};
\node[below] at (3,-0.6) {$1.5 \text{ kg}$};

\filldraw[fill=yellow!80!orange, draw=black, thick] (6,0) circle (0.4) node {\large C};
\draw[->, thick, green!60!black] (6,0.5) -- (6,1.5) node[left, midway] {$3 \text{ m/s}$};
\node[below] at (6,-0.6) {$0.8 \text{ kg}$};
\end{tikzpicture}
\end{center}

\subsection*{Solución}

\subsubsection*{Significado Físico del Signo de la Velocidad}

En el modelado unidimensional de la cinemática, la dirección del movimiento está codificada explícitamente en el signo del vector velocidad. Habiendo definido nuestro marco de referencia con el eje positivo apuntando hacia el cielo en oposición a la gravedad terrestre, podemos establecer reglas de diagnóstico directas:
\begin{itemize}
    \item Si $v > 0$, el objeto posee desplazamiento positivo y se encuentra ascendiendo.
    \item Si $v = 0$, el objeto ha consumido toda su energía cinética ascendente, se encuentra en el vértice de su trayectoria parabólica y está detenido momentáneamente.
    \item Si $v < 0$, el objeto posee desplazamiento negativo respecto a su punto más alto y se encuentra descendiendo en caída libre.
\end{itemize}

\subsubsection*{Clasificación de los Estados de Movimiento}

Recuperamos los valores exactos de los vectores de velocidad calculados para el instante temporal $t = 1 \text{ s}$ utilizando la constante de aceleración gravitatoria:

\begin{equation}
v_A = 0 \text{ m/s}
\end{equation}
La Bola A es la entidad cuya velocidad evaluada es exactamente cero. Esto ocurre porque la desaceleración disipó su velocidad inicial de forma exacta en el marco de tiempo evaluado. Se encuentra en su punto de altura máxima.

\begin{equation}
v_B = +5 \text{ m/s}
\end{equation}
La Bola B mantiene un vector velocidad con signo positivo. Esto indica físicamente que la bola todavía avanza en la dirección original del lanzamiento porque la gravedad no ha tenido el tiempo suficiente para revertir por completo su gran impulso inicial de quince metros por segundo.

\begin{equation}
v_C = -7 \text{ m/s}
\end{equation}
La Bola C exhibe un vector velocidad con signo negativo. Su escasa velocidad inicial de tres metros por segundo fue neutralizada en apenas tres décimas de segundo. Durante los siete décimos de segundo restantes, ha estado acelerando ininterrumpidamente hacia el centro de la Tierra.

\subsubsection*{Conclusión}

\begin{itemize}
    \item \textbf{Bola A:} Se encuentra detenida momentáneamente porque su velocidad ha sido reducida a $0 \text{ m/s}$.
    \item \textbf{Bola B:} Aún se encuentra ascendiendo porque conserva una velocidad positiva de $5 \text{ m/s}$.
    \item \textbf{Bola C:} Ya se encuentra descendiendo porque posee una velocidad negativa de $-7 \text{ m/s}$.
\end{itemize}

\newpage

\subsection*{Enunciado}

Desde un mismo punto, se deja caer simultáneamente dos cuerpos. El primero se suelta desde el reposo, mientras que el segundo es lanzado hacia abajo con una rapidez inicial de $3 \text{ m/s}$. Analiza la situación y determina en qué momento la distancia entre ambos cuerpos es de $9 \text{ m}$. Antes de realizar cálculos, ¿cuál de los dos cuerpos consideras que irá más abajo conforme transcurre el tiempo? Explica tu respuesta a partir de sus condiciones iniciales.

\vspace{0.5cm}

\begin{center}
\begin{tikzpicture}
\draw[dashed, gray] (-2,0) -- (4,0) node[right] {Nivel inicial};

\filldraw[fill=blue!50, draw=black, thick] (0,0) circle (0.3) node[above=0.4cm] {Cuerpo 1};
\node[right] at (0.3, 0) {$v_{01} = 0 \text{ m/s}$};

\filldraw[fill=red!50, draw=black, thick] (2,0) circle (0.3) node[above=0.4cm] {Cuerpo 2};
\draw[->, thick, red!80!black] (2,-0.4) -- (2,-1.5) node[right, midway] {$v_{02} = 3 \text{ m/s}$};

\draw[->, thick, black] (-1.5,-0.5) -- (-1.5,-2) node[left, midway] {$\vec{g}$};
\end{tikzpicture}
\end{center}

\subsection*{Solución}

\subsubsection*{Conceptos Físicos Involucrados}

En el estudio de la cinemática, el estado de movimiento de un cuerpo depende críticamente de sus condiciones iniciales. Según los principios expuestos en la física conceptual, un objeto que se suelta desde el reposo inicia su recorrido con una velocidad de cero y debe acumular velocidad paulatinamente gracias a la aceleración de la gravedad. Por otro lado, un objeto que es lanzado hacia abajo ya posee un "impulso" inicial; inicia su caída con una velocidad preexistente a la cual la gravedad sumará aún más velocidad.

\subsubsection*{Análisis Cualitativo del Movimiento}

Evaluemos mentalmente el escenario. En el instante inicial ($t = 0$), ambos cuerpos se encuentran en la misma posición, pero el Cuerpo 2 ya tiene una rapidez de tres metros por segundo hacia abajo. 

Al transcurrir el primer segundo, la gravedad (aproximada conceptualmente a $10 \text{ m/s}^2$) habrá incrementado la rapidez de ambos cuerpos en la misma proporción. El Cuerpo 1 pasará de $0$ a $10 \text{ m/s}$, mientras que el Cuerpo 2 pasará de $3$ a $13 \text{ m/s}$. 

Como podemos observar, debido a que la gravedad afecta a ambos por igual, la diferencia de velocidades se mantiene constante en todo momento. El Cuerpo 2 siempre viajará tres metros por segundo más rápido que el Cuerpo 1. Al moverse a una mayor velocidad en cada instante de la caída, necesariamente cubrirá una mayor distancia temporal y espacial.

\subsubsection*{Conclusión}

El segundo cuerpo irá más abajo conforme transcurre el tiempo porque parte con una velocidad inicial mayor, lo que garantiza que en cualquier instante de la caída su velocidad siempre sea superior a la del primer cuerpo, permitiéndole recorrer más distancia.

\newpage

\subsection*{Enunciado}

Desde un mismo punto, se deja caer simultáneamente dos cuerpos. El primero se suelta desde el reposo, mientras que el segundo es lanzado hacia abajo con una rapidez inicial de $3 \text{ m/s}$. Analiza la situación y determina en qué momento la distancia entre ambos cuerpos es de $9 \text{ m}$. Ambos cuerpos están sometidos a la acción de la gravedad. ¿Tienen la misma aceleración o aceleraciones diferentes? Justifica tu respuesta y explica si esta igualdad influye en la separación entre ellos.

\vspace{0.5cm}

\begin{center}
\begin{tikzpicture}
\draw[dashed, gray] (-2,0) -- (4,0) node[right] {Nivel inicial};

\filldraw[fill=blue!50, draw=black, thick] (0,0) circle (0.3) node[above=0.4cm] {Cuerpo 1};
\draw[->, thick, blue] (0,-0.4) -- (0,-1.0) node[left] {$\vec{v}_1$};

\filldraw[fill=red!50, draw=black, thick] (2,0) circle (0.3) node[above=0.4cm] {Cuerpo 2};
\draw[->, thick, red] (2,-0.4) -- (2,-1.8) node[right] {$\vec{v}_2$};

\draw[->, thick, black] (-1.5,-0.5) -- (-1.5,-2) node[left, midway] {$\vec{g}$};
\end{tikzpicture}
\end{center}

\subsection*{Solución}

\subsubsection*{El Principio de Equivalencia y la Caída Libre}

Uno de los postulados fundamentales establecidos por Galileo Galilei y posteriormente formalizado por Isaac Newton es que, en ausencia de resistencia del aire (o considerando que sus efectos son despreciables), todos los objetos en caída libre cerca de la superficie de la Tierra experimentan exactamente la misma aceleración, sin importar su masa, su tamaño o su velocidad inicial.

La aceleración de la gravedad es una propiedad del campo gravitatorio del planeta, no de los objetos que caen en él. 

\subsubsection*{Análisis de la Aceleración de los Cuerpos}

Dado que tanto el cuerpo que se deja caer como el que es lanzado hacia abajo están sujetos únicamente a la fuerza de atracción gravitacional terrestre una vez que abandonan la mano, ambos se encuentran en un estado de caída libre perfecta. Por lo tanto, poseen la misma aceleración, denotada por la constante $g$ (aproximadamente $9.8 \text{ m/s}^2$ apuntando hacia el centro de la Tierra).

\subsubsection*{Efecto sobre la Separación Relativa}

El concepto de velocidad y posición relativa es clave aquí. Si calculamos la aceleración del Cuerpo 2 respecto al Cuerpo 1, debemos restar sus aceleraciones individuales:
\begin{equation}
a_{\text{relativa}} = a_2 - a_1 = g - g = 0 \text{ m/s}^2
\end{equation}

El hecho de que la aceleración relativa sea nula tiene una consecuencia matemática profunda: significa que, desde la perspectiva de uno de los cuerpos, el otro no está acelerando, sino alejándose a una velocidad constante. La igualdad en la aceleración hace que el término cuadrático del desplazamiento (que provoca la caída acelerada hacia el suelo) sea idéntico para ambos y se anule al comparar sus posiciones. 

La separación crece uniformemente, dictada única y exclusivamente por la diferencia en sus velocidades iniciales.

\subsubsection*{Conclusión}

Ambos cuerpos tienen la misma aceleración porque se encuentran en caída libre bajo la influencia del mismo campo gravitatorio. Esta igualdad hace que la aceleración no influya en el aumento de la separación entre ellos, permitiendo que se alejen el uno del otro a un ritmo constante equivalente a la diferencia de sus velocidades iniciales.

\newpage

\subsection*{Enunciado}

Desde un mismo punto, se deja caer simultáneamente dos cuerpos. El primero se suelta desde el reposo, mientras que el segundo es lanzado hacia abajo con una rapidez inicial de $3 \text{ m/s}$. Determina en qué instante la distancia entre ambos cuerpos es de $9 \text{ m}$. Explica el procedimiento utilizado e indica por qué en este caso resulta útil comparar los desplazamientos de ambos cuerpos.

\vspace{0.5cm}

\begin{center}
\begin{tikzpicture}
\draw[dashed, gray] (-2,0) -- (4,0) node[right] {Nivel inicial $y=0$};

\filldraw[fill=blue!50, draw=black, thick] (0,-2) circle (0.3) node[left=0.2cm] {$y_1$};
\filldraw[fill=red!50, draw=black, thick] (2,-4) circle (0.3) node[right=0.2cm] {$y_2$};

\draw[<->, thick, green!60!black] (0.5,-2) -- (0.5,-4) node[midway, right] {$\Delta y = 9 \text{ m}$};
\draw[dashed, gray] (0,-2) -- (0.5,-2);
\draw[dashed, gray] (2,-4) -- (0.5,-4);

\end{tikzpicture}
\end{center}

\subsection*{Solución}

\subsubsection*{Definición del Sistema de Referencia y Ecuaciones}

Para resolver el problema analíticamente, establecemos un sistema de coordenadas vertical unidimensional. Ubicaremos el origen ($y = 0$) en el punto exacto de lanzamiento y definiremos la dirección hacia abajo como positiva. Esta convención matemática simplifica los signos de las ecuaciones, ya que tanto las velocidades como las aceleraciones apuntan en la misma dirección positiva.

La posición de un cuerpo con aceleración constante está regida por la ecuación:
\begin{equation}
y(t) = y_0 + v_0 t + \frac{1}{2} a t^2
\end{equation}

Para el Cuerpo 1 (que parte del reposo):
\begin{itemize}
    \item Posición inicial $y_{01} = 0$
    \item Velocidad inicial $v_{01} = 0$
    \item Aceleración $a_1 = g$
\end{itemize}
\begin{equation}
y_1(t) = \frac{1}{2} g t^2
\end{equation}

Para el Cuerpo 2 (lanzado hacia abajo):
\begin{itemize}
    \item Posición inicial $y_{02} = 0$
    \item Velocidad inicial $v_{02} = 3 \text{ m/s}$
    \item Aceleración $a_2 = g$
\end{itemize}
\begin{equation}
y_2(t) = 3t + \frac{1}{2} g t^2
\end{equation}

\subsubsection*{Condición de Separación y Desarrollo Matemático}

Se solicita encontrar el instante en el que la distancia vertical entre ellos es de nueve metros. Geométricamente, esta distancia se define como la diferencia entre sus coordenadas espaciales:
\begin{equation}
\Delta y = y_2(t) - y_1(t)
\end{equation}

Sustituimos la separación exigida por el problema y las ecuaciones de movimiento que hemos modelado:
\begin{equation}
9 = \left( 3t + \frac{1}{2} g t^2 \right) - \left( \frac{1}{2} g t^2 \right)
\end{equation}

Como se abordó en la justificación cualitativa, la aceleración de la gravedad es idéntica para ambos cuerpos. Al restar sus posiciones espaciales, los términos matemáticos que contienen la aceleración gravitatoria se anulan mutuamente:
\begin{equation}
9 = 3t + \frac{1}{2} g t^2 - \frac{1}{2} g t^2
\end{equation}

\begin{equation}
9 = 3t
\end{equation}

Despejamos la incógnita del tiempo dividiendo por la velocidad relativa:
\begin{equation}
t = \frac{9}{3}
\end{equation}

\begin{equation}
t = 3 \text{ s}
\end{equation}

\subsubsection*{Representación Gráfica del Movimiento Relativo}

Visualizar las gráficas de desplazamiento en función del tiempo para ambos cuerpos ilustra cómo el término de aceleración genera una parábola, pero la distancia entre ellas se incrementa de forma lineal y constante.

\begin{lstlisting}[language=Python]
import matplotlib.pyplot as plt
import numpy as np

t_array = np.linspace(0, 4, 100)
g = 9.81
pos_cuerpo1 = 0.5 * g * t_array**2
pos_cuerpo2 = 3 * t_array + 0.5 * g * t_array**2

plt.figure(figsize=(8, 5))
plt.plot(t_array, pos_cuerpo1, label='Cuerpo 1 (v0 = 0)', color='blue', linewidth=2)
plt.plot(t_array, pos_cuerpo2, label='Cuerpo 2 (v0 = 3 m/s)', color='red', linewidth=2)

plt.axvline(3, color='black', linestyle='dashed', alpha=0.5)

plt.annotate('', xy=(3, 0.5*g*(3**2)), xytext=(3, 3*3 + 0.5*g*(3**2)),
             arrowprops=dict(arrowstyle='<->', color='green', lw=2))

plt.text(3.1, (0.5*g*(3**2) + 3*3 + 0.5*g*(3**2))/2, 'Separacion: 9 m', color='green', rotation=-90, verticalalignment='center')

plt.title('Desplazamiento vs Tiempo de los dos Cuerpos')
plt.xlabel('Tiempo (s)')
plt.ylabel('Posicion desde el origen hacia abajo (m)')
plt.legend()
plt.gca().invert_yaxis()
plt.grid(True)
plt.show()
\end{lstlisting}

\subsubsection*{Conclusión}

El instante en que la distancia entre ambos cuerpos es de $9 \text{ m}$ es a los $3 \text{ s}$. El procedimiento consistió en restar las funciones de posición. Resulta extremadamente útil comparar los desplazamientos de ambos cuerpos porque los términos de aceleración se anulan, transformando una ecuación cinemática cuadrática compleja en un simple problema de velocidad constante.

\newpage

\subsection*{Enunciado}

Desde la ventana de un edificio se lanzan dos piedras, A y B, con la misma rapidez inicial. La piedra A se lanza verticalmente hacia arriba, mientras que la piedra B se lanza verticalmente hacia abajo. Despreciando la resistencia del aire, antes de realizar cualquier cálculo, ¿cuál crees que llegará al suelo con mayor rapidez: la piedra A, la piedra B, o ambas con la misma rapidez? Explica tu razonamiento.

\vspace{0.5cm}

\begin{center}
\begin{tikzpicture}
\draw[thick, fill=gray!20] (0,0) rectangle (2,4);
\draw[thick, fill=cyan!10] (1,3) rectangle (2,4);
\draw[thick] (0,0) -- (5,0) node[right] {Suelo};

\filldraw[fill=yellow!80!orange, draw=black, thick] (2.5, 3.5) circle (0.2) node[right=0.2cm] {A};
\draw[->, thick, green!60!black] (2.5, 3.8) -- (2.5, 5) node[above] {$v_0$};

\filldraw[fill=yellow!80!orange, draw=black, thick] (3.5, 3.5) circle (0.2) node[right=0.2cm] {B};
\draw[->, thick, green!60!black] (3.5, 3.2) -- (3.5, 2) node[below] {$v_0$};

\draw[dashed, gray] (2, 3.5) -- (4, 3.5);
\end{tikzpicture}
\end{center}

\subsection*{Solución}

\subsubsection*{Conceptos Físicos Involucrados}

La física conceptual enfatiza la simetría inherente en los movimientos de caída libre. Cuando un objeto es lanzado verticalmente hacia arriba en un campo gravitatorio uniforme y sin resistencia del aire, este pierde velocidad a un ritmo constante hasta detenerse momentáneamente en su punto más alto. A partir de ese instante, comienza a caer, ganando velocidad exactamente al mismo ritmo. 

Esta simetría implica que, en cualquier punto de su trayectoria de descenso, el objeto tendrá exactamente la misma rapidez que tenía cuando pasó por esa misma altura durante su ascenso, pero con la dirección de la velocidad invertida.

\subsubsection*{Análisis del Movimiento de la Piedra A}

La piedra A es lanzada hacia arriba con una rapidez inicial determinada. Alcanzará una altura máxima y luego comenzará a caer libremente. Al analizar su trayecto descendente, habrá un instante exacto en el que la piedra A pasará nuevamente por el nivel de la ventana desde donde fue lanzada. 

Debido a la simetría del movimiento descrita anteriormente, en ese preciso instante en que cruza el nivel de la ventana hacia abajo, su rapidez será idéntica a su rapidez inicial de lanzamiento.

\subsubsection*{Comparación de Estados Cinemáticos}

Consideremos el estado de la piedra A justo en el momento en que cruza el nivel de la ventana en su caída. Se encuentra a la misma altura inicial del problema, y se mueve verticalmente hacia abajo con la rapidez inicial original. 

Este estado es físicamente indistinguible del estado inicial de la piedra B, la cual fue lanzada desde esa misma altura directamente hacia abajo con esa misma rapidez. Puesto que a partir de ese nivel ambas piedras experimentan las mismas condiciones iniciales efectivas (misma altura y misma velocidad inicial hacia abajo), la gravedad las acelerará de la misma manera durante el resto del trayecto hasta el suelo.

\subsubsection*{Conclusión}

Ambas piedras llegarán al suelo con la misma rapidez.

\newpage

\subsection*{Enunciado}

Desde la ventana de un edificio se lanzan dos piedras, A y B. La piedra A se lanza verticalmente hacia arriba, mientras que la piedra B se lanza verticalmente hacia abajo. Aunque la piedra A primero sube y luego baja, mientras que la piedra B baja directamente, ¿las dos parten desde la misma altura y con la misma rapidez inicial? Explica por qué esta información es importante para comparar sus rapideces al llegar al suelo.

\vspace{0.5cm}

\begin{center}
\begin{tikzpicture}
\draw[thick, fill=gray!20] (0,0) rectangle (2,4);
\draw[thick] (0,0) -- (6,0) node[right] {Suelo ($h=0$)};

\draw[dashed, thick, blue] (2, 3.5) -- (5, 3.5);
\draw[<->, thick, blue] (4.5, 0) -- (4.5, 3.5) node[midway, fill=white] {Altura inicial ($h$)};

\filldraw[fill=yellow!80!orange, draw=black, thick] (2.5, 3.5) circle (0.2) node[right=0.2cm] {A};
\draw[->, thick, green!60!black] (2.5, 3.8) -- (2.5, 5) node[above] {$|v_{0A}| = v_0$};

\filldraw[fill=yellow!80!orange, draw=black, thick] (3.5, 3.5) circle (0.2) node[right=0.2cm] {B};
\draw[->, thick, green!60!black] (3.5, 3.2) -- (3.5, 2) node[below] {$|v_{0B}| = v_0$};

\end{tikzpicture}
\end{center}

\subsection*{Solución}

\subsubsection*{Análisis de las Condiciones Físicas Iniciales}

El enunciado establece de manera explícita que ambas piedras son lanzadas desde la misma ventana del edificio. En un sistema de referencia físico, esto garantiza que la coordenada de posición vertical en el instante inicial es idéntica para ambos cuerpos. Por lo tanto, sí parten desde la misma altura.

Respecto a la velocidad, es fundamental distinguir conceptualmente entre el vector velocidad y su magnitud escalar, denominada rapidez. La piedra A tiene una velocidad inicial positiva (hacia arriba) y la piedra B tiene una velocidad inicial negativa (hacia abajo). Sin embargo, la magnitud absoluta de ambos vectores es igual. Por lo tanto, sí parten con la misma rapidez inicial.

\subsubsection*{Importancia en la Dinámica del Movimiento}

Esta información geométrica y escalar es de vital importancia porque dictamina la cantidad total de energía que el sistema transfiere a cada piedra en el instante cero. 

La energía potencial gravitatoria depende exclusivamente de la altura respecto a un nivel de referencia. Al compartir la misma altura, ambas piedras poseen exactamente la misma energía potencial inicial.

La energía cinética, que representa la energía del movimiento, se define matemáticamente mediante el cuadrado de la rapidez. Al elevar al cuadrado, el signo direccional del vector velocidad pierde su relevancia, y la energía depende únicamente de la magnitud escalar. Al poseer la misma rapidez inicial, ambas piedras reciben exactamente la misma cantidad de energía cinética.

\subsubsection*{Conclusión}

Sí, las dos parten desde la misma altura y con la misma rapidez inicial. Esta información es fundamental porque demuestra que ambas piedras inician su movimiento con una cantidad idéntica de energía mecánica total, lo cual determinará de forma inequívoca su estado energético al impactar el suelo.

\newpage

\subsection*{Enunciado}

Desde la ventana de un edificio se lanzan dos piedras, A y B, con la misma rapidez inicial. La piedra A se lanza verticalmente hacia arriba, mientras que la piedra B se lanza verticalmente hacia abajo. Despreciando la resistencia del aire, analiza la situación usando el principio de conservación de la energía mecánica y determina la relación entre las rapideces finales de ambas piedras al llegar al suelo.

\vspace{0.5cm}

\begin{center}
\begin{tikzpicture}
\draw[thick] (0,0) -- (8,0) node[right] {Suelo ($U=0$)};
\draw[dashed, thick, gray] (1,4) -- (7,4) node[right] {Nivel inicial ($U=mgh$)};

\filldraw[fill=yellow!80!orange, draw=black, thick] (2,4) circle (0.2) node[left=0.2cm] {A};
\draw[->, thick, blue] (2,4.3) -- (2,5.5) node[above] {$v_0$};
\draw[dashed, ->, thick, gray] (2, 5.5) .. controls (2.5, 6) .. (3, 5.5) -- (3, 0.3);
\filldraw[fill=yellow!80!orange, draw=black, thick, opacity=0.5] (3,0.2) circle (0.2);
\draw[->, thick, red] (3,0) -- (3,-1.2) node[below] {$v_{fA}$};

\filldraw[fill=yellow!80!orange, draw=black, thick] (5,4) circle (0.2) node[right=0.2cm] {B};
\draw[->, thick, blue] (5,3.7) -- (5,2.5) node[below] {$v_0$};
\draw[dashed, ->, thick, gray] (5, 2.5) -- (5, 0.3);
\filldraw[fill=yellow!80!orange, draw=black, thick, opacity=0.5] (5,0.2) circle (0.2);
\draw[->, thick, red] (5,0) -- (5,-1.2) node[below] {$v_{fB}$};

\end{tikzpicture}
\end{center}

\subsection*{Solución}

\subsubsection*{Conceptos Físicos Involucrados}

El principio de conservación de la energía mecánica es una herramienta pedagógica sumamente poderosa que permite predecir los estados finales de un sistema sin necesidad de calcular las variables intermedias, como el tiempo o las posiciones parciales. 

Este principio establece que, si sobre un cuerpo actúan únicamente fuerzas conservativas (como la gravedad terrestre), la suma de su energía cinética y su energía potencial gravitatoria permanece constante a lo largo de toda su trayectoria. Al omitir la resistencia del aire, este escenario cumple perfectamente con dicha condición.

\subsubsection*{Planteamiento de la Energía Inicial}

Definimos el nivel del suelo como nuestro punto de referencia topográfico, asignándole una energía potencial nula ($h = 0$). La ventana del edificio se encuentra a una altura positiva $h$. 

La energía mecánica inicial ($E_i$) para cualquier piedra es la suma de su energía cinética ($K_i$) y su energía potencial ($U_i$):

\begin{equation}
E_i = K_i + U_i
\end{equation}

Sustituyendo las definiciones matemáticas de la mecánica clásica:

\begin{equation}
E_i = \frac{1}{2} m v_0^2 + m g h
\end{equation}

Como se demostró en el análisis previo, el término de la energía cinética depende del cuadrado de la rapidez inicial ($v_0^2$). Por lo tanto, el hecho de que la piedra A suba y la piedra B baje no altera en absoluto este valor. Ambas piedras poseen exactamente la misma energía mecánica inicial.

\subsubsection*{Planteamiento de la Energía Final}

En el instante inmediatamente anterior a impactar el suelo, la altura de las piedras respecto a nuestro nivel de referencia es cero. Esto anula por completo la energía potencial gravitatoria del sistema. Toda la energía mecánica inicial se ha transformado íntegramente en energía cinética de traslación.

\begin{equation}
E_f = K_f + 0
\end{equation}

\begin{equation}
E_f = \frac{1}{2} m v_f^2
\end{equation}

\subsubsection*{Desarrollo Matemático}

Aplicamos el principio de conservación igualando la energía en el estado inicial con la energía en el estado final:

\begin{equation}
E_i = E_f
\end{equation}

\begin{equation}
\frac{1}{2} m v_0^2 + m g h = \frac{1}{2} m v_f^2
\end{equation}

Dividimos ambos lados de la ecuación por la masa inercial ($m$), demostrando que la trayectoria y las rapideces resultantes son independientes de las proporciones del objeto:

\begin{equation}
\frac{1}{2} v_0^2 + g h = \frac{1}{2} v_f^2
\end{equation}

Despejamos la rapidez final multiplicando toda la ecuación por dos y extrayendo la raíz cuadrada:

\begin{equation}
v_0^2 + 2 g h = v_f^2
\end{equation}

\begin{equation}
v_f = \sqrt{v_0^2 + 2 g h}
\end{equation}

\subsubsection*{Representación Analítica de la Energía}

La elegancia del enfoque energético radica en que las trayectorias divergentes terminan confluyendo en la misma ecuación final. El siguiente código en Python modela gráficamente las energías de ambos objetos, validando que los depósitos de energía final son matemáticamente indistinguibles.

\begin{lstlisting}[language=Python]
import matplotlib.pyplot as plt
import numpy as np

# Definicion de parametros arbitrarios para la demostracion
m = 1.0        # Masa (kg)
v0 = 10.0      # Rapidez inicial (m/s)
h = 20.0       # Altura inicial (m)
g = 9.81       # Gravedad (m/s^2)

# Calculo de energias Iniciales
K_i = 0.5 * m * v0**2
U_i = m * g * h
E_total = K_i + U_i

# Calculo de energias Finales (Suelo)
U_f = 0
K_f = E_total

categorias = ['Estado Inicial\n(Ventana)', 'Estado Final\n(Suelo)']
energia_potencial = [U_i, U_f]
energia_cinetica = [K_i, K_f]

plt.figure(figsize=(8, 5))
bar_width = 0.5

# Grafico de barras apiladas
plt.bar(categorias, energia_potencial, color='lightblue', label='Energia Potencial (U)', width=bar_width)
plt.bar(categorias, energia_cinetica, bottom=energia_potencial, color='salmon', label='Energia Cinetica (K)', width=bar_width)

plt.axhline(E_total, color='black', linestyle='dashed', linewidth=1.5, label='Energia Mecanica Total')

plt.title('Conservacion de la Energia para ambas Piedras (A y B)')
plt.ylabel('Energia (Joules)')
plt.legend()
plt.grid(axis='y', linestyle='--', alpha=0.7)
plt.show()
\end{lstlisting}

\subsubsection*{Conclusión}

Al resolver la ecuación de conservación, observamos que la rapidez final ($v_f$) depende exclusivamente de la magnitud de la velocidad inicial ($v_0$), la gravedad ($g$) y la altura ($h$). Dado que ambas piedras comparten estos tres parámetros exactos, la relación entre sus rapideces finales es una igualdad absoluta ($v_{fA} = v_{fB}$).

\newpage

\subsection*{Enunciado}

Un cuerpo cae libremente desde el reposo durante $4 \text{ s}$. A partir de esta situación, analiza la relación entre la distancia recorrida en el primer segundo y la distancia recorrida en el último segundo del movimiento. Antes de realizar cálculos, ¿crees que el cuerpo recorre la misma distancia en cada segundo de caída o distancias diferentes? Explica tu razonamiento a partir del efecto de la gravedad sobre la rapidez del cuerpo.

\vspace{0.5cm}

\begin{center}
\begin{tikzpicture}
\draw[thick, fill=gray!30] (-2,0) rectangle (2,0.5);
\filldraw[fill=blue!50, draw=black, thick] (0,-0.3) circle (0.3) node[right=0.4cm] {$v_0 = 0 \text{ m/s}$};
\draw[->, thick, black] (-1.5,-0.5) -- (-1.5,-2) node[left, midway] {$\vec{g}$};
\draw[dashed, thick, gray] (0,-0.6) -- (0,-4);
\node[right] at (0.2, -2) {Caída libre ($t = 4 \text{ s}$)};
\end{tikzpicture}
\end{center}

\subsection*{Solución}

\subsubsection*{Conceptos Físicos Involucrados}

En la física conceptual, la caída libre se define como el movimiento de un objeto bajo la influencia exclusiva de la gravedad. La aceleración gravitatoria es constante, lo que significa que la velocidad del objeto no es estática, sino que cambia a un ritmo constante. Específicamente, la rapidez del cuerpo aumenta uniformemente conforme transcurre el tiempo.

La distancia que recorre un objeto en un intervalo de tiempo depende directamente de su rapidez media durante ese intervalo.

\subsubsection*{Análisis Cualitativo del Movimiento}

Al soltar el cuerpo desde el reposo, su rapidez inicial es cero. Durante el primer segundo, la gravedad acelera el cuerpo, por lo que su rapidez aumenta gradualmente desde cero hasta un valor determinado (aproximadamente $9.8 \text{ m/s}$ al final de ese primer segundo). La distancia cubierta dependerá de la rapidez media en ese lapso.

En el segundo siguiente, el cuerpo ya no parte del reposo, sino que inicia con la rapidez acumulada en el paso anterior. La gravedad sigue actuando, por lo que al final del segundo intervalo la rapidez será aún mayor. 

Dado que en cada segundo sucesivo el cuerpo viaja a una rapidez media estrictamente mayor que en el segundo anterior, es físicamente imposible que cubra la misma porción de espacio.

\subsubsection*{Conclusión}

El cuerpo recorre distancias diferentes en cada segundo de caída. Al estar acelerado por la gravedad, su rapidez aumenta continuamente, lo que provoca que en cada segundo sucesivo recorra una distancia mayor que en el segundo anterior.

\newpage

\subsection*{Enunciado}

Un cuerpo cae libremente desde el reposo durante $4 \text{ s}$. A partir de esta situación, analiza la relación entre la distancia recorrida en el primer segundo y la distancia recorrida en el último segundo del movimiento. Determina la distancia recorrida en el primer segundo y en el cuarto segundo del movimiento. Explica qué ecuaciones utilizas y por qué no debe confundirse la distancia total recorrida con la distancia recorrida en un intervalo específico de tiempo.

\vspace{0.5cm}

\begin{center}
\begin{tikzpicture}
\draw[thick, fill=gray!30] (-2,0) rectangle (2,0.5);
\filldraw[fill=blue!50, draw=black, thick] (0,-0.3) circle (0.3);

\draw[->, thick, black] (-1.5,-0.5) -- (-1.5,-6) node[left, midway] {$\vec{g}$};

\filldraw[fill=blue!50, draw=black, thick, opacity=0.3] (0,-1.5) circle (0.3);
\filldraw[fill=blue!50, draw=black, thick, opacity=0.3] (0,-4.5) circle (0.3);
\filldraw[fill=blue!50, draw=black, thick] (0,-6.5) circle (0.3);

\draw[<->, thick, red] (1,-0.3) -- (1,-1.5) node[midway, right] {$\Delta y_1$ (1er segundo)};
\draw[<->, thick, green!60!black] (1,-4.5) -- (1,-6.5) node[midway, right] {$\Delta y_4$ (4to segundo)};

\draw[dashed, gray] (0.3,-0.3) -- (1.5,-0.3) node[right, black] {$t=0$};
\draw[dashed, gray] (0.3,-1.5) -- (1.5,-1.5) node[right, black] {$t=1$};
\draw[dashed, gray] (0.3,-4.5) -- (1.5,-4.5) node[right, black] {$t=3$};
\draw[dashed, gray] (0.3,-6.5) -- (1.5,-6.5) node[right, black] {$t=4$};
\end{tikzpicture}
\end{center}

\subsection*{Solución}

\subsubsection*{Modelado Matemático y Conceptos}

Para un objeto que cae partiendo del reposo, la posición en función del tiempo está dada por la ecuación cinemática:
\begin{equation}
y(t) = \frac{1}{2} g t^2
\end{equation}
Donde $y(t)$ representa la distancia total recorrida desde el instante de soltar el objeto ($t=0$) hasta un instante de tiempo $t$. 

Es crucial no confundir la "distancia total recorrida" con la "distancia en un intervalo específico". La distancia total es acumulativa desde el inicio. Por el contrario, la distancia en un intervalo temporal específico (por ejemplo, entre el tiempo $t_a$ y el tiempo $t_b$) es la diferencia espacial entre las posiciones que el cuerpo ocupaba en esos dos instantes exactos.
\begin{equation}
\Delta y = y(t_b) - y(t_a)
\end{equation}

\subsubsection*{Cálculo de la Distancia en el Primer Segundo}

El primer segundo de movimiento comprende el intervalo de tiempo desde $t=0 \text{ s}$ hasta $t=1 \text{ s}$. Evaluamos la función de posición para encontrar el desplazamiento en este lapso, utilizando el valor estándar de la gravedad $g = 9.8 \text{ m/s}^2$.

\begin{equation}
\Delta y_1 = y(1) - y(0)
\end{equation}

\begin{equation}
\Delta y_1 = \frac{1}{2}(9.8)(1)^2 - 0
\end{equation}

\begin{equation}
\Delta y_1 = 4.9 \text{ m}
\end{equation}

\subsubsection*{Cálculo de la Distancia en el Cuarto Segundo}

El cuarto segundo (y último, según el problema) comprende el intervalo temporal estrictamente desde que el cronómetro marca $t=3 \text{ s}$ hasta que marca $t=4 \text{ s}$. Calculamos las posiciones totales para ambos instantes y las restamos:

\begin{equation}
y(3) = \frac{1}{2}(9.8)(3)^2 = 4.9 \cdot 9 = 44.1 \text{ m}
\end{equation}

\begin{equation}
y(4) = \frac{1}{2}(9.8)(4)^2 = 4.9 \cdot 16 = 78.4 \text{ m}
\end{equation}

La distancia recorrida específicamente en ese cuarto segundo es:

\begin{equation}
\Delta y_4 = y(4) - y(3) = 78.4 - 44.1
\end{equation}

\begin{equation}
\Delta y_4 = 34.3 \text{ m}
\end{equation}

\subsubsection*{Representación Gráfica del Desplazamiento}

Para ilustrar pedagógicamente la diferencia entre distancia total y distancia por intervalo, generamos una gráfica del comportamiento cuadrático. El área resaltada en el eje espacial representa los incrementos de distancia ($\Delta y$) para intervalos de tiempo iguales ($\Delta t = 1 \text{ s}$).

\begin{lstlisting}[language=Python]
import matplotlib.pyplot as plt
import numpy as np

t = np.linspace(0, 4.5, 100)
g = 9.8
y = 0.5 * g * t**2

plt.figure(figsize=(9, 6))
plt.plot(t, y, label='Trayectoria de caida', color='blue', linewidth=2)

t_1 = 1
y_1 = 0.5 * g * t_1**2
plt.vlines(t_1, 0, y_1, colors='orange', linestyles='dashed')
plt.hlines(y_1, 0, t_1, colors='orange', linestyles='dashed')
plt.fill_between(t[t<=1], y[t<=1], color='orange', alpha=0.3, label='Distancia 1er seg (4.9 m)')

t_3 = 3
y_3 = 0.5 * g * t_3**2
t_4 = 4
y_4 = 0.5 * g * t_4**2

plt.vlines(t_3, 0, y_3, colors='red', linestyles='dashed')
plt.hlines(y_3, 0, t_3, colors='red', linestyles='dashed')
plt.vlines(t_4, 0, y_4, colors='red', linestyles='dashed')
plt.hlines(y_4, 0, t_4, colors='red', linestyles='dashed')

t_interval = t[(t>=3) & (t<=4)]
plt.fill_between(t_interval, y_3, y[(t>=3) & (t<=4)], color='red', alpha=0.3, label='Distancia 4to seg (34.3 m)')

plt.gca().invert_yaxis()
plt.title('Distancia Total vs Tiempo en Caida Libre')
plt.xlabel('Tiempo (s)')
plt.ylabel('Posicion desde el origen (m)')
plt.legend(loc='lower left')
plt.grid(True, alpha=0.4)
plt.show()
\end{lstlisting}

\subsubsection*{Conclusión}

La distancia recorrida en el primer segundo es de $4.9 \text{ m}$, y la distancia recorrida en el cuarto segundo es de $34.3 \text{ m}$.

\newpage

\subsection*{Enunciado}

Un cuerpo cae libremente desde el reposo durante $4 \text{ s}$. A partir de esta situación, analiza la relación entre la distancia recorrida en el primer segundo y la distancia recorrida en el último segundo del movimiento. Establece la relación entre la distancia recorrida en el último segundo y la del primer segundo. Interpreta físicamente el resultado obtenido.

\vspace{0.5cm}

\begin{center}
\begin{tikzpicture}
\draw[thick, fill=gray!30] (-2,0) rectangle (2,0.5);
\filldraw[fill=blue!50, draw=black, thick] (0,-0.3) circle (0.3);

\draw[<->, thick, orange] (1,-0.3) -- (1,-1.5) node[midway, right] {$d_1 = 4.9 \text{ m}$};
\draw[<->, thick, red] (1,-4.5) -- (1,-6.5) node[midway, right] {$d_4 = 34.3 \text{ m}$};

\draw[dashed, gray] (0.3,-0.3) -- (1.5,-0.3);
\draw[dashed, gray] (0.3,-1.5) -- (1.5,-1.5);
\draw[dashed, gray] (0.3,-4.5) -- (1.5,-4.5);
\draw[dashed, gray] (0.3,-6.5) -- (1.5,-6.5);
\end{tikzpicture}
\end{center}

\subsection*{Solución}

\subsubsection*{Cálculo de la Proporción Matemática}

Para establecer una "relación" matemática rigurosa entre dos magnitudes, se debe efectuar el cociente (o razón) entre ellas. Esto nos permite saber cuántas veces cabe la magnitud menor dentro de la magnitud mayor. 

Procedemos a dividir la distancia del cuarto segundo ($d_4 = 34.3 \text{ m}$) entre la distancia del primer segundo ($d_1 = 4.9 \text{ m}$).

\begin{equation}
Razón = \frac{d_4}{d_1}
\end{equation}

\begin{equation}
Razón = \frac{34.3 \text{ m}}{4.9 \text{ m}}
\end{equation}

\begin{equation}
Razón = 7
\end{equation}

Este resultado es un número adimensional, exacto e independiente del valor de la aceleración de la gravedad utilizado.

\subsubsection*{Interpretación Física (Ley de Galileo)}

El resultado obtenido es una manifestación directa de la Ley de los Números Impares descubierta por Galileo Galilei. Esta ley establece que, si un objeto parte del reposo y experimenta una aceleración uniforme, las distancias recorridas durante intervalos de tiempo iguales y sucesivos son proporcionales a los números impares consecutivos (1, 3, 5, 7, 9...).

Interpretación cronológica del fenómeno:
\begin{itemize}
    \item En el 1er segundo, recorre una distancia base proporcional a $1$.
    \item En el 2do segundo, recorre una distancia proporcional a $3$.
    \item En el 3er segundo, recorre una distancia proporcional a $5$.
    \item En el 4to segundo, recorre una distancia proporcional a $7$.
\end{itemize}

Físicamente, esto significa que la rapidez media del objeto durante el cuarto segundo de su caída es exactamente siete veces mayor que la rapidez media que tuvo durante el primer segundo. Como el intervalo de tiempo es el mismo (un segundo de duración en ambos casos), viajar siete veces más rápido en promedio se traduce inevitablemente en recorrer siete veces más distancia.

\subsubsection*{Conclusión}

La relación es de $7$ a $1$. Es decir, la distancia recorrida en el último segundo es exactamente 7 veces mayor que la distancia recorrida en el primer segundo. Físicamente, esto ilustra la ley de Galileo de los números impares para el movimiento uniformemente acelerado.

\newpage

\subsection*{Enunciado}

Desde la azotea de un edificio se lanza una pelota verticalmente hacia arriba con una rapidez inicial de $20 \text{ m/s}$. Tres segundos después, ¿la pelota aún asciende, está en su altura máxima o ya desciende?

\vspace{0.5cm}

\begin{center}
\begin{tikzpicture}
\draw[thick, fill=gray!20] (0,0) rectangle (3,4);
\draw[thick] (-1,0) -- (5,0) node[right] {Suelo};

\draw[dashed, gray, thick] (1.5, 4.2) -- (3.5, 4.2) node[right, black] {Nivel de la azotea};

\filldraw[fill=red!60, draw=black, thick] (1.5, 4.2) circle (0.2);
\draw[->, thick, blue] (1.5, 4.5) -- (1.5, 6.5) node[above] {$v_0 = 20 \text{ m/s}$};

\draw[->, thick, black] (0.5, 5.5) -- (0.5, 4.5) node[left, midway] {$\vec{g}$};
\end{tikzpicture}
\end{center}

\subsection*{Solución}

\subsubsection*{Conceptos Físicos Involucrados}

El movimiento vertical de un objeto bajo la influencia exclusiva de la gravedad terrestre constituye un Movimiento Rectilíneo Uniformemente Variado (MRUV). En el marco de la física conceptual de Paul G. Hewitt, la aceleración debida a la gravedad es una constante que altera la velocidad de los cuerpos en caída libre. Para facilitar la comprensión pedagógica y el cálculo mental, esta magnitud se aproxima a $10 \text{ m/s}^2$ dirigida hacia el centro de la Tierra.

Esto significa físicamente que la fuerza gravitacional resta $10 \text{ m/s}$ a la velocidad del objeto por cada segundo que este se mueve en dirección ascendente, en contra de la atracción gravitatoria. En el análisis unidimensional, el signo de la velocidad determina el estado del movimiento:
\begin{itemize}
    \item \textbf{Velocidad positiva ($v > 0$):} El objeto asciende (se aleja del centro de la Tierra).
    \item \textbf{Velocidad nula ($v = 0$):} El objeto alcanza su altura máxima y se encuentra momentáneamente en reposo antes de invertir su dirección.
    \item \textbf{Velocidad negativa ($v < 0$):} El objeto desciende (se acerca al centro de la Tierra).
\end{itemize}

\subsubsection*{Análisis del Tiempo de Ascenso}

Para determinar el estado de la pelota a los tres segundos, primero debemos averiguar cuánto tiempo tarda en agotar su energía cinética y alcanzar su altura máxima. Utilizamos la ecuación cinemática de la velocidad:

\begin{equation}
v_f = v_0 - g \cdot t
\end{equation}

En el punto de altura máxima, la velocidad final es matemáticamente cero. Sustituyendo el valor de la rapidez inicial de $20 \text{ m/s}$ y la constante gravitatoria de $10 \text{ m/s}^2$:

\begin{equation}
0 = 20 - 10 \cdot t_{\text{subida}}
\end{equation}

Despejamos algebraicamente el tiempo necesario para la subida:

\begin{equation}
10 \cdot t_{\text{subida}} = 20
\end{equation}

\begin{equation}
t_{\text{subida}} = 2 \text{ s}
\end{equation}

Esto demuestra que la pelota requiere exactamente dos segundos para que la gravedad disipe por completo su impulso inicial hacia arriba.

\subsubsection*{Evaluación del Estado Cinemático a los Tres Segundos}

Dado que el evento por el que se nos consulta ocurre en $t = 3 \text{ s}$, y sabemos que la pelota alcanza su vértice a los $2 \text{ s}$, lógicamente se deduce que el objeto ya ha comenzado su fase de caída libre. Podemos corroborar esta deducción lógica calculando el vector velocidad exacto en ese instante:

\begin{equation}
v(3) = v_0 - g \cdot (3)
\end{equation}

\begin{equation}
v(3) = 20 - 10 \cdot (3)
\end{equation}

\begin{equation}
v(3) = 20 - 30
\end{equation}

\begin{equation}
v(3) = -10 \text{ m/s}
\end{equation}

El signo negativo del resultado confirma de manera analítica que el vector velocidad apunta hacia abajo. Durante ese tercer segundo, la pelota ha estado acelerando a favor de la gravedad partiendo desde el reposo en su altura máxima.

\subsubsection*{Representación Gráfica del Movimiento}

La representación de la velocidad en función del tiempo permite visualizar claramente el cruce de la función por el eje horizontal (cambio de dirección del movimiento). El código en Python adjunto genera un gráfico dual que relaciona la pérdida de velocidad con la trayectoria parabólica espacial.

\begin{lstlisting}[language=Python]
import matplotlib.pyplot as plt
import numpy as np

t = np.linspace(0, 4, 100)
v = 20 - 10 * t
y = 20 * t - 0.5 * 10 * t**2

fig, ax1 = plt.subplots(figsize=(9, 6))

color_v = 'tab:blue'
ax1.set_xlabel('Tiempo (s)')
ax1.set_ylabel('Velocidad (m/s)', color=color_v)
ax1.plot(t, v, color=color_v, linewidth=2.5, label='Velocidad')
ax1.tick_params(axis='y', labelcolor=color_v)
ax1.axhline(0, color='black', linewidth=1, linestyle='dashed')

ax2 = ax1.twinx()
color_y = 'tab:red'
ax2.set_ylabel('Desplazamiento desde la azotea (m)', color=color_y)
ax2.plot(t, y, color=color_y, linewidth=2.5, linestyle='-.', label='Posicion')
ax2.tick_params(axis='y', labelcolor=color_y)

ax1.axvline(2, color='gray', linestyle='dotted', alpha=0.7)
ax1.axvline(3, color='gray', linestyle='dotted', alpha=0.7)

ax1.plot(2, 0, 'ko', markersize=6)
ax1.text(2.1, 1, 'Altura maxima (t=2s, v=0)', fontsize=10)

ax1.plot(3, -10, 'ko', markersize=6)
ax1.text(3.1, -9, 'Descenso (t=3s, v=-10 m/s)', fontsize=10)

plt.title('Evolucion de la Velocidad y Posicion de la Pelota')
fig.tight_layout()
plt.grid(True, alpha=0.3)
plt.show()
\end{lstlisting}

\subsubsection*{Conclusión}

Tres segundos después del lanzamiento, la pelota ya desciende.

\newpage

\subsection*{Enunciado}

Un hombre lanza una pelota verticalmente hacia arriba. Dos segundos más tarde, lanza una segunda pelota con la misma velocidad inicial que la primera. Se observa que ambas pelotas chocan 0.4 s después de que la segunda pelota fue lanzada. ¿Cuál es la velocidad inicial de ambas pelotas?

\vspace{0.5cm}

\begin{center}
\begin{tikzpicture}
\draw[thick] (-2,0) -- (4,0) node[right] {Suelo};

\filldraw[fill=blue!70, draw=black, thick] (0,0.2) circle (0.2) node[below=0.3cm] {Pelota A ($t=0$)};
\draw[->, thick, orange] (0,0.5) -- (0,2.5) node[left] {$v_0$};
\draw[dashed, orange] (0,2.5) .. controls (0,4) and (1,4) .. (1,2.5);

\filldraw[fill=purple!70, draw=black, thick] (2,0.2) circle (0.2) node[below=0.3cm] {Pelota B ($t=2 \text{ s}$)};
\draw[->, thick, purple] (2,0.5) -- (2,2.5) node[right] {$v_0$};

\filldraw[fill=blue!70, draw=black, thick] (0.8, 2.7) circle (0.2);
\filldraw[fill=purple!70, draw=black, thick] (1.2, 2.7) circle (0.2);
\node[above] at (1, 3.0) {Choque a los $0.4 \text{ s}$ de lanzar B};

\end{tikzpicture}
\end{center}

\subsection*{Solución}

\subsubsection*{Conceptos Físicos Involucrados}

De acuerdo con el enfoque conceptual, cuando un objeto es lanzado verticalmente hacia arriba, este experimenta una aceleración constante hacia abajo debido a la gravedad terrestre. Esta aceleración reduce uniformemente la velocidad del objeto mientras asciende y, posteriormente, aumenta su velocidad mientras desciende. 

En este problema interactúan dos proyectiles idénticos sujetos al mismo campo gravitatorio. La clave pedagógica reside en comprender que, aunque comparten el mismo espacio físico y la misma velocidad inicial, operan en diferentes marcos temporales debido al retraso en el lanzamiento del segundo objeto. Para resolverlo, se debe establecer una sincronización temporal que permita comparar sus posiciones en el mismo instante absoluto.

\subsubsection*{Sincronización del Sistema de Referencia Temporal}

Se establece el nivel del suelo como el origen de coordenadas espaciales ($y = 0$). Para el dominio del tiempo, se define el instante $t = 0$ como el momento exacto en el que el hombre lanza la primera pelota (Pelota A). 

Dado que la segunda pelota (Pelota B) se lanza dos segundos después, su tiempo de vuelo siempre será dos segundos menor que el de la Pelota A. El problema indica que el choque ocurre $0.4 \text{ s}$ después de que la segunda pelota es lanzada. Esto nos permite deducir el tiempo total de vuelo de cada pelota hasta el instante de la colisión:

Para la Pelota B:
\begin{equation}
t_B = 0.4 \text{ s}
\end{equation}

Para la Pelota A (que lleva dos segundos adicionales en el aire):
\begin{equation}
t_A = 2.0 \text{ s} + 0.4 \text{ s} = 2.4 \text{ s}
\end{equation}

\subsubsection*{Ecuaciones de Movimiento}

La posición vertical de cualquier objeto en caída libre o tiro vertical está regida por la ecuación cuadrática de la cinemática:
\begin{equation}
y(t) = y_0 + v_0 t - \frac{1}{2} g t^2
\end{equation}

Aplicamos esta ecuación a cada pelota en su respectivo tiempo de vuelo. Se utilizará la aproximación conceptual de la gravedad $g = 10 \text{ m/s}^2$, por lo que el factor $\frac{1}{2}g$ equivale a $5 \text{ m/s}^2$.

Posición de la Pelota A en el momento del choque:
\begin{equation}
y_A = v_0 (2.4) - 5 (2.4)^2
\end{equation}

Posición de la Pelota B en el momento del choque:
\begin{equation}
y_B = v_0 (0.4) - 5 (0.4)^2
\end{equation}

\subsubsection*{Condición de Choque y Desarrollo Matemático}

Para que ocurra una colisión, ambas pelotas deben ocupar exactamente la misma coordenada espacial en el mismo instante. Matemáticamente, esto significa igualar sus posiciones finales:

\begin{equation}
y_A = y_B
\end{equation}

Sustituimos las ecuaciones estructuradas previamente:

\begin{equation}
v_0 (2.4) - 5 (2.4)^2 = v_0 (0.4) - 5 (0.4)^2
\end{equation}

Desarrollamos los cuadrados de los tiempos:

\begin{equation}
2.4 v_0 - 5 (5.76) = 0.4 v_0 - 5 (0.16)
\end{equation}

Realizamos las multiplicaciones aritméticas:

\begin{equation}
2.4 v_0 - 28.8 = 0.4 v_0 - 0.8
\end{equation}

Agrupamos los términos que contienen la incógnita $v_0$ en el lado izquierdo de la ecuación, y las constantes numéricas en el lado derecho:

\begin{equation}
2.4 v_0 - 0.4 v_0 = 28.8 - 0.8
\end{equation}

\begin{equation}
2.0 v_0 = 28.0
\end{equation}

Finalmente, despejamos la velocidad inicial dividiendo entre dos:

\begin{equation}
v_0 = \frac{28}{2}
\end{equation}

\begin{equation}
v_0 = 14 \text{ m/s}
\end{equation}

\subsubsection*{Representación Gráfica del Encuentro}

Graficar la altura de ambas pelotas en función del tiempo maestro (el tiempo desde que se lanzó la primera pelota) ilustra cómo las dos trayectorias parabólicas se intersecan. El siguiente código en Python modela este cruce cinemático.

\begin{lstlisting}[language=Python]
import matplotlib.pyplot as plt
import numpy as np

t_maestro = np.linspace(0, 2.6, 200)

y_A = 14 * t_maestro - 5 * t_maestro**2

y_B = np.zeros_like(t_maestro)
mask = t_maestro >= 2.0
y_B[mask] = 14 * (t_maestro[mask] - 2.0) - 5 * (t_maestro[mask] - 2.0)**2

plt.figure(figsize=(8, 5))
plt.plot(t_maestro, y_A, label='Pelota A (lanzada en t=0)', color='blue', linewidth=2)
plt.plot(t_maestro[mask], y_B[mask], label='Pelota B (lanzada en t=2s)', color='purple', linewidth=2)

plt.axvline(2.4, color='black', linestyle='dashed', alpha=0.5)
plt.axhline(4.8, color='black', linestyle='dashed', alpha=0.5)

plt.plot(2.4, 4.8, 'ro', markersize=8)
plt.text(1.3, 5.2, 'Choque (t=2.4 s, y=4.8 m)', fontsize=11, weight='bold', color='red')

plt.title('Altura vs Tiempo de las Pelotas A y B')
plt.xlabel('Tiempo desde el primer lanzamiento (s)')
plt.ylabel('Altura (m)')
plt.legend()
plt.grid(True)
plt.show()
\end{lstlisting}

\subsubsection*{Conclusión}

La velocidad inicial con la que se lanzaron ambas pelotas es de 14 m/s.

\newpage

\subsection*{Enunciado}

Después de soltarse de un helicóptero, un paracaidista cae 80 m en forma libre, y abre en ese instante el paracaídas, lo cual le produce un retroceso en su velocidad de 2 m/s cada segundo, llegando al suelo con una rapidez de 2 m/s hacia abajo. ¿Cuánto tiempo estuvo en el aire?

\vspace{0.5cm}

\begin{center}
\begin{tikzpicture}
\filldraw[black] (0, 0) circle (2pt) node[right] {A (Helicóptero, $t=0$)};
\filldraw[black] (0, -3) circle (2pt) node[right] {B (Abre paracaídas)};
\filldraw[black] (0, -6) circle (2pt) node[right] {C (Suelo)};

\draw[dashed, thick, gray] (0,0) -- (0,-6);

\draw[<->, thick] (-1, 0) -- (-1, -3) node[midway, left] {$\Delta y_1 = 80 \text{ m}$};
\draw[<->, thick, dashed] (-1, -3) -- (-1, -6) node[midway, left] {$\Delta y_2$};

\draw[->, thick, red] (1, -1) -- (1, -2) node[midway, right] {$\vec{g} = 10 \text{ m/s}^2$ (hacia abajo)};
\draw[->, thick, blue] (0, -1.2) -- (0, -2.2) node[midway, left] {$\vec{v}$};

\draw[->, thick, red] (1, -4.5) -- (1, -3.5) node[midway, right] {$\vec{a} = 2 \text{ m/s}^2$ (hacia arriba)};
\draw[->, thick, blue] (0, -4) -- (0, -5) node[midway, left] {$\vec{v}$};

\draw[thick] (-2, -6) -- (2, -6) node[right] {Nivel del suelo};
\end{tikzpicture}
\end{center}

\subsection*{Solución}

\subsubsection*{Conceptos Físicos Involucrados}

El movimiento del paracaidista se divide en dos fases cinemáticas distintas. De acuerdo con la física conceptual, la primera fase es una caída libre pura, donde la única fuerza que actúa sobre el sistema es la gravedad de la Tierra. En esta fase, la aceleración es constante y está dirigida hacia el centro del planeta.

La segunda fase inicia cuando se despliega el paracaídas. El área superficial masiva del paracaídas interactúa con las moléculas del aire, generando una enorme fuerza de resistencia aerodinámica que supera momentáneamente al peso del paracaidista. Según la Segunda Ley de Newton, esta fuerza neta hacia arriba produce una aceleración ascendente (referida en el problema como un retroceso). Esta aceleración no invierte la dirección del movimiento, sino que reduce progresivamente la magnitud de la velocidad de caída hasta alcanzar una velocidad final más segura para el impacto.

\subsubsection*{Definición del Sistema de Referencia}

Para modelar matemáticamente este fenómeno, estableceremos un eje de coordenadas vertical. Situaremos el origen en el punto exacto donde el paracaidista se suelta del helicóptero. Definiremos la dirección hacia arriba como positiva y la dirección hacia abajo como negativa. Esta convención de signos es fundamental para evitar errores algebraicos en las ecuaciones cinemáticas.

\subsubsection*{Análisis de la Primera Fase (Caída Libre)}

Durante los primeros ochenta metros, el paracaidista acelera por efecto exclusivo de la gravedad. Extraemos los parámetros con sus signos correspondientes:
\begin{itemize}
    \item Velocidad inicial $v_A = 0 \text{ m/s}$
    \item Desplazamiento $\Delta y_1 = -80 \text{ m}$
    \item Aceleración gravitatoria $a_y = -10 \text{ m/s}^2$
\end{itemize}

Utilizamos la ecuación cinemática independiente del tiempo para calcular la velocidad que alcanza justo antes de abrir el paracaídas (punto B):
\begin{equation}
v_B^2 = v_A^2 + 2 \cdot a_y \cdot \Delta y_1
\end{equation}

Sustituimos los valores numéricos:
\begin{equation}
v_B^2 = 0^2 + 2 \cdot (-10) \cdot (-80)
\end{equation}

\begin{equation}
v_B^2 = 1600
\end{equation}

Al extraer la raíz cuadrada, obtenemos dos posibles soluciones matemáticas. Físicamente, dado que el paracaidista se mueve hacia abajo, seleccionamos la raíz negativa:
\begin{equation}
v_B = -\sqrt{1600} = -40 \text{ m/s}
\end{equation}

Ahora, calculamos el tiempo que duró esta primera fase utilizando la ecuación de definición de la aceleración:
\begin{equation}
v_B = v_A + a_y \cdot t_1
\end{equation}

\begin{equation}
-40 = 0 + (-10) \cdot t_1
\end{equation}

Despejamos el tiempo dividiendo la ecuación:
\begin{equation}
t_1 = \frac{-40}{-10} = 4 \text{ s}
\end{equation}

\subsubsection*{Análisis de la Segunda Fase (Frenado Aerodinámico)}

La segunda fase comienza con las condiciones finales de la primera fase. El paracaídas induce una aceleración hacia arriba, por lo que su signo debe ser estrictamente positivo.

\begin{itemize}
    \item Velocidad inicial de la fase $v_B = -40 \text{ m/s}$
    \item Velocidad final en el suelo $v_C = -2 \text{ m/s}$
    \item Aceleración de frenado $a = +2 \text{ m/s}^2$
\end{itemize}

Buscamos el tiempo que requiere esta reducción de velocidad. Aplicamos nuevamente la ecuación lineal de cinemática:
\begin{equation}
v_C = v_B + a \cdot t_2
\end{equation}

Sustituimos los valores con sus signos correctos:
\begin{equation}
-2 = -40 + 2 \cdot t_2
\end{equation}

Transponemos el término de la velocidad inicial al lado izquierdo sumando cuarenta:
\begin{equation}
-2 + 40 = 2 \cdot t_2
\end{equation}

\begin{equation}
38 = 2 \cdot t_2
\end{equation}

Despejamos el tiempo de la segunda fase:
\begin{equation}
t_2 = \frac{38}{2} = 19 \text{ s}
\end{equation}

\subsubsection*{Cálculo del Tiempo Total}

El tiempo total que el paracaidista estuvo suspendido en el aire es la suma aritmética de los intervalos temporales de ambas fases.

\begin{equation}
t_{\text{total}} = t_1 + t_2
\end{equation}

\begin{equation}
t_{\text{total}} = 4 + 19
\end{equation}

\begin{equation}
t_{\text{total}} = 23 \text{ s}
\end{equation}

\subsubsection*{Representación Gráfica del Movimiento}

La gráfica de la velocidad vertical en función del tiempo permite visualizar el dramático cambio en la pendiente que ocurre en el instante en que se despliega el paracaídas. La pendiente de la primera recta es muy pronunciada debido a la intensidad de la gravedad, mientras que la segunda presenta una inclinación inversa y más suave que estabiliza la caída.

\begin{lstlisting}[language=Python]
import matplotlib.pyplot as plt
import numpy as np

t1 = np.linspace(0, 4, 100)
v1 = -10 * t1

t2 = np.linspace(4, 23, 100)
v2 = -40 + 2 * (t2 - 4)

plt.figure(figsize=(9, 6))

plt.plot(t1, v1, label='Fase 1: Caida libre (g = -10 m/s^2)', color='blue', linewidth=2.5)
plt.plot(t2, v2, label='Fase 2: Frenado (a = +2 m/s^2)', color='red', linewidth=2.5)

plt.axvline(4, color='black', linestyle='dashed', alpha=0.5)
plt.axhline(0, color='black', linewidth=1)

plt.plot(4, -40, 'ko', markersize=6)
plt.text(4.5, -38, 'Apertura del paracaidas\n(4 s, -40 m/s)', fontsize=10)

plt.plot(23, -2, 'ko', markersize=6)
plt.text(15, -4, 'Llegada al suelo\n(23 s, -2 m/s)', fontsize=10)

plt.title('Velocidad Vertical vs Tiempo de Vuelo del Paracaidista')
plt.xlabel('Tiempo (s)')
plt.ylabel('Velocidad (m/s)')
plt.legend()
plt.grid(True, alpha=0.4)
plt.show()
\end{lstlisting}

\subsubsection*{Conclusión}

El paracaidista estuvo en el aire un tiempo total de 23 s.

\newpage

Se lanza un proyectil con una velocidad $\vec{v}_0 = 10\vec{i} + 5\vec{j} \text{ m/s}$. A cierto instante su velocidad es $\vec{v} = 10\vec{i} - 5\vec{j} \text{ m/s}$. Determine para este instante el vector posición con respecto al punto de lanzamiento. Considere la aceleración de la gravedad como $g = 9.8 \text{ m/s}^2$.

\vspace{0.5cm}

\begin{center}
\begin{tikzpicture}
\draw[thick, ->] (0,0) -- (6,0) node[right] {$x$};
\draw[thick, ->] (0,0) -- (0,3) node[above] {$y$};
\draw[dashed, blue, thick] (0,0) parabola bend (2.5, 2.0) (5,0);
\filldraw[black] (0,0) circle (2pt) node[below left] {Inicio};
\draw[->, thick, green!60!black] (0,0) -- (1.5,1) node[above left] {$\vec{v}_0 = 10\vec{i} + 5\vec{j}$};
\filldraw[black] (5,0) circle (2pt) node[below right] {Instante final};
\draw[->, thick, red] (5,0) -- (6.5,-1) node[below left] {$\vec{v} = 10\vec{i} - 5\vec{j}$};
\end{tikzpicture}
\end{center}

\subsection*{Solución}

\subsubsection*{Conceptos Físicos Involucrados}

El movimiento de proyectiles es bidimensional y se analiza de forma pedagógica separando sus ejes espaciales. Horizontalmente, el proyectil experimenta un Movimiento Rectilíneo Uniforme porque la resistencia del aire es nula, manteniendo su velocidad constante. Verticalmente, experimenta un Movimiento Rectilíneo Uniformemente Variado bajo la influencia de la aceleración constante de la gravedad.

\subsubsection*{Análisis de Simetría}

La velocidad inicial del lanzamiento tiene una componente vertical ascendente de $5\vec{j} \text{ m/s}$. En el instante que se nos solicita evaluar, el vector velocidad posee una componente vertical descendente de $-5\vec{j} \text{ m/s}$. 

De acuerdo con el principio de conservación de la energía mecánica y la simetría cinemática en el eje vertical, cuando un proyectil alcanza una velocidad de igual magnitud pero en dirección diametralmente opuesta, significa que ha retornado físicamente a su nivel topográfico de origen. Por ende, deducimos directamente que el desplazamiento en el eje de las ordenadas es nulo ($y = 0$).

\subsubsection*{Cálculo del Tiempo y Desplazamiento}

Para estructurar el vector posición, debemos calcular la coordenada del desplazamiento horizontal. Primero encontramos el tiempo transcurrido evaluando la ecuación de velocidad vertical:

$$v_y = v_{0y} - g \cdot t$$

Sustituimos las componentes verticales conocidas:

$$-5 = 5 - 9.8 \cdot t$$

$$-10 = -9.8 \cdot t$$

$$t = \frac{10}{9.8} \approx 1.02 \text{ s}$$

Teniendo la variable temporal del instante, calculamos la coordenada horizontal sabiendo que la velocidad en dicho eje se ha mantenido constante a $10 \text{ m/s}$:

$$x = v_{0x} \cdot t$$

$$x = 10 \cdot 1.02$$

$$x = 10.20 \text{ m}$$

\subsubsection*{Representación Gráfica Computacional}

El modelado computacional permite visualizar la parábola del desplazamiento y cómo el vector posición solicitado nace en el origen y se posa sobre el eje de las abscisas debido al retorno altimétrico del proyectil.

\begin{lstlisting}[language=Python]
import matplotlib.pyplot as plt
import numpy as np

t_array = np.linspace(0, 1.02, 100)
x_pos = 10 * t_array
y_pos = 5 * t_array - 0.5 * 9.8 * t_array**2

plt.figure(figsize=(8, 4))
plt.plot(x_pos, y_pos, label='Trayectoria parabolica', color='blue', linewidth=2)
plt.quiver(0, 0, 10.2, 0, angles='xy', scale_units='xy', scale=1, color='green', width=0.005, label='Vector Posicion')

plt.title('Vector Posicion en el Instante de Retorno')
plt.xlabel('Posicion horizontal x (m)')
plt.ylabel('Posicion vertical y (m)')
plt.axhline(0, color='black', linewidth=1)
plt.legend()
plt.grid(True)
plt.show()
\end{lstlisting}

\subsubsection*{Conclusión}

El vector posición con respecto al punto de lanzamiento es $\vec{r} = 10.20\vec{i} + 0\vec{j} \text{ m}$.

\newpage

\subsection*{Enunciado}

Se lanza un proyectil con una velocidad $\vec{v}_0 = 10\vec{i} + 5\vec{j} \text{ m/s}$. A cierto instante su velocidad es $\vec{v} = 10\vec{i} - 5\vec{j} \text{ m/s}$. Determine para este instante el tiempo de subida. Considere la aceleración de la gravedad como $g = 9.8 \text{ m/s}^2$.

\vspace{0.5cm}

\begin{center}
\begin{tikzpicture}
\draw[thick, ->] (0,0) -- (6,0) node[right] {$x$};
\draw[dashed, blue, thick] (0,0) parabola bend (2.5, 2.0) (5,0);
\filldraw[black] (2.5,2.0) circle (2pt) node[above] {Altura Máxima ($v_y = 0$)};
\draw[->, thick, orange] (0,0) parabola bend (2.5, 2.0) (2.5, 2.0);
\draw[->, thick, green!60!black] (2.5,2.0) -- (4,2.0) node[above] {$\vec{v}_x = 10\vec{i}$};
\end{tikzpicture}
\end{center}

\subsection*{Solución}

\subsubsection*{Condición de Altura Máxima}

El tiempo de subida se conceptualiza como el lapso absoluto que invierte el proyectil en escalar hasta el vértice de su trayectoria parabólica. En esta cúspide geométrica, el vector gravitacional ha logrado contrarrestar y agotar el total del impulso cinético ascendente. Esto establece una frontera de modelado donde la componente vertical de la velocidad se asume momentáneamente nula.

\subsubsection*{Desarrollo Matemático}

Planteamos la ecuación cinemática para la velocidad vertical, estableciendo su valor final en cero:

$$v_y = v_{0y} - g \cdot t_{\text{subida}}$$

Sustituimos el valor de la componente inicial del lanzamiento y el de la gravedad terrestre:

$$0 = 5 - 9.8 \cdot t_{\text{subida}}$$

Trasnponemos el factor acelerado hacia el lado opuesto de la igualdad:

$$9.8 \cdot t_{\text{subida}} = 5$$

Despejamos el tiempo dividiendo ambas entidades numéricas:

$$t_{\text{subida}} = \frac{5}{9.8}$$

$$t_{\text{subida}} \approx 0.51 \text{ s}$$

\subsubsection*{Conclusión}

El tiempo de subida es de $0.51 \text{ s}$.

\newpage

\subsection*{Enunciado}

Se lanza un proyectil con una velocidad $\vec{v}_0 = 10\vec{i} + 5\vec{j} \text{ m/s}$. A cierto instante su velocidad es $\vec{v} = 10\vec{i} - 5\vec{j} \text{ m/s}$. Determine para este instante el tiempo de bajada. Considere la aceleración de la gravedad como $g = 9.8 \text{ m/s}^2$.

\vspace{0.5cm}

\begin{center}
\begin{tikzpicture}
\draw[thick, ->] (0,0) -- (6,0) node[right] {$x$};
\draw[dashed, blue, thick] (0,0) parabola bend (2.5, 2.0) (5,0);
\filldraw[black] (2.5,2.0) circle (2pt) node[above] {Inicio bajada};
\filldraw[black] (5,0) circle (2pt) node[right] {Fin bajada};
\draw[->, thick, purple] (2.5,2.0) parabola bend (2.5, 2.0) (5,0);
\end{tikzpicture}
\end{center}

\subsection*{Solución}

\subsubsection*{Simetría de la Trayectoria}

El tiempo de bajada enmarca el intervalo cronológico en el cual el proyectil abandona su altura máxima y desciende hasta igualar el nivel donde se encuentra el instante estudiado por el enunciado. Al iniciar el descenso, la velocidad vertical es efectivamente cero, y el cuerpo se precipita en caída libre pura a favor del vector de la gravedad hasta recuperar una velocidad escalar de $5 \text{ m/s}$ dirigida hacia la Tierra.

\subsubsection*{Desarrollo Matemático}

Utilizamos nuevamente la ecuación que rige la velocidad vertical lineal, adaptando las condiciones limítrofes. La velocidad inicial del tramo será el reposo vertical de la cúspide, y la velocidad final será la componente descendente dada por la métrica del problema.

$$v_{y\text{final}} = v_{y\text{pico}} - g \cdot t_{\text{bajada}}$$

$$-5 = 0 - 9.8 \cdot t_{\text{bajada}}$$

$$-5 = -9.8 \cdot t_{\text{bajada}}$$

Despejamos el tiempo ejecutando el cociente algebraico, el cual neutraliza los signos dimensionales:

$$t_{\text{bajada}} = \frac{-5}{-9.8}$$

$$t_{\text{bajada}} \approx 0.51 \text{ s}$$

Esta demostración matemática evidencia la armonía del movimiento parabólico: el tiempo necesario para mermar una velocidad es rigurosamente idéntico al tiempo para reconstituirla.

\subsubsection*{Conclusión}

El tiempo de bajada es de $0.51 \text{ s}$.

\newpage

\subsection*{Enunciado}

Se lanza un proyectil con una velocidad $\vec{v}_0 = 10\vec{i} + 5\vec{j} \text{ m/s}$. A cierto instante su velocidad es $\vec{v} = 10\vec{i} - 5\vec{j} \text{ m/s}$. Determine para este instante el tiempo de vuelo. Considere la aceleración de la gravedad como $g = 9.8 \text{ m/s}^2$.

\vspace{0.5cm}

\begin{center}
\begin{tikzpicture}
\draw[thick, ->] (0,0) -- (6,0) node[right] {$x$};
\draw[->, thick, cyan!80!blue] (0,0) parabola bend (2.5, 2.0) (5,0);
\node[above] at (2.5, 2.2) {Tiempo total de vuelo};
\filldraw[black] (0,0) circle (2pt) node[below] {$t = 0$};
\filldraw[black] (5,0) circle (2pt) node[below] {$t = t_{\text{vuelo}}$};
\end{tikzpicture}
\end{center}

\subsection*{Solución}

\subsubsection*{Composición del Tiempo Total}

El tiempo de vuelo total que requiere un objeto volador para transitar entre dos etapas definidas de su arco es la consolidación aditiva de todos los micro-lapsos que conforman su trayecto inercial. Se entiende pedagógicamente como la variable cronológica que soporta todo el trabajo modificador de la gravedad sobre el vector original.

\subsubsection*{Desarrollo Matemático}

Teniendo el fraccionamiento de las fases cinemáticas ya resueltas en la simetría del arco, simplemente ejecutamos una suma aritmética para englobar el evento global:

$$t_{\text{vuelo}} = t_{\text{subida}} + t_{\text{bajada}}$$

$$t_{\text{vuelo}} = 0.51 + 0.51$$

$$t_{\text{vuelo}} = 1.02 \text{ s}$$

Alternativamente, para dotar al análisis de una mayor robustez científica, evaluamos el impacto total directo de la gravedad sobre el salto vectorial absoluto, midiendo el intervalo de punta a punta:

$$v_{y\text{final}} = v_{0y} - g \cdot t_{\text{vuelo}}$$

$$-5 = 5 - 9.8 \cdot t_{\text{vuelo}}$$

$$-10 = -9.8 \cdot t_{\text{vuelo}}$$

$$t_{\text{vuelo}} = \frac{10}{9.8} \approx 1.02 \text{ s}$$

Ambos enfoques conceptuales convergen rigurosamente en el mismo valor cronológico absoluto.

\subsubsection*{Conclusión}

El tiempo de vuelo para este instante es de $1.02 \text{ s}$.

\newpage

\subsection*{Enunciado}

Desde lo alto de un edificio de $100 \text{ m}$ de altura, se dispara un proyectil hacia un terreno horizontal con una rapidez de $60 \text{ m/s}$ y $30^\circ$ sobre la horizontal. Determine el tiempo de movimiento.

\vspace{0.5cm}

\begin{center}
\begin{tikzpicture}
\draw[thick, fill=gray!20] (0,0) rectangle (1.5,4);
\draw[thick] (0,0) -- (8,0) node[right] {Suelo};

\draw[dashed, blue, thick] (1.5,4) parabola bend (3.5, 5.5) (7.5,0);

\filldraw[black] (1.5,4) circle (2pt);
\draw[->, thick, green!60!black] (1.5,4) -- (3,4.866) node[above left] {$v_0 = 60 \text{ m/s}$};
\draw[dashed, gray] (1.5,4) -- (3.5,4);
\draw (2.3,4) arc (0:30:0.8);
\node at (2.6,4.2) {$30^\circ$};

\draw[<->, thick] (-0.5,0) -- (-0.5,4) node[midway, fill=white] {$100 \text{ m}$};
\draw[dashed, gray] (0,4) -- (-0.8,4);
\draw[dashed, gray] (0,0) -- (-0.8,0);
\end{tikzpicture}
\end{center}

\subsection*{Solución}

\subsubsection*{Conceptos Físicos Involucrados}

De acuerdo con la física conceptual, el movimiento de un proyectil es bidimensional y debe analizarse aplicando el principio de independencia de los movimientos. Este principio establece que el movimiento horizontal (que ocurre a velocidad constante por carecer de aceleración) y el movimiento vertical (que ocurre bajo la aceleración constante de la gravedad) ocurren simultáneamente pero no se afectan entre sí. 

El "tiempo de movimiento" está dictado exclusivamente por el eje vertical. El proyectil permanecerá en el aire el tiempo exacto que le tome a la gravedad frenar su ascenso, invertir su dirección y obligarlo a recorrer toda la distancia vertical hasta chocar con el suelo.

\subsubsection*{Descomposición Vectorial de la Velocidad}

El primer paso analítico es descomponer el vector de velocidad inicial en sus componentes ortogonales utilizando trigonometría básica.

Para la componente vertical inicial:
\begin{equation}
v_{0y} = v_0 \cdot \sin(\theta)
\end{equation}

\begin{equation}
v_{0y} = 60 \cdot \sin(30^\circ) = 60 \cdot 0.5 = 30 \text{ m/s}
\end{equation}

Para la componente horizontal inicial:
\begin{equation}
v_{0x} = v_0 \cdot \cos(\theta)
\end{equation}

\begin{equation}
v_{0x} = 60 \cdot \cos(30^\circ) = 60 \cdot \frac{\sqrt{3}}{2} = 30\sqrt{3} \approx 51.96 \text{ m/s}
\end{equation}

\subsubsection*{Definición del Sistema de Referencia y Ecuación de Posición}

Situamos el origen de coordenadas ($y = 0$) en el nivel del suelo. Bajo este marco, la posición inicial del proyectil es $y_0 = 100 \text{ m}$. La dirección hacia arriba es positiva, por lo que la aceleración de la gravedad adopta un signo negativo ($g = -9.8 \text{ m/s}^2$).

Planteamos la ecuación cinemática para la posición vertical:
\begin{equation}
y(t) = y_0 + v_{0y}t - \frac{1}{2}gt^2
\end{equation}

El movimiento termina exactamente cuando el proyectil toca el suelo, es decir, cuando su posición final es $y(t) = 0$. Sustituimos los valores conocidos:
\begin{equation}
0 = 100 + 30t - \frac{1}{2}(9.8)t^2
\end{equation}

\begin{equation}
0 = 100 + 30t - 4.9t^2
\end{equation}

\subsubsection*{Desarrollo Matemático}

Reordenamos la ecuación cuadrática a su forma estándar multiplicando por menos uno para facilitar la resolución:
\begin{equation}
4.9t^2 - 30t - 100 = 0
\end{equation}

Aplicamos la fórmula resolvente para ecuaciones de segundo grado:
\begin{equation}
t = \frac{-b \pm \sqrt{b^2 - 4ac}}{2a}
\end{equation}

\begin{equation}
t = \frac{-(-30) \pm \sqrt{(-30)^2 - 4(4.9)(-100)}}{2(4.9)}
\end{equation}

\begin{equation}
t = \frac{30 \pm \sqrt{900 + 1960}}{9.8}
\end{equation}

\begin{equation}
t = \frac{30 \pm \sqrt{2860}}{9.8}
\end{equation}

Evaluamos la raíz cuadrada ($\sqrt{2860} \approx 53.48$). Como el tiempo debe ser una magnitud estrictamente positiva, descartamos la resta y operamos con la suma:
\begin{equation}
t = \frac{30 + 53.48}{9.8}
\end{equation}

\begin{equation}
t = \frac{83.48}{9.8} \approx 8.518 \text{ s}
\end{equation}

\subsubsection*{Representación Gráfica Computacional}

Visualizar la altura en función del tiempo permite comprobar que la curva inicia en $100 \text{ m}$, asciende brevemente hasta su cúspide y luego cruza el eje temporal pasados los ocho segundos y medio.

\begin{lstlisting}[language=Python]
import matplotlib.pyplot as plt
import numpy as np

t_array = np.linspace(0, 9, 100)
y_pos = 100 + 30 * t_array - 4.9 * t_array**2

plt.figure(figsize=(8, 5))
plt.plot(t_array, y_pos, color='blue', linewidth=2.5)

plt.axhline(0, color='black', linewidth=1.5)
plt.axvline(8.52, color='red', linestyle='dashed', alpha=0.6)

plt.plot(8.52, 0, 'ro', markersize=7)
plt.text(6.5, 10, 'Impacto\nt = 8.52 s', fontsize=11, color='red')

plt.title('Altura del Proyectil en funcion del Tiempo')
plt.xlabel('Tiempo de vuelo (s)')
plt.ylabel('Altura desde el suelo (m)')
plt.grid(True, alpha=0.4)
plt.show()
\end{lstlisting}

\subsubsection*{Conclusión}

El tiempo de movimiento del proyectil es de $8.52 \text{ s}$.

\newpage

\subsection*{Enunciado}

Desde lo alto de un edificio de $100 \text{ m}$ de altura, se dispara un proyectil hacia un terreno horizontal con una rapidez de $60 \text{ m/s}$ y $30^\circ$ sobre la horizontal. Determine el tiempo de vuelo.

\vspace{0.5cm}

\begin{center}
\begin{tikzpicture}
\draw[thick, fill=gray!20] (0,0) rectangle (1.5,4);
\draw[thick] (0,0) -- (8,0) node[right] {Suelo};

\draw[dashed, blue, thick] (1.5,4) parabola bend (3.5, 5.5) (7.5,0);

\filldraw[black] (1.5,4) circle (2pt);
\draw[->, thick, green!60!black] (1.5,4) -- (3,4.866) node[above left] {$v_0 = 60 \text{ m/s}$};
\draw[dashed, gray] (1.5,4) -- (3.5,4);
\draw (2.3,4) arc (0:30:0.8);
\node at (2.6,4.2) {$30^\circ$};

\draw[<->, thick] (-0.5,0) -- (-0.5,4) node[midway, fill=white] {$100 \text{ m}$};
\draw[dashed, gray] (0,4) -- (-0.8,4);
\draw[dashed, gray] (0,0) -- (-0.8,0);
\end{tikzpicture}
\end{center}

\subsection*{Solución}

\subsubsection*{Conceptos Físicos Involucrados}

En la cinemática de proyectiles, la terminología es crucial. El "tiempo de vuelo" hace referencia a la duración cronológica total en la que el objeto se encuentra en el aire, desvinculado de las superficies de soporte. Físicamente, este concepto es un sinónimo exacto del "tiempo de movimiento" calculado para la trayectoria completa.

El objetivo pedagógico de separar estas terminologías suele ser evaluar la comprensión del estudiante sobre las fórmulas prefabricadas. Es un error común intentar calcular el tiempo de vuelo utilizando la ecuación simplificada $T = \frac{2v_{0y}}{g}$. Dicha fórmula es un caso particular derivado de la simetría geométrica y es válida única y exclusivamente cuando el proyectil aterriza al mismo nivel topográfico desde el cual fue lanzado ($y_0 = y_f$). 

Dado que nuestro proyectil es lanzado desde un edificio de cien metros y aterriza en el suelo, la trayectoria es radicalmente asimétrica. El tiempo de bajada será notablemente mayor que el tiempo de subida. Por consiguiente, el tiempo de vuelo debe ser extraído mediante la resolución general de la ecuación cuadrática de posición vertical.

\subsubsection*{Validación Analítica}

Retomamos el modelo cuadrático estructurado a partir del suelo ($y=0$):
\begin{equation}
y(t) = 100 + 30t - 4.9t^2
\end{equation}

Como ya se demostró rigurosamente, la única raíz física válida para la condición de impacto ($y = 0$) es la raíz positiva del discriminante.

\begin{equation}
t = \frac{30 + \sqrt{2860}}{9.8} \approx 8.52 \text{ s}
\end{equation}

Esta magnitud temporal representa el tiempo de vuelo absoluto de la partícula.

\subsubsection*{Conclusión}

El tiempo de vuelo es el mismo tiempo total de movimiento, equivalente a $8.52 \text{ s}$.

\newpage

\subsection*{Enunciado}

Desde lo alto de un edificio de $100 \text{ m}$ de altura, se dispara un proyectil hacia un terreno horizontal con una rapidez de $60 \text{ m/s}$ y $30^\circ$ sobre la horizontal. Determine el desplazamiento del proyectil cuando toca tierra.

\vspace{0.5cm}

\begin{center}
\begin{tikzpicture}
\draw[thick, fill=gray!20] (0,0) rectangle (1.5,4);
\draw[thick] (0,0) -- (8,0) node[right] {Suelo};

\draw[dashed, blue, thick] (1.5,4) parabola bend (3.5, 5.5) (7.5,0);
\filldraw[black] (7.5,0) circle (2pt) node[below] {Punto de impacto};

\filldraw[black] (1.5,4) circle (2pt) node[left=0.2cm] {Lanzamiento};
\draw[->, thick, red] (1.5,4) -- (7.5,0) node[midway, below left] {$\Delta\vec{r}$ (Desplazamiento)};

\end{tikzpicture}
\end{center}

\subsection*{Solución}

\subsubsection*{Conceptos Físicos Involucrados}

El desplazamiento no debe confundirse con la distancia total recorrida a lo largo de la parábola. El desplazamiento es una magnitud estrictamente vectorial que conecta en línea recta la posición espacial inicial con la posición espacial final del objeto, independientemente de la trayectoria curva que haya seguido.

Para definir matemáticamente este vector ($\Delta\vec{r}$), necesitamos cuantificar su componente horizontal ($\Delta x$) y su componente vertical ($\Delta y$).

\subsubsection*{Cálculo de la Componente Vertical}

La coordenada vertical es directa a partir de la geometría del problema. El proyectil parte de la azotea ubicada a cien metros de altura y aterriza en la base plana. Si establecemos nuestro punto de referencia en la azotea, el proyectil desciende cien metros.
\begin{equation}
\Delta y = -100 \text{ m}
\end{equation}

\subsubsection*{Cálculo de la Componente Horizontal}

La componente horizontal, también conocida como el alcance, depende del Movimiento Rectilíneo Uniforme que experimenta el proyectil a lo largo del eje de las abscisas. Se calcula multiplicando la velocidad horizontal constante por el tiempo de vuelo previamente deducido ($t = 8.518 \text{ s}$).

Recuperamos la velocidad horizontal:
\begin{equation}
v_{0x} = 30\sqrt{3} \approx 51.96 \text{ m/s}
\end{equation}

Aplicamos la ecuación fundamental del MRU:
\begin{equation}
\Delta x = v_{0x} \cdot t
\end{equation}

\begin{equation}
\Delta x = 30\sqrt{3} \cdot (8.518)
\end{equation}

\begin{equation}
\Delta x \approx 51.961 \cdot 8.518 \approx 442.6 \text{ m}
\end{equation}

\subsubsection*{Construcción del Vector Desplazamiento y Magnitud}

El vector desplazamiento se expresa sumando ambas componentes direccionales mediante los vectores unitarios $\hat{i}$ y $\hat{j}$:
\begin{equation}
\Delta\vec{r} = 442.6\hat{i} - 100\hat{j} \text{ m}
\end{equation}

Para obtener la magnitud absoluta del desplazamiento (la longitud de la línea recta de color rojo en el diagrama), aplicamos el teorema de Pitágoras sobre las componentes:
\begin{equation}
|\Delta\vec{r}| = \sqrt{(\Delta x)^2 + (\Delta y)^2}
\end{equation}

\begin{equation}
|\Delta\vec{r}| = \sqrt{(442.6)^2 + (-100)^2}
\end{equation}

\begin{equation}
|\Delta\vec{r}| = \sqrt{195894.76 + 10000} = \sqrt{205894.76}
\end{equation}

\begin{equation}
|\Delta\vec{r}| \approx 453.8 \text{ m}
\end{equation}

\subsubsection*{Representación Gráfica del Desplazamiento}

El siguiente modelado en Python proyecta la simulación espacial del evento. Superpone la trayectoria curva real del proyectil con el vector resultante que mide el desplazamiento geométrico neto de la masa.

\begin{lstlisting}[language=Python]
import matplotlib.pyplot as plt
import numpy as np

t_array = np.linspace(0, 8.518, 100)
x_pos = 30 * np.sqrt(3) * t_array
y_pos = 100 + 30 * t_array - 4.9 * t_array**2

plt.figure(figsize=(9, 5))
plt.plot(x_pos, y_pos, color='blue', linewidth=2, label='Trayectoria parabolica')

plt.quiver(0, 100, 442.6, -100, angles='xy', scale_units='xy', scale=1, 
           color='red', width=0.005, label='Vector Desplazamiento')

plt.axhline(0, color='black', linewidth=1)
plt.axvline(0, color='black', linewidth=1)

plt.title('Trayectoria Espacial y Vector Desplazamiento Resultante')
plt.xlabel('Distancia Horizontal x (m)')
plt.ylabel('Altura y (m)')
plt.legend()
plt.grid(True, alpha=0.4)
plt.show()
\end{lstlisting}

\subsubsection*{Conclusión}

El desplazamiento del proyectil es el vector $\Delta\vec{r} = 442.6\hat{i} - 100\hat{j} \text{ m}$, y la magnitud lineal que los separa es de $453.8 \text{ m}$.

\newpage

\subsection*{Enunciado}

Una pelota es disparada con un ángulo de $60^{\circ}$ sobre la horizontal y choca con un árbol, cuando todavía está subiendo con una rapidez de $7\text{ m/s}$ que forma $30^{\circ}$ sobre la horizontal. Considere la magnitud de la aceleración debida a la gravedad como $10\text{ m/s}^{2}$. Determine la distancia horizontal a la que se encuentra el árbol.

\begin{center}
\begin{tikzpicture}[scale=1.2]
\draw[->, thick] (-0.5,0) -- (7,0) node[below] {$x(m)$};
\draw[->, thick] (0,-0.5) -- (0,4) node[left] {$y(m)$};
\draw[fill=black] (0,0) circle (2pt) node[below left] {$0$};

\draw[->, thick, blue] (0,0) -- (1.5, 2.598) node[above left] {$\vec{v}_0$};
\draw (0.8,0) arc (0:60:0.8);
\node at (1.1,0.4) {$60^{\circ}$};

\draw[thick, domain=0:4.24, smooth, variable=\x, green!60!black] plot ({\x}, {1.732*\x - 0.272*\x*\x});

\fill[brown] (4.0,0) rectangle (4.4,2.4);
\fill[green!70!black] (4.2,3.2) circle (1.0);
\fill[green!70!black] (3.6,2.6) circle (0.8);
\fill[green!70!black] (4.8,2.6) circle (0.8);

\draw[->, thick, black] (4.24, 2.45) -- (5.74, 3.31) node[above] {$v = 7\text{ m/s}$};
\draw[dashed] (4.24, 2.45) -- (5.74, 2.45);
\draw (5.24, 2.45) arc (0:30:1.0);
\node at (5.5, 2.65) {$30^{\circ}$};
\end{tikzpicture}
\end{center}

\subsection*{Solución}

\subsubsection*{Análisis conceptual del movimiento de proyectiles}

De acuerdo con los principios de la Física Conceptual, el movimiento en dos dimensiones de un proyectil en caída libre puede entenderse como la superposición de dos movimientos rectilíneos independientes. En el eje horizontal (eje $x$), al no existir fuerzas externas como la resistencia del aire, la aceleración es nula y la componente horizontal de la velocidad se mantiene constante durante toda la trayectoria. En el eje vertical (eje $y$), actúa la fuerza de gravedad de la Tierra, produciendo una aceleración constante dirigida hacia abajo. Esta separación de ejes es fundamental para estructurar la resolución matemática del problema.

\subsubsection*{Cálculo de las componentes de la velocidad en el punto de impacto}

Dado que conocemos la rapidez y el ángulo de la pelota en el instante exacto en que impacta contra el árbol, podemos descomponer este vector velocidad $\vec{v}$ en sus componentes ortogonales utilizando trigonometría básica. La componente adyacente al ángulo corresponde a la velocidad horizontal $v_x$, y la opuesta a la velocidad vertical $v_y$.

\begin{align*}
v_y &= |\vec{v}| \sin(\theta) \\
v_y &= 7 \sin(30^{\circ}) \\
v_y &= 7 (0.5) \\
v_y &= 3.5 \text{ m/s}
\end{align*}

\begin{align*}
v_x &= |\vec{v}| \cos(\theta) \\
v_x &= 7 \cos(30^{\circ}) \\
v_x &= 7 \left(\frac{\sqrt{3}}{2}\right) \\
v_x &\approx 6.06 \text{ m/s}
\end{align*}

\subsubsection*{Determinación de la velocidad inicial de lanzamiento}

Basándonos en el principio de independencia de los movimientos, sabemos que la velocidad horizontal no cambia en el tiempo. Por lo tanto, la componente horizontal de la velocidad inicial $v_{0x}$ es exactamente igual a la componente horizontal en el punto de impacto $v_x$.

\begin{equation*}
v_{0x} = v_x = 6.06 \text{ m/s}
\end{equation*}

Conociendo $v_{0x}$ y el ángulo de lanzamiento original de $60^{\circ}$, podemos encontrar la componente vertical inicial $v_{0y}$ mediante la función tangente, que relaciona el cateto opuesto con el adyacente.

\begin{align*}
\tan(60^{\circ}) &= \frac{v_{0y}}{v_{0x}} \\
v_{0y} &= v_{0x} \tan(60^{\circ}) \\
v_{0y} &= 6.06 \tan(60^{\circ}) \\
v_{0y} &\approx 10.5 \text{ m/s}
\end{align*}

\subsubsection*{Cálculo del tiempo de vuelo hasta el impacto}

Para encontrar la distancia horizontal, primero necesitamos conocer el tiempo total $\Delta t$ que la pelota estuvo en el aire. Utilizamos la ecuación cinemática de la velocidad para el eje vertical, donde la velocidad final es igual a la velocidad inicial más el producto de la aceleración por el tiempo. Al definir hacia arriba como positivo, la aceleración de la gravedad es negativa ($g = -10 \text{ m/s}^2$).

\begin{align*}
v_y &= v_{0y} - g \Delta t \\
3.5 &= 10.5 - 10 \Delta t \\
10 \Delta t &= 10.5 - 3.5 \\
10 \Delta t &= 7 \\
\Delta t &= 0.7 \text{ s}
\end{align*}

\subsubsection*{Cálculo de la distancia horizontal}

Finalmente, dado que el movimiento en el eje $x$ es rectilíneo uniforme (velocidad constante), el desplazamiento horizontal $x$ es simplemente el producto de la velocidad horizontal por el tiempo de vuelo.

\begin{align*}
x &= v_{0x} \Delta t \\
x &= (6.06)(0.7) \\
x &\approx 4.24 \text{ m}
\end{align*}

\subsubsection*{Representación computacional del modelo cinemático}

Para consolidar la comprensión pedagógica de la independencia de los ejes, el siguiente script en Python genera una simulación visual de la trayectoria y grafica los vectores de velocidad en el lanzamiento y en el impacto.

\begin{lstlisting}[language=Python]
import matplotlib.pyplot as plt
import numpy as np

v0x = 6.06
v0y = 10.5
g = 10.0
t_vuelo = 0.7

t = np.linspace(0, t_vuelo, 100)
x = v0x * t
y = v0y * t - 0.5 * g * t**2

plt.figure(figsize=(8, 5))
plt.plot(x, y, label='Trayectoria parabolica', color='green')

plt.quiver(0, 0, v0x, v0y, angles='xy', scale_units='xy', scale=1, color='blue', label='V. inicial')

vy_impacto = v0y - g * t_vuelo
plt.quiver(x[-1], y[-1], v0x, vy_impacto, angles='xy', scale_units='xy', scale=1, color='red', label='V. impacto')

plt.title('Cinematica del proyectil y vectores de velocidad')
plt.xlabel('Distancia x (m)')
plt.ylabel('Altura y (m)')
plt.axhline(0, color='black', linewidth=0.8)
plt.legend()
plt.grid(True, linestyle='--')
plt.show()
\end{lstlisting}

\subsubsection*{Conclusión}

La distancia horizontal a la que se encuentra el árbol es de $4.24\text{ m}$.

\newpage

\subsection*{Enunciado}

Una pelota es disparada con un ángulo de $60^{\circ}$ sobre la horizontal y choca con un árbol, cuando todavía está subiendo con una rapidez de $7\text{ m/s}$ que forma $30^{\circ}$ sobre la horizontal. Considere la magnitud de la aceleración debida a la gravedad como $10\text{ m/s}^{2}$. Determine la altura a la que se estrella la pelota.

\begin{center}
\begin{tikzpicture}[scale=1.2]
\draw[->, thick] (-0.5,0) -- (7,0) node[below] {$x(m)$};
\draw[->, thick] (0,-0.5) -- (0,4) node[left] {$y(m)$};
\draw[fill=black] (0,0) circle (2pt) node[below left] {$0$};

\draw[->, thick, blue] (0,0) -- (1.5, 2.598) node[above left] {$\vec{v}_0$};
\draw (0.8,0) arc (0:60:0.8);
\node at (1.1,0.4) {$60^{\circ}$};

\draw[thick, domain=0:4.24, smooth, variable=\x, green!60!black] plot ({\x}, {1.732*\x - 0.272*\x*\x});

\fill[brown] (4.0,0) rectangle (4.4,2.4);
\fill[green!70!black] (4.2,3.2) circle (1.0);
\fill[green!70!black] (3.6,2.6) circle (0.8);
\fill[green!70!black] (4.8,2.6) circle (0.8);

\draw[->, thick, black] (4.24, 2.45) -- (5.74, 3.31) node[above] {$v = 7\text{ m/s}$};
\draw[dashed] (4.24, 2.45) -- (5.74, 2.45);
\draw (5.24, 2.45) arc (0:30:1.0);
\node at (5.5, 2.65) {$30^{\circ}$};

\draw[<->, thick, red] (4.6, 0) -- (4.6, 2.45) node[midway, right] {$y = ?$};
\draw[dashed, red] (4.24, 2.45) -- (4.6, 2.45);
\end{tikzpicture}
\end{center}

\subsection*{Solución}

\subsubsection*{Análisis conceptual del desplazamiento vertical}

Para determinar la altura de impacto, debemos centrarnos exclusivamente en el eje vertical. De acuerdo con el marco teórico de la Física Conceptual, el movimiento vertical de un proyectil es idéntico al de un objeto lanzado verticalmente hacia arriba. La altura alcanzada depende de la velocidad inicial vertical y de la acción desaceleradora de la gravedad durante el tiempo de vuelo. Para calcular esta altura, necesitamos primero deducir la velocidad vertical inicial y el tiempo exacto que transcurre hasta el choque.

\subsubsection*{Deducción de parámetros cinemáticos preliminares}

El problema establece un vínculo entre el eje horizontal y vertical a través del tiempo. Primero, calculamos la velocidad en el instante del impacto mediante descomposición vectorial.

\begin{align*}
v_x &= 7 \cos(30^{\circ}) \approx 6.06 \text{ m/s} \\
v_y &= 7 \sin(30^{\circ}) = 3.5 \text{ m/s}
\end{align*}

Como la velocidad horizontal es constante en ausencia de resistencia del aire, la velocidad inicial en el eje $x$ es $v_{0x} = 6.06 \text{ m/s}$. Usando el ángulo de lanzamiento inicial de $60^{\circ}$, encontramos la velocidad vertical inicial.

\begin{align*}
v_{0y} &= v_{0x} \tan(60^{\circ}) \\
v_{0y} &= 6.06 \tan(60^{\circ}) \approx 10.5 \text{ m/s}
\end{align*}

Con la velocidad vertical inicial y final conocidas, empleamos la ecuación de cambio de velocidad para despejar el tiempo de vuelo $\Delta t$, asumiendo una aceleración gravitacional de $g = -10 \text{ m/s}^2$.

\begin{align*}
\Delta t &= \frac{v_{0y} - v_y}{g} \\
\Delta t &= \frac{10.5 - 3.5}{10} \\
\Delta t &= 0.7 \text{ s}
\end{align*}

\subsubsection*{Cálculo de la altura del impacto}

Una vez establecido el tiempo de vuelo y la velocidad vertical inicial, aplicamos la ecuación cinemática de posición para objetos con aceleración constante. La altura $y$ es el resultado del desplazamiento provocado por la velocidad inicial menos la pérdida de altura debida a la gravedad en ese intervalo de tiempo.

\begin{align*}
y &= v_{0y} \Delta t - \frac{1}{2} g \Delta t^2 \\
y &= (10.5)(0.7) - \frac{1}{2}(10)(0.7)^2 \\
y &= 7.35 - 5(0.49) \\
y &= 7.35 - 2.45 \\
y &= 4.9 \text{ m}
\end{align*}

\subsubsection*{Conclusión}

La altura a la que se estrella la pelota contra el árbol es de $4.9\text{ m}$.

\newpage

\subsection*{Enunciado}

Se lanza un objeto de tal manera que pasa justamente sobre dos obstáculos, cada uno de $11\text{ m}$ de altura y que están separados $52\text{ m}$. Sabiendo que el tiempo empleado en recorrer la distancia que hay entre los dos obstáculos es de $2.6\text{ s}$ y considerando la magnitud de la aceleración de la gravedad como $10\text{ m/s}^2$, calcule la velocidad inicial del objeto.

\begin{center}
\begin{tikzpicture}[scale=0.15]
\draw[->, thick] (-5,0) -- (85,0) node[below] {x};
\draw[->, thick] (0,-5) -- (0,30) node[left] {y};
\fill[black] (0,0) circle (10pt) node[below left] {0};

\draw[dashed, thick, orange] (0,0) parabola bend (39.4, 19.4) (78.9, 0);

\draw[->, thick, green!60!black] (0,0) -- (15, 14.8) node[above left, black] {$\vec{v}_0$};
\draw (8,0) arc (0:44.6:8);
\node at (11, 4) {$\theta$};

\fill[gray!40] (13.4,0) rectangle (16.4,11);
\draw[thick] (13.4,0) rectangle (16.4,11);
\fill[gray!40] (62.4,0) rectangle (65.4,11);
\draw[thick] (62.4,0) rectangle (65.4,11);

\node[above] at (14.9, 11) {$11\text{ m}$};
\node[above] at (63.9, 11) {$11\text{ m}$};

\draw[<->] (17, 12) -- (61.4, 12) node[midway, fill=white] {$52\text{ m}$};
\node at (39.4, 5) {$t = 2.6\text{ s}$};
\fill[orange] (5,0) circle (10pt);
\end{tikzpicture}
\end{center}

\subsection*{Solución}

\subsubsection*{Resultado Principal}
La velocidad inicial del lanzamiento posee una magnitud de $28.09\text{ m/s}$ y un ángulo de elevación de $44.6^{\circ}$ respecto a la horizontal.

\subsubsection*{Fundamentación del Vuelo Simétrico}
Bajo los principios de la Física Conceptual de Hewitt, el movimiento parabólico se compone de ejes independientes. Los parámetros iniciales entre las cotas elevadas quedan establecidos como contexto base: el desplazamiento uniforme entre los obstáculos separados por $52\text{ m}$ en $2.6\text{ s}$ define irrevocablemente una velocidad horizontal constante $v_x = 20\text{ m/s}$. Por la simetría de la parábola respecto a su vértice, la componente vertical de la velocidad exactamente a los $11\text{ m}$ de altura se fija de manera estándar en $v_{y(11)} = 13\text{ m/s}$.

\subsubsection*{Determinación Directa del Vector Inicial}
Aplicando una transformación directa hacia el origen topológico mediante la cinemática vertical, la velocidad inicial en el eje $y$ se obtiene al evaluar el nivel base:
\begin{equation*}
v_{0y} = \sqrt{13^2 + 2(10)(11)} = 19.72\text{ m/s}
\end{equation*}
Por composición vectorial ortogonal, la rapidez inicial total resulta en:
\begin{equation*}
v_0 = \sqrt{20^2 + 19.72^2} = 28.09\text{ m/s}
\end{equation*}

\subsubsection*{Conclusión}
La magnitud de la velocidad inicial es de $28.09\text{ m/s}$.

\newpage

\subsection*{Enunciado}

Se lanza un objeto de tal manera que pasa justamente sobre dos obstáculos, cada uno de $11\text{ m}$ de altura y que están separados $52\text{ m}$. Sabiendo que el tiempo empleado en recorrer la distancia que hay entre los dos obstáculos es de $2.6\text{ s}$ y considerando la magnitud de la aceleración de la gravedad como $10\text{ m/s}^2$, calcule el alcance máximo del movimiento.

\begin{center}
\begin{tikzpicture}[scale=0.15]
\draw[->, thick] (-5,0) -- (85,0) node[below] {x};
\draw[->, thick] (0,-5) -- (0,30) node[left] {y};
\fill[black] (0,0) circle (10pt) node[below left] {0};

\draw[dashed, thick, orange] (0,0) parabola bend (39.4, 19.4) (78.9, 0);

\draw[->, thick, green!60!black] (0,0) -- (15, 14.8) node[above left, black] {$\vec{v}_0$};
\draw (8,0) arc (0:44.6:8);
\node at (11, 4) {$\theta$};

\fill[gray!40] (13.4,0) rectangle (16.4,11);
\draw[thick] (13.4,0) rectangle (16.4,11);
\fill[gray!40] (62.4,0) rectangle (65.4,11);
\draw[thick] (62.4,0) rectangle (65.4,11);

\node[above] at (14.9, 11) {$11\text{ m}$};
\node[above] at (63.9, 11) {$11\text{ m}$};

\draw[<->] (17, 12) -- (61.4, 12) node[midway, fill=white] {$52\text{ m}$};
\node at (39.4, 5) {$t = 2.6\text{ s}$};
\end{tikzpicture}
\end{center}

\subsection*{Solución}

\subsubsection*{Resultado Principal}
El alcance máximo del proyectil sobre el plano horizontal es de $78.9\text{ m}$.

\subsubsection*{Consolidación del Tiempo Global}
A partir de la componente vertical validada previamente en $19.72\text{ m/s}$, el procesamiento temporal estándar para el vuelo completo sobre el horizonte determina directamente que el objeto permanecerá suspendido durante $3.94\text{ s}$.

\subsubsection*{Cálculo Directo de Alcance}
Dado el principio de inercia horizontal postulado por Hewitt, el alcance $x$ es una progresión lineal inalterada. Extrapolando a la duración total del vuelo:
\begin{equation*}
x = 20 \times 3.94 = 78.9\text{ m}
\end{equation*}

\subsubsection*{Representación Computacional de la Trayectoria}
Para ilustrar la morfología parabólica en relación con los obstáculos, el siguiente modelo en Python dibuja la trayectoria exacta y verifica geométricamente el paso preciso sobre las cotas restrictivas de $11\text{ m}$.

\begin{lstlisting}[language=Python]
import matplotlib.pyplot as plt
import numpy as np

v0x = 20.0
v0y = 19.72
g = 10.0
t_vuelo = 3.94

t = np.linspace(0, t_vuelo, 200)
x = v0x * t
y = v0y * t - 0.5 * g * t**2

plt.figure(figsize=(9, 4))
plt.plot(x, y, label='Proyeccion parabolica', color='orange', linestyle='--')

plt.bar([13.4, 62.4], [11, 11], width=3, color='gray', align='edge', label='Obstaculos (11m)')

plt.title('Modelo cinematico del alcance')
plt.xlabel('Distancia x (m)')
plt.ylabel('Altura y (m)')
plt.axhline(0, color='black')
plt.legend()
plt.grid(True, alpha=0.3)
plt.show()
\end{lstlisting}

\subsubsection*{Conclusión}
El alcance máximo es de $78.9\text{ m}$.

\newpage

\subsection*{Enunciado}

Desde lo alto de un edificio de $100\text{ m}$ de altura, se dispara un proyectil hacia un terreno horizontal con una rapidez de $60\text{ m/s}$ y un ángulo de $30^{\circ}$ sobre la horizontal. Determine el vector desplazamiento del proyectil desde el lanzamiento hasta que impacta contra el suelo.

\begin{center}
\begin{tikzpicture}[scale=0.08]
\draw[->, thick] (0,0) -- (100,0) node[below] {$x$};
\draw[->, thick] (0,0) -- (0,35) node[left] {$y$};
\draw[thick, fill=gray!30] (0,0) rectangle (10,25);
\node at (5,12.5) {$100\text{m}$};

\draw[dashed, thick, blue] (10,25) parabola bend (35, 36) (75,0);
\draw[->, thick, red] (10,25) -- (20, 31) node[right] {$\vec{v}_0$};
\draw (20,25) arc (0:30:10);
\node at (24,27) {$30^{\circ}$};

\draw[->, thick, black] (10,25) -- (75,0) node[midway, above, sloped] {$\Delta\vec{r}$};
\end{tikzpicture}
\end{center}

\subsection*{Solución}

\subsubsection*{Análisis del sistema de referencia}
Para resolver el desplazamiento de un proyectil lanzado desde una altura, definimos el punto de lanzamiento como el origen en el eje $x$ ($x_0 = 0$) y en una altura inicial $y_0 = 100\text{ m}$. El suelo, por tanto, se encuentra en $y = 0\text{ m}$. El desplazamiento vertical $\Delta y$ será entonces $y_f - y_0 = 0 - 100 = -100\text{ m}$.

\subsubsection*{Descomposición de la velocidad inicial}
Siguiendo los principios de la Física Conceptual, descomponemos el vector velocidad inicial $\vec{v}_0$ en sus componentes rectangulares independientes:
\begin{align*}
v_{0x} &= v_0 \cos(30^{\circ}) = 60 \cos(30^{\circ}) \approx 51.96\text{ m/s} \\
v_{0y} &= v_0 \sin(30^{\circ}) = 60 \sin(30^{\circ}) = 30\text{ m/s}
\end{align*}

\subsubsection*{Cálculo del tiempo de vuelo}
Utilizamos la ecuación cinemática para el eje vertical, donde el proyectil cae bajo la influencia de la gravedad ($g = -10\text{ m/s}^2$):
\begin{equation*}
\Delta y = v_{0y} \Delta t + \frac{1}{2} g \Delta t^2
\end{equation*}
Sustituyendo los valores:
\begin{equation*}
-100 = 30 \Delta t - 5 \Delta t^2 \implies 5 \Delta t^2 - 30 \Delta t - 100 = 0
\end{equation*}
Resolviendo la ecuación cuadrática $\Delta t^2 - 6 \Delta t - 20 = 0$, obtenemos un tiempo de vuelo $\Delta t \approx 8.38\text{ s}$.

\subsubsection*{Determinación del vector desplazamiento}
El desplazamiento horizontal se obtiene mediante movimiento rectilíneo uniforme:
\begin{equation*}
\Delta x = v_{0x} \Delta t = 51.96 \times 8.38 \approx 435.42\text{ m}
\end{equation*}
El vector desplazamiento total $\Delta\vec{r}$ se expresa como:
\begin{equation*}
\Delta\vec{r} = \Delta x \hat{i} + \Delta y \hat{j} = 435.42 \hat{i} - 100 \hat{j} \text{ (m)}
\end{equation*}

\subsubsection*{Representación mediante modelado computacional}
El siguiente script permite visualizar la trayectoria parabólica desde la altura del edificio:
\begin{lstlisting}[language=Python]
import matplotlib.pyplot as plt
import numpy as np

v0x, v0y, g, h0 = 51.96, 30.0, 10.0, 100.0
t = np.linspace(0, 8.38, 100)
x = v0x * t
y = h0 + v0y * t - 0.5 * g * t**2

plt.plot(x, y)
plt.axhline(0, color='black')
plt.title('Trayectoria desde edificio')
plt.grid(True)
plt.show()
\end{lstlisting}

\subsubsection*{Conclusión}
El desplazamiento total es de $\Delta\vec{r} = 435.42\hat{i} - 100\hat{j} \text{ m}$.

\end{document}